% Electrostatics and Electric Fields — Physics Upper Sixth — Cours signé OnBuch+
% Official MINESEC syllabus. Compiler avec : tectonic PHY03-electrostatics-and-electric-fields.tex
\newcommand{\DOCMATIERE}{Physics}
\newcommand{\DOCNIVEAU}{Upper Sixth}
\newcommand{\DOCMODULE}{Upper Sixth — Physics}
\newcommand{\DOCLECON}{3}
\newcommand{\DOCTITRE}{Electrostatics and Electric Fields}
% OnBuch+ — préambule « cours signés » ENRICHI (compilé avec tectonic / XeLaTeX)
% Système visuel « Pop » d'OnBuch+ : fond crème (jamais blanc), bordures encre
% épaisses, ombres « dures » décalées, titres Archivo Black en capitales.
% Version MATHÉMATIQUES : graphes (pgfplots/tikz), boîtes pédagogiques
% (définition, propriété, méthode, attention, le savais-tu, exercice résolu,
% exercice, TP, bilan), page de garde et signature OnBuch+ ancrée en bas.
%
% Macros attendues dans main.tex AVANT \input{preamble} :
%   \DOCMATIERE  \DOCNIVEAU  \DOCMODULE  \DOCLECON (numéro)  \DOCTITRE
\documentclass[11pt]{article}

\usepackage{fontspec}
\usepackage[english]{babel}
\usepackage[a4paper,margin=2cm,top=3.4cm,bottom=2.8cm,headheight=1.4cm,footskip=1.3cm]{geometry}
\usepackage{amsmath,amssymb,amsfonts}
\usepackage{mathtools} % \xmapsto, \coloneqq... souvent utilisés par les modèles
\usepackage{cancel} % simplifications barrées (bilans de forces)
% \abs{z} (module, valeur absolue) et \norm{u} (norme), utilisés par les modèles
\providecommand{\abs}{}\let\abs\relax\DeclarePairedDelimiter{\abs}{\lvert}{\rvert}
\providecommand{\norm}{}\let\norm\relax\DeclarePairedDelimiter{\norm}{\lVert}{\rVert}
\usepackage[table]{xcolor} % [table] : fournit aussi \rowcolors (alternance de lignes)
\usepackage{pagecolor}
\usepackage{tikz}
\usetikzlibrary{arrows.meta,positioning,calc,shapes.geometric,shapes.misc,shapes.arrows,decorations.pathmorphing,decorations.pathreplacing,decorations.markings,patterns,shadows,mindmap,trees,fit,backgrounds,matrix,intersections,angles,quotes}
\usepackage{pgfplots}
\usepackage[version=4]{mhchem} % équations nucléaires \ce{^{235}_{92}U}
\usepackage[european,siunitx]{circuitikz} % schémas électriques
\pgfplotsset{compat=1.18}
\usepgfplotslibrary{fillbetween,groupplots} % aires entre courbes (calcul intégral)
\usepackage{enumitem}
\usepackage{fancyhdr}
% [most] charge toutes les bibliothèques tcolorbox (listings, minted, raster,
% documentation...) et gonfle inutilement la mémoire TeX sur un document long
% (~300 boîtes) : on ne charge que ce qui est réellement utilisé.
\usepackage[skins,breakable]{tcolorbox}
\usepackage{graphicx}
\usepackage{colortbl}
\usepackage{booktabs}
\usepackage{tabularx}
\usepackage{diagbox} % case d'en-tête diagonale des tableaux à double entrée
% Colonnes extensibles centrée / gauche / droite (C, L, R), souvent utilisées par les modèles
\newcolumntype{C}{>{\centering\arraybackslash}X}
\newcolumntype{L}{>{\raggedright\arraybackslash}X}
\newcolumntype{R}{>{\raggedleft\arraybackslash}X}
\usepackage{array}
\usepackage{multicol}
\usepackage{siunitx}
% parse-numbers=false : accepte aussi \SI{2,5 \times 10^{-3}}{...}
\sisetup{locale=UK,output-decimal-marker={.},group-minimum-digits=4,parse-numbers=false}
\DeclareSIUnit\ohm{\ensuremath{\symup{\Omega}}} % U+2126 absent de la police
\DeclareSIUnit\division{div} % réglages d'oscilloscope (ms/div, V/div)
\DeclareSIUnit\tr{tr} % tours (vitesse de rotation, ex. \SI{1500}{\tr\per\minute})
\usepackage{qrcode}
% \degree : commande fréquemment utilisée par les modèles pour un angle ou une
% température en dehors de \SI{}{} (non fournie par les paquets chargés).
\providecommand{\degree}{^{\circ}}
% \qed (fin de démonstration), utilisé par les modèles sans amsthm
\providecommand{\qed}{\hfill$\square$}
% \barre : double barre des valeurs interdites dans un tableau de variations (tkz-tab)
\providecommand{\barre}{\ensuremath{\|}}
% environnement proof (amsthm non chargé) utilisé par les modèles
\ifdefined\proof\else\newenvironment{proof}[1][Proof]{\par\noindent\textit{#1.}\ }{\qed\par}\fi
\usepackage{lastpage}
\usepackage{hyperref}
\hypersetup{hidelinks}

% ============================================================
% Palette « Pop » OnBuch+ — mêmes hex que la classe P (plus_ui.dart)
% ============================================================
\definecolor{poporange}{HTML}{F59321}
\definecolor{popdark}{HTML}{E07A0C}
\definecolor{popcream}{HTML}{FFF6E8}
\definecolor{popink}{HTML}{1C1714}
\definecolor{poppurple}{HTML}{7A5AE0}
\definecolor{popgreen}{HTML}{2E9E6A}
\definecolor{poppink}{HTML}{E0457B}
\definecolor{popblue}{HTML}{2F7DE1}
\definecolor{popgold}{HTML}{C98A12}
\definecolor{popmuted}{HTML}{8A7F76}
\definecolor{popline}{HTML}{EADFD2}
\definecolor{popgray}{HTML}{8A7F76}
\colorlet{popgrey}{popgray}
% Alias tolérants (le modèle nomme parfois les couleurs différemment)
\colorlet{popwhite}{white}
\colorlet{popblack}{popink}
\colorlet{popred}{poppink}
\colorlet{popyellow}{popgold}
% Teintes claires pour fonds de tableaux / zones
\colorlet{popblueL}{popblue!10!white}
\colorlet{poporangeL}{poporange!14!white}
\colorlet{popgreenL}{popgreen!12!white}
\colorlet{poppurpleL}{poppurple!12!white}
\colorlet{poppinkL}{poppink!12!white}
\colorlet{popgoldL}{popgold!14!white}

\pagecolor{popcream}

% ============================================================
% Typographie
% ============================================================
\defaultfontfeatures{Ligatures=TeX}
\setmainfont{PlusJakartaSans-Medium.ttf}[Path=../../../fonts/,BoldFont=PlusJakartaSans-Bold.ttf]
% Maths en sans-serif (Fira Math) avec chiffres et lettres droites de Plus
% Jakarta, pour que les formules se fondent dans le texte.
\usepackage[math-style=ISO,bold-style=ISO]{unicode-math}
\setmathfont{FiraMath-Regular.otf}[Path=../../../fonts/]
\setmathfont{PlusJakartaSans-Medium.ttf}[Path=../../../fonts/,range={up/num,up/{Latin,latin},\mathup}]
% (icomma retiré : décimales avec point en anglais)
% Caractères mathématiques Unicode absents des polices de texte (Plus Jakarta,
% Archivo Black) : rendus par leur équivalent mathématique, en texte comme en
% formule — sinon ils s'affichent en carré vide (ex. « Arithmétique dans ℤ »).
\usepackage{newunicodechar}
\newunicodechar{ℝ}{\ensuremath{\mathbb{R}}}
\newunicodechar{ℤ}{\ensuremath{\mathbb{Z}}}
\newunicodechar{ℂ}{\ensuremath{\mathbb{C}}}
\newunicodechar{ℕ}{\ensuremath{\mathbb{N}}}
\newunicodechar{ℚ}{\ensuremath{\mathbb{Q}}}
\newunicodechar{𝔻}{\ensuremath{\mathbb{D}}}
\newunicodechar{α}{\ensuremath{\alpha}}
\newunicodechar{β}{\ensuremath{\beta}}
\newunicodechar{γ}{\ensuremath{\gamma}}
\newunicodechar{δ}{\ensuremath{\delta}}
\newunicodechar{ε}{\ensuremath{\varepsilon}}
\newunicodechar{θ}{\ensuremath{\theta}}
\newunicodechar{λ}{\ensuremath{\lambda}}
\newunicodechar{μ}{\ensuremath{\mu}}
\newunicodechar{σ}{\ensuremath{\sigma}}
\newunicodechar{τ}{\ensuremath{\tau}}
\newunicodechar{φ}{\ensuremath{\varphi}}
\newunicodechar{ψ}{\ensuremath{\psi}}
\newunicodechar{ω}{\ensuremath{\omega}}
\newunicodechar{Σ}{\ensuremath{\Sigma}}
% majuscules produites par \MakeUppercase dans les titres (α-aminés -> Α)
\newunicodechar{Α}{\ensuremath{\alpha}}
\newunicodechar{Β}{\ensuremath{\beta}}
\newunicodechar{Γ}{\ensuremath{\Gamma}}
\newunicodechar{Δ}{\ensuremath{\Delta}}
\newunicodechar{Ε}{\ensuremath{\varepsilon}}
\newunicodechar{Θ}{\ensuremath{\Theta}}
\newunicodechar{Λ}{\ensuremath{\Lambda}}
\newunicodechar{Μ}{\ensuremath{\mu}}
\newunicodechar{Τ}{\ensuremath{\tau}}
\newunicodechar{Φ}{\ensuremath{\Phi}}
\newunicodechar{Ψ}{\ensuremath{\Psi}}
\newunicodechar{Ω}{\ensuremath{\Omega}}
\newunicodechar{↦}{\ensuremath{\mapsto}}
\newunicodechar{∈}{\ensuremath{\in}}
\newunicodechar{∉}{\ensuremath{\notin}}
\newunicodechar{∧}{\ensuremath{\wedge}}
\newunicodechar{⇔}{\ensuremath{\Leftrightarrow}}
\newunicodechar{⟺}{\ensuremath{\Longleftrightarrow}}
\newunicodechar{⇒}{\ensuremath{\Rightarrow}}
\newunicodechar{⟹}{\ensuremath{\Longrightarrow}}
\newunicodechar{∘}{\ensuremath{\circ}}
\newunicodechar{∪}{\ensuremath{\cup}}
\newunicodechar{∩}{\ensuremath{\cap}}
\newunicodechar{≡}{\ensuremath{\equiv}}
\newunicodechar{⊂}{\ensuremath{\subset}}
\newunicodechar{⊥}{\ensuremath{\perp}}
\newunicodechar{∃}{\ensuremath{\exists}}
\newunicodechar{∀}{\ensuremath{\forall}}
\newunicodechar{∛}{\ensuremath{\sqrt[3]{\,}}}
\newunicodechar{∜}{\ensuremath{\sqrt[4]{\,}}}
\newunicodechar{⃗}{\ensuremath{{}^{\to}}}
\newunicodechar{ⁿ}{\ifmmode^{n}\else\textsuperscript{\ensuremath{n}}\fi}
\newunicodechar{ˣ}{\ifmmode^{x}\else\textsuperscript{\ensuremath{x}}\fi}
\newunicodechar{ᵇ}{\ifmmode^{b}\else\textsuperscript{\ensuremath{b}}\fi}
\newunicodechar{ʳ}{\ifmmode^{r}\else\textsuperscript{\ensuremath{r}}\fi}
\newunicodechar{ᵘ}{\ifmmode^{u}\else\textsuperscript{\ensuremath{u}}\fi}
\newunicodechar{ʸ}{\ifmmode^{y}\else\textsuperscript{\ensuremath{y}}\fi}
\newunicodechar{ᵃ}{\ifmmode^{a}\else\textsuperscript{\ensuremath{a}}\fi}
\newunicodechar{ᵉ}{\ifmmode^{e}\else\textsuperscript{\ensuremath{e}}\fi}
\newunicodechar{ᵏ}{\ifmmode^{k}\else\textsuperscript{\ensuremath{k}}\fi}
\newunicodechar{ᵖ}{\ifmmode^{p}\else\textsuperscript{\ensuremath{p}}\fi}
\newunicodechar{ᴺ}{\ifmmode^{N}\else\textsuperscript{\ensuremath{N}}\fi}
\newunicodechar{ᶜ}{\ifmmode^{c}\else\textsuperscript{\ensuremath{c}}\fi}
\newunicodechar{ʰ}{\ifmmode^{h}\else\textsuperscript{\ensuremath{h}}\fi}
\newunicodechar{ᵠ}{\ifmmode^{q}\else\textsuperscript{\ensuremath{q}}\fi}
\newunicodechar{ₙ}{\ifmmode_{n}\else\textsubscript{\ensuremath{n}}\fi}
\newunicodechar{ₐ}{\ifmmode_{a}\else\textsubscript{\ensuremath{a}}\fi}
\newunicodechar{ᵢ}{\ifmmode_{i}\else\textsubscript{\ensuremath{i}}\fi}
\newunicodechar{ᵣ}{\ifmmode_{r}\else\textsubscript{\ensuremath{r}}\fi}
\newunicodechar{ₖ}{\ifmmode_{k}\else\textsubscript{\ensuremath{k}}\fi}
\newunicodechar{ᵦ}{\ifmmode_{\beta}\else\textsubscript{\ensuremath{\beta}}\fi}
\newunicodechar{ₓ}{\ifmmode_{x}\else\textsubscript{\ensuremath{x}}\fi}
\newfontfamily\popdisplay{ArchivoBlack-Regular.ttf}[Path=../../../fonts/]
\newfontfamily\poplogo{SpaceGrotesk-Bold.ttf}[Path=../../../fonts/]
\newfontfamily\popsemi{PlusJakartaSans-SemiBold.ttf}[Path=../../../fonts/]
\newfontfamily\popxbold{PlusJakartaSans-ExtraBold.ttf}[Path=../../../fonts/]
\newfontfamily\popxboldls{PlusJakartaSans-ExtraBold.ttf}[Path=../../../fonts/,LetterSpace=3.0]

\setlist[enumerate]{leftmargin=1.7em,itemsep=5pt,topsep=4pt}
\setlist[itemize]{leftmargin=1.7em,itemsep=4pt,topsep=4pt,label={\color{poporange}\textbullet}}
\setlength{\parindent}{0pt}
\setlength{\parskip}{5pt}
\renewcommand{\arraystretch}{1.25}

% pgfplots : style Pop par défaut
\pgfplotsset{
  every axis/.append style={axis line style={popink,line width=1pt},tick style={popink},
    grid style={popline,line width=0.5pt},label style={font=\small},tick label style={font=\footnotesize}},
  popcurve/.style={poporange,line width=1.8pt,no markers,smooth},
}
% Même style disponible dans un \draw TikZ simple (hors pgfplots)
\tikzset{popcurve/.style={poporange,line width=1.8pt}}
\tikzset{popgrid/.style={popline,very thin}}
\colorlet{popcurve}{poporange} % le nom du style est parfois utilisé comme couleur

% ============================================================
% Pill / squiggle
% ============================================================
\newcommand{\poppill}[3]{%
  \tikz[baseline=-0.55ex]\node[draw=popink,line width=1.2pt,rounded corners=2.6pt,%
    fill=#2,inner xsep=6.5pt,inner ysep=3pt]{%
    \popxboldls\fontsize{7}{7}\selectfont\color{#3}\MakeUppercase{#1}};%
}
\newcommand{\popsquiggle}[1]{%
  \tikz[baseline=-0.4ex]{
    \draw[#1,line width=2.6pt,line cap=round]
      plot [smooth] coordinates {(0,0.09) (0.34,0.01) (0.68,0.21) (1.02,0.01) (1.36,0.10)};
  }%
}
% Mot-clé surligné (marqueur orange sous le texte)
\newcommand{\cle}[1]{{\popxbold\color{popdark}#1}}

% ============================================================
% En-tête / pied
% ============================================================
\pagestyle{fancy}
\fancyhf{}
\renewcommand{\headrulewidth}{0pt}
\renewcommand{\footrulewidth}{0pt}
\lhead{\raisebox{-0.15ex}{\poplogo\fontsize{13}{13}\selectfont\color{popink}OnBuch\color{poporange}+}}
\rhead{\poppill{\DOCMATIERE\ \textbullet\ \DOCNIVEAU\ \textbullet\ Chapter \DOCLECON}{white}{popink}}
\lfoot{\popxbold\fontsize{7.6}{7.6}\selectfont\color{popmuted}\MakeUppercase{Signed course OnBuch+}\ \textbullet\ {\popsemi\DOCTITRE}}
\rfoot{\tikz[baseline=-0.6ex]\node[rounded corners=8.5pt,fill=popink,minimum height=17pt,minimum width=17pt,inner xsep=5pt,inner sep=0]{\popxbold\fontsize{8}{8}\selectfont\color{popcream}\thepage\,/\,\pageref*{LastPage}};}

% ============================================================
% Titres
% ============================================================
\newcommand{\chaptitre}[2]{%
  \begin{tcolorbox}[enhanced,colback=poporange,colframe=popink,boxrule=2.6pt,arc=11pt,
      drop shadow=popink,left=15pt,right=15pt,top=12pt,bottom=11pt,before skip=3pt,after skip=18pt]
    \poppill{#2}{popink}{popcream}\par\vspace{9pt}
    {\raggedright\hyphenpenalty=10000\exhyphenpenalty=10000\popdisplay\fontsize{22}{24}\selectfont\color{popcream}\MakeUppercase{#1}\par}
  \end{tcolorbox}
}
% Section numérotée : pastille orange avec le numéro + titre Archivo
\newcounter{popsec}
\newcommand{\coursec}[1]{%
  \stepcounter{popsec}\setcounter{popsub}{0}%
  \par\vspace{16pt}\noindent
  \begin{minipage}{\linewidth}
  \tikz[baseline=-0.7ex]\node[draw=popink,line width=1.6pt,fill=poporange,rounded corners=4pt,minimum size=22pt,inner sep=0]{\popdisplay\fontsize{12}{12}\selectfont\color{popcream}\thepopsec};%
  \hspace{8pt}{\popdisplay\fontsize{15}{17}\selectfont\color{popink}\MakeUppercase{#1}}\par
  \vspace{4pt}\noindent{\color{poporange}\rule{36pt}{3.6pt}}
  \end{minipage}\par\nopagebreak\vspace{8pt}
}
\newcounter{popsub}
\newcommand{\courssub}[1]{%
  \stepcounter{popsub}%
  \par\vspace{9pt}\noindent{\popxbold\fontsize{11.5}{13}\selectfont\color{popink}{\color{poporange}\thepopsec.\thepopsub}\hspace{6pt}#1}\par\nopagebreak\vspace{4pt}
}

% ============================================================
% Boîtes pédagogiques — même recette : fond blanc, bordure épaisse colorée,
% ombre dure encre, badge plein. Titre optionnel [#1].
% ============================================================
\tcbset{popbase/.style={breakable,enhanced,colback=white,boxrule=2.3pt,arc=9pt,
  shadow={3pt}{-3pt}{0pt}{popink},left=14pt,right=14pt,top=11pt,bottom=11pt,before skip=11pt,after skip=15pt,
  fontupper=\popsemi\fontsize{10.6}{14.5}\selectfont\color{popink},
  fontlower=\popsemi\fontsize{10.6}{14.5}\selectfont\color{popink}}}
% Coche dessinée (le glyphe ✓ n'existe pas dans les polices du document)
\newcommand{\popcheck}[1]{\tikz[baseline=-0.5ex]\draw[#1,line width=1.6pt,line cap=round,line join=round](-0.13,0.0)--(-0.04,-0.09)--(0.13,0.1);}
\providecommand{\coche}{\popcheck{popgreen}}
\providecommand{\resultat}[1]{\par\smallskip\noindent\fcolorbox{popgreen}{popgreenL}{\parbox{\dimexpr\linewidth-2\fboxsep-2\fboxrule\relax}{\textbf{Result.} #1}}\par\smallskip}
\newcommand{\popboxhead}[3]{\poppill{#1}{#2}{white}\ifx\relax#3\relax\else\hspace{6pt}{\popxbold\fontsize{10.6}{12}\selectfont\color{popink}#3}\fi\par\vspace{7pt}}

\newtcolorbox{aretenir}[1][]{popbase,colframe=popgreen,before upper={\popboxhead{Key points}{popgreen}{#1}}}
\newtcolorbox{exemplebox}[1][]{popbase,colframe=poporange,before upper={\popboxhead{Exemple}{poporange}{#1}}}
\newtcolorbox{definition}[1][]{popbase,colframe=popblue,before upper={\popboxhead{Definition}{popblue}{#1}}}
\newtcolorbox{propriete}[1][]{popbase,colframe=poppurple,before upper={\popboxhead{Property}{poppurple}{#1}}}
\newtcolorbox{methode}[1][]{popbase,colframe=popgold,before upper={\popboxhead{Method}{popgold}{#1}}}
\newtcolorbox{attention}[1][]{popbase,colframe=poppink,before upper={\popboxhead{Warning}{poppink}{#1}}}
\newtcolorbox{experience}[1][]{popbase,colframe=popblue,colback=popblueL,before upper={\popboxhead{Experiment}{popblue}{#1}}}
% « Le savais-tu ? » : carte violette pleine (contexte, culture, Cameroun)
\newtcolorbox{savaistu}[1][]{popbase,colback=poppurple,colframe=popink,
  fontupper=\popsemi\fontsize{10.4}{14.2}\selectfont\color{white},
  before upper={\poppill{Did you know?}{white}{poppurple}\ifx\relax#1\relax\else\hspace{6pt}{\popxbold\color{white}#1}\fi\par\vspace{7pt}}}
% Exercice résolu : énoncé en haut, \tcblower puis la solution sur fond crème
\newcounter{popexo}\newcounter{popexr}
\newtcolorbox{exoresolu}[1][]{popbase,colframe=poporange,colbacklower=popcream,
  segmentation style={popink,line width=1.2pt,dashed},
  before upper={\stepcounter{popexr}\popboxhead{Worked example \thepopexr}{poporange}{#1}},
  before lower={\poppill{Solution}{popink}{popcream}\par\vspace{6pt}}}
% Exercice (non résolu en place) — la correction est dans la partie Corrigés
\newtcolorbox{exercice}[1][]{popbase,colframe=popink,
  before upper={\stepcounter{popexo}\popboxhead{Exercise \thepopexo}{popink}{#1}}}
% Corrigé
\newtcolorbox{corrige}[1][]{popbase,colframe=popgreen,colback=popgreenL,
  before upper={\popboxhead{Solution}{popgreen}{#1}}}
% Objectifs / prérequis / bilan
\newtcolorbox{objectifs}{popbase,colframe=popink,colback=poporangeL,
  before upper={\popboxhead{Chapter objectives}{popink}{}}}
\newtcolorbox{prerequis}{popbase,colframe=popink,colback=popblueL,
  before upper={\popboxhead{Prerequisites}{popblue}{}}}
\newtcolorbox{bilan}{popbase,colframe=popink,colback=white,boxrule=2.8pt,
  before upper={\popboxhead{Fiche bilan}{poporange}{L'essentiel à connaître pour l'examen}}}
% Figure encadrée avec légende
\newtcolorbox{popfigure}[1][]{enhanced,breakable=false,colback=white,colframe=popline,boxrule=1.4pt,arc=8pt,
  left=10pt,right=10pt,top=10pt,bottom=8pt,before skip=10pt,after skip=12pt,halign=center,
  lower separated=false,halign lower=center,
  fontlower=\popsemi\fontsize{9}{11.5}\selectfont\color{popmuted},
  #1}
\newcommand{\legende}[1]{\par\vspace{6pt}{\popsemi\fontsize{9}{11.5}\selectfont\color{popmuted}#1}}

% ============================================================
% Signature OnBuch+
% ============================================================
\newcommand{\poplogoinline}[1]{{\poplogo\fontsize{#1}{#1}\selectfont\color{popink}OnBuch\color{poporange}+}}
\newcommand{\signatureonbuch}{%
  \begin{tcolorbox}[enhanced,colback=popink,colframe=popink,boxrule=2.2pt,arc=10pt,
      drop shadow=poporange,left=14pt,right=14pt,top=10pt,bottom=10pt,before skip=0pt,after skip=0pt]
    \tikz[baseline=-0.7ex]\node[circle,fill=popgreen,draw=popcream,line width=1.3pt,minimum size=22pt,inner sep=0]{\popcheck{white}};%
    \hspace{9pt}\begin{minipage}[c]{0.62\linewidth}
      {\popxbold\fontsize{12}{13}\selectfont\color{popcream}Signed\ \textbullet\ {\poplogo OnBuch\color{poporange}+}}\par\vspace{2pt}
      {\raggedright\popsemi\fontsize{8}{10.5}\selectfont\color{popline}\MakeUppercase{OnBuch+ teaching team}\\\MakeUppercase{Follows the official MINESEC syllabus}\par}
    \end{minipage}\hfill
    {\poplogo\fontsize{22}{22}\selectfont\color{popcream}OnBuch\color{poporange}+}
  \end{tcolorbox}%
}

% ============================================================
% Page de garde
% #1 titre · #2 matière · #3 niveau · #4 module · #5 n° leçon · #6 durée ·
% #7 description / objectifs (texte ou itemize)
% ============================================================
\newcommand{\pagegarde}[7]{%
  \thispagestyle{empty}
  \newgeometry{a4paper,margin=1.7cm,top=1.6cm,bottom=1.4cm}%
  \begin{tikzpicture}[remember picture,overlay]
    \fill[poporange!18] ([xshift=-3.2cm,yshift=-3.6cm]current page.north east) circle (4.6cm);
    \fill[poppurple!14] ([xshift=2.4cm,yshift=7.6cm]current page.south west) circle (3.8cm);
  \end{tikzpicture}%
  {\poplogo\fontsize{30}{30}\selectfont\color{popink}OnBuch\color{poporange}+}\hfill
  \raisebox{0.6ex}{\poppill{Signed course \textbullet\ Premium}{popink}{popcream}}\par
  \vspace{1pt}\popsquiggle{poppurple}\par
  \vspace{8pt}
  \begin{tcolorbox}[enhanced,colback=poporange,colframe=popink,boxrule=2.8pt,arc=13pt,
      drop shadow=popink,left=18pt,right=18pt,top=12pt,bottom=13pt,after skip=10pt]
    \poppill{#2 \textbullet\ #3}{popink}{popcream}\hspace{5pt}\poppill{#4}{white}{popink}\par\vspace{8pt}
    {\popdisplay\fontsize{14}{14}\selectfont\color{popink}CHAPTER #5}\par\vspace{4pt}
    {\raggedright\hyphenpenalty=10000\exhyphenpenalty=10000\popdisplay\fontsize{25}{27}\selectfont\color{popcream}\MakeUppercase{#1}\par}
    \vspace{8pt}
    {\popxbold\fontsize{9.5}{9.5}\selectfont\color{popink}\MakeUppercase{Suggested time: #6}}
  \end{tcolorbox}
  \begin{tcolorbox}[enhanced,colback=white,colframe=popink,boxrule=2.2pt,arc=10pt,
      drop shadow=popink,left=16pt,right=16pt,top=10pt,bottom=10pt,after skip=8pt]
    \poppill{In this chapter}{poppurple}{white}\par\vspace{5pt}
    \popsemi\fontsize{9.3}{12.6}\selectfont\color{popink}#7
  \end{tcolorbox}
  \noindent\begin{minipage}[t]{0.487\linewidth}
    \begin{tcolorbox}[enhanced,colback=popgreen,colframe=popink,boxrule=2.2pt,arc=10pt,
        drop shadow=popink,left=11pt,right=11pt,top=10pt,bottom=10pt,height=3.1cm,valign=center]
      \tcbox[colback=white,colframe=popink,boxrule=1.3pt,arc=5pt,boxsep=0pt,nobeforeafter,
        left=3pt,right=3pt,top=3pt,bottom=3pt,valign=center]{\qrcode[height=1.9cm]{https://play.google.com/store/apps/details?id=com.luvvix.onbuch&hl=en}}%
      \hspace{8pt}\begin{minipage}[c]{0.5\linewidth}
        {\popdisplay\fontsize{11}{12.5}\selectfont\color{white}\MakeUppercase{Download OnBuch}}\par\vspace{4pt}
        {\raggedright\popsemi\fontsize{8.4}{11}\selectfont\color{white}Free \textbullet\ Google Play\\AI tutor, past papers, results\par}
      \end{minipage}
    \end{tcolorbox}
  \end{minipage}\hfill
  \begin{minipage}[t]{0.487\linewidth}
    \begin{tcolorbox}[enhanced,colback=poppurple,colframe=popink,boxrule=2.2pt,arc=10pt,
        drop shadow=popink,left=11pt,right=11pt,top=10pt,bottom=10pt,height=3.1cm,valign=center]
      \tikz[baseline=-0.7ex]\node[circle,fill=white,draw=popink,line width=1.3pt,minimum size=1.6cm,inner sep=0]{\popdisplay\fontsize{24}{24}\selectfont\color{poppurple}+};%
      \hspace{8pt}\begin{minipage}[c]{0.6\linewidth}
        {\popdisplay\fontsize{11}{12.5}\selectfont\color{white}\MakeUppercase{Go OnBuch+}}\par\vspace{4pt}
        {\raggedright\popsemi\fontsize{8.4}{11}\selectfont\color{white}All signed courses, unlimited AI tutor, PDF sheets — OnBuch+ tab in the app\par}
      \end{minipage}
    \end{tcolorbox}
  \end{minipage}
  \vspace{8pt}
  \popcontenu
  \vfill
  \signatureonbuch
  \restoregeometry
  \newpage
}

% Rangée « Ce cours contient » (page de garde)
\newcommand{\poptile}[3]{% #1 couleur #2 gros texte #3 légende
  \begin{tcolorbox}[enhanced,colback=#1,colframe=popink,boxrule=2pt,arc=9pt,shadow={2.5pt}{-2.5pt}{0pt}{popink},
      left=6pt,right=6pt,top=9pt,bottom=9pt,halign=center,valign=center,height=1.75cm,width=\linewidth,nobeforeafter]
    {\popdisplay\fontsize{12.5}{13}\selectfont\color{white}\MakeUppercase{#2}}\par\vspace{3pt}
    {\centering\popsemi\fontsize{7.8}{9.5}\selectfont\color{white}#3\par}
  \end{tcolorbox}}
\newcommand{\popcontenu}{%
  \par\noindent{\popxboldls\fontsize{7.5}{7.5}\selectfont\color{popmuted}THIS SIGNED COURSE CONTAINS}\par\vspace{7pt}
  \noindent\begin{minipage}[t]{0.235\linewidth}\poptile{popblue}{Course}{complete, illustrated, step by step}\end{minipage}\hfill
  \begin{minipage}[t]{0.235\linewidth}\poptile{poporange}{Methods}{and worked examples}\end{minipage}\hfill
  \begin{minipage}[t]{0.235\linewidth}\poptile{poppink}{Exercises}{exam style + solutions}\end{minipage}\hfill
  \begin{minipage}[t]{0.235\linewidth}\poptile{popgreen}{Summary sheet}{the essentials to remember}\end{minipage}\par}

% Sommaire
\newtcolorbox{sommaire}{enhanced,colback=white,colframe=popink,boxrule=2.2pt,arc=10pt,
  drop shadow=popink,left=18pt,right=18pt,top=13pt,bottom=13pt,before skip=6pt,after skip=20pt,
  before upper={\poppill{Contents}{poppurple}{white}\par\vspace{9pt}\popsemi\fontsize{10.6}{14.5}\selectfont\color{popink}}}

% Signature de fin de cours — poussée tout en bas de la dernière page
\newcommand{\findecours}{%
  \par\vfill\nopagebreak
  \begin{center}\popsquiggle{poporange}\end{center}\vspace{4pt}
  \signatureonbuch
}
\AtBeginDocument{\def\llbracket{\mathopen{[\mkern-3.5mu[}}\def\rrbracket{\mathclose{]\mkern-3.5mu]}}} % intervalles d'entiers (glyphe ⟦ absent de la police)
% Glyphes absents de Fira Math : redéfinis à partir de glyphes disponibles
\AtBeginDocument{%
  \def\setminus{\mathbin{\backslash}}\let\smallsetminus\setminus
  \def\vdots{\mathord{\vbox{\baselineskip3.2pt\lineskiplimit0pt\kern4pt\hbox{.}\hbox{.}\hbox{.}}}}%
  \def\ddots{\mathinner{\mkern1mu\raise6.5pt\hbox{.}\mkern2mu\raise3.6pt\hbox{.}\mkern2mu\raise0.7pt\hbox{.}\mkern1mu}}%
  \def\triangle{\mathord{\scalebox{0.9}{$\symup{\Delta}$}}}%
  \def\nabla{\mathord{\raisebox{\depth}{\scalebox{1}[-1]{$\symup{\Delta}$}}}}%
  \def\checkmark{\popcheck{popdark}}%
  \def\star{\mathbin{\ast}}\def\bigstar{\mathord{\ast}}%
  \def\diamond{\mathbin{\scalebox{0.6}{$\Diamond$}}}%
  \def\bigcup{\mathop{\mathchoice{\raisebox{-0.3em}{\scalebox{1.7}{$\cup$}}}{\scalebox{1.25}{$\cup$}}{\cup}{\cup}}\displaylimits}%
  \def\bigcap{\mathop{\mathchoice{\raisebox{-0.3em}{\scalebox{1.7}{$\cap$}}}{\scalebox{1.25}{$\cap$}}{\cap}{\cap}}\displaylimits}%
  \def\lesseqqgtr{\mathrel{\lessgtr}}\def\bigoplus{\mathop{\oplus}\displaylimits}%
}
\providecommand{\ssi}{if and only if} % abréviation fréquente
\AtBeginDocument{\providecommand{\wideparen}{\overparen}} % arc de cercle

%% FIGFIX-BEGIN v4 — correctifs globaux des figures (docs/onbuch-generation/tools/figfix)
\usepackage{adjustbox}\usepackage{etoolbox}
% toute figure TikZ/pgfplots plus large que la ligne est réduite à la largeur disponible
\BeforeBeginEnvironment{tikzpicture}{\begin{adjustbox}{max width=\linewidth}}
\AfterEndEnvironment{tikzpicture}{\end{adjustbox}}
% légendes pgfplots lisibles par-dessus les courbes
\pgfplotsset{every axis/.append style={legend style={fill=white,fill opacity=0.92,text opacity=1,draw=black!15,inner sep=3pt}}}
% étiquettes d'arêtes (syntaxe edge["..."]) avec fond blanc
\tikzset{every edge quotes/.append style={fill=white,inner sep=1.2pt}}
%% FIGFIX-END
\begin{document}

\pagegarde{Electrostatics and Electric Fields}{Physics}{Upper Sixth}{Upper Sixth — Physics}{3}{2 periods}{This chapter introduces electrostatics: the origin and types of electric charge, the behaviour of conductors and insulators, and the methods of charging a body. It then establishes Coulomb's law and the concept of electric field, the essential tools for analysing electrical interactions at GCE Advanced Level.
\vspace{4pt}
\begin{multicols}{2}\small
\begin{itemize}
\item By the end of this chapter I can explain the atomic origin of electric charge and distinguish positive from negative charge.
\item By the end of this chapter I can state and apply the basic law of electrostatics (attraction and repulsion between charges).
\item By the end of this chapter I can distinguish conductors from insulators and explain earthing.
\item By the end of this chapter I can describe and explain the three methods of charging: friction, contact and induction.
\item By the end of this chapter I can explain point action and describe the construction and role of a lightning conductor.
\item By the end of this chapter I can state Coulomb's law, use it quantitatively, and show its inverse-square nature.
\end{itemize}\end{multicols}}

\begin{sommaire}
\begin{multicols}{2}
\begin{enumerate}
\item Electric Charge and the Basic Law of Electrostatics
\item Conductors, Insulators and Methods of Charging
\item Point Action and the Lightning Conductor
\item Coulomb's Law and the Electric Field
\item Integration activity
\item Methods and worked examples
\item Exercises
\item Solutions to the exercises
\item Summary sheet
\end{enumerate}
\end{multicols}
\end{sommaire}

\begin{objectifs}
\begin{itemize}
\item By the end of this chapter I can explain the atomic origin of electric charge and distinguish positive from negative charge.
\item By the end of this chapter I can state and apply the basic law of electrostatics (attraction and repulsion between charges).
\item By the end of this chapter I can distinguish conductors from insulators and explain earthing.
\item By the end of this chapter I can describe and explain the three methods of charging: friction, contact and induction.
\item By the end of this chapter I can explain point action and describe the construction and role of a lightning conductor.
\item By the end of this chapter I can state Coulomb's law, use it quantitatively, and show its inverse-square nature.
\item By the end of this chapter I can explain how the medium between charges affects the electrostatic force (permittivity).
\item By the end of this chapter I can define electric field strength E, calculate it for point charges, and represent field lines.
\item By the end of this chapter I can solve GCE-style problems combining superposition of forces and fields from several charges.
\end{itemize}
\end{objectifs}

\begin{prerequis}
\begin{itemize}
\item Structure of the atom: protons, neutrons, electrons (Lower Sixth / Form 5 chemistry and physics)
\item Forces, vector addition and Newton's laws (Lower Sixth mechanics)
\item Inverse-square relationships and manipulation of powers of ten in calculations
\item The elementary charge e = 1.6 × 10⁻¹⁹ C and SI units
\item Basic graph interpretation (plotting F against 1/r²)
\end{itemize}
\end{prerequis}

\newpage

\coursec{Electric Charge and the Basic Law of Electrostatics}

During the rainy season in Bamenda or Douala, lightning regularly strikes tall buildings, radio masts and isolated trees, sometimes destroying electrical installations worth millions of francs CFA. Yet a simple pointed metal rod fixed on a roof and connected to the ground protects a whole house: this is the lightning conductor, a direct application of electrostatics. Why does a rubbed plastic ruler attract small pieces of paper? Why do sparks jump when you remove a nylon pullover on a dry evening in Bafoussam? All these phenomena come from one fundamental entity: \cle{electric charge}. In this chapter we uncover the nature of charge, the forces between charges, and the invisible field that surrounds them.

\courssub{Historical and Everyday Observations}

More than two thousand years ago, the Greeks observed that a piece of \emph{amber} (a fossilised tree resin; the Greek word is \emph{elektron}) rubbed with fur acquires the power to attract light objects such as straw, feathers or small seeds. For centuries this remained a curiosity. Today you can reproduce the same effect in a few seconds at your desk:

\begin{itemize}
\item Rub a plastic ruler or biro vigorously on your sleeve or on dry hair, then bring it near small pieces of paper: the pieces jump up and stick to the ruler.
\item Run a plastic comb through dry hair, then hold it near a thin stream of water from a tap: the stream bends towards the comb.
\item Pull off a nylon pullover in the dark on a dry evening: you hear faint crackling and may even see tiny sparks.
\end{itemize}

\begin{popfigure}
\begin{center}
\begin{tikzpicture}[scale=1.0]
% rod
\draw[fill=popblueL, draw=popblue, thick, rotate around={-20:(0,0)}] (-0.3,-0.25) rectangle (4.2,0.25);
\node[popdark] at (1.95,0.75) {rubbed plastic rod ($-$)};
% minus signs on rod
\foreach \x in {0.3,1.0,1.7,2.4,3.1,3.8} \node[popblue] at (\x,-0.62) {$-$};
% hand motion arrow
\draw[->, thick, poporange] (-1.3,0.9) -- (-0.2,0.45) node[midway, above left, popdark]{\small rub};
% pieces of paper
\foreach \p/\a in {(1.0,-2.2)/15,(1.9,-2.5)/-20,(2.8,-2.1)/30,(3.6,-2.4)/-10}
  \draw[fill=white, draw=popmuted, rotate around={\a:\p}] \p rectangle ++(0.45,0.3);
% attraction arrows
\draw[->, thick, popgreen] (1.2,-1.85) -- (1.15,-1.05);
\draw[->, thick, popgreen] (2.05,-2.15) -- (2.0,-1.1);
\draw[->, thick, popgreen] (2.9,-1.75) -- (2.85,-1.05);
\draw[->, thick, popgreen] (3.7,-2.05) -- (3.6,-1.1);
\node[popgreen] at (4.6,-1.6) {\small attraction};
% table
\draw[thick, popmuted] (0.2,-2.65) -- (4.6,-2.65);
\end{tikzpicture}
\end{center}
\legende{A plastic rod rubbed on cloth attracts small neutral pieces of paper.}
\end{popfigure}

In all these cases, \emph{rubbing} (friction between two different materials) produces a new state of the bodies: we say they have become \cle{electrified} or \cle{charged}. The branch of physics that studies charges at rest and the forces between them is called \cle{electrostatics}.

\begin{definition}[Electric charge]
\cle{Electric charge} is the fundamental property of matter responsible for electrostatic interactions (attractions and repulsions between rubbed bodies). Charge is a scalar physical quantity, denoted $q$ (or $Q$). Its SI unit is the \cle{coulomb}, symbol $\mathrm{C}$.
\end{definition}

\courssub{The Atomic Origin of Charge}

Where does the charge come from? Rubbing does not \emph{create} anything; it \emph{transfers} something. To understand what, recall the structure of the atom:

\begin{itemize}
\item a tiny central \textbf{nucleus}, containing \textbf{protons} (positively charged) and \textbf{neutrons} (uncharged);
\item \textbf{electrons} (negatively charged) moving around the nucleus, held by electrical attraction.
\end{itemize}

In a neutral atom, the number of electrons exactly equals the number of protons, so the total charge is zero. Protons are locked inside the nucleus and essentially never move in solid materials. Electrons, however — especially the outermost ones — can be torn away by friction.

\begin{propriete}[Electrification by friction]
When two different insulating materials are rubbed together, \cle{electrons} are transferred from one material to the other:
\begin{itemize}
\item the body that \textbf{gains} electrons becomes \textbf{negatively} charged;
\item the body that \textbf{loses} electrons becomes \textbf{positively} charged.
\end{itemize}
A charged body is therefore simply a body with an \emph{excess} or a \emph{deficit} of electrons. Charging never creates or destroys charge; it only moves electrons from one body to another.
\end{propriete}

\begin{exemplebox}[The ruler and the paper]
When you rub a plastic ruler on a woollen sleeve, electrons pass from the wool to the plastic. The ruler carries an excess of electrons (negative charge); the sleeve, having lost electrons, carries an equal positive charge. The two charges are equal in magnitude: total charge is unchanged.
\end{exemplebox}

\courssub{The Two Types of Charge}

Systematic experiments in the 18th century showed that there are exactly \textbf{two} kinds of electric charge, conventionally named:

\begin{itemize}
\item \textbf{positive charge} ($+$): the charge acquired by \emph{glass} rubbed with \emph{silk} (the glass loses electrons);
\item \textbf{negative charge} ($-$): the charge acquired by \emph{ebonite} or \emph{plastic} rubbed with \emph{fur} (the ebonite gains electrons).
\end{itemize}

Every charged body ever observed behaves either like rubbed glass or like rubbed ebonite: there is no third type. The choice of which is ``positive'' is a historical convention (Benjamin Franklin); the physics is unchanged if the signs are swapped everywhere.

\courssub{The Basic Law of Electrostatics}

The fundamental experimental fact governing the interaction between charges is the \cle{basic law of electrostatics}.

\begin{propriete}[Basic law of electrostatics]
Two electric charges of the \textbf{same sign repel} each other; two charges of \textbf{opposite sign attract} each other. In symbols:
\[ (+)\ \text{repels}\ (+), \qquad (-)\ \text{repels}\ (-), \qquad (+)\ \text{attracts}\ (-). \]
\end{propriete}

\begin{popfigure}
\begin{center}
\begin{tikzpicture}[scale=1.0]
% ===== left: repulsion =====
\draw[thick, popmuted] (-0.6,2.6) -- (3.6,2.6);
\draw[thick] (0.7,2.6) -- (0.7,1.5);
\draw[thick] (2.3,2.6) -- (2.3,1.5);
\draw[fill=popblueL, draw=popblue, thick] (0.7,1.2) circle (0.3);
\draw[fill=popblueL, draw=popblue, thick] (2.3,1.2) circle (0.3);
\node[popblue] at (0.7,1.2) {$+$};
\node[popblue] at (2.3,1.2) {$+$};
\draw[->, very thick, poporange] (0.35,1.2) -- (-0.35,1.2);
\draw[->, very thick, poporange] (2.65,1.2) -- (3.35,1.2);
\node[popdark] at (1.5,0.35) {\small like charges repel};
% ===== right: attraction =====
\begin{scope}[xshift=6.5cm]
\draw[thick, popmuted] (-0.6,2.6) -- (3.6,2.6);
\draw[thick] (0.7,2.6) -- (1.05,1.55);
\draw[thick] (2.3,2.6) -- (1.95,1.55);
\draw[fill=popblueL, draw=popblue, thick] (1.05,1.25) circle (0.3);
\draw[fill=poppink!25, draw=poppink, thick] (1.95,1.25) circle (0.3);
\node[popblue] at (1.05,1.25) {$+$};
\node[poppink] at (1.95,1.25) {$-$};
\draw[->, very thick, popgreen] (1.05,1.6) -- (1.05,2.25);
\draw[->, very thick, popgreen] (1.95,1.6) -- (1.95,2.25);
\node[popdark] at (1.5,0.35) {\small unlike charges attract};
\end{scope}
\end{tikzpicture}
\end{center}
\legende{Two light charged spheres suspended by threads. Left: same sign — the threads diverge (repulsion). Right: opposite signs — the spheres swing together (attraction).}
\end{popfigure}

\begin{attention}[Attraction is not a proof of opposite charge]
If a charged rod \emph{repels} a light suspended ball, the ball is certainly charged \emph{with the same sign} as the rod: repulsion is conclusive. But if the rod \emph{attracts} the ball, the ball may be oppositely charged \textbf{or simply neutral}: a charged body always attracts a nearby neutral light object by \cle{electrostatic induction} (it redistributes the charges inside the neutral object, bringing opposite charges closer). Attraction alone therefore proves nothing — only repulsion is a reliable test of charge.
\end{attention}

\courssub{Quantisation and Conservation of Charge}

Careful measurements (most famously Millikan's oil-drop experiment, 1909--1913) show that charge never appears in arbitrary amounts. Every charge found in nature is a whole-number multiple of one smallest unit, the \cle{elementary charge} $e$.

\begin{propriete}[Quantisation of charge]
Electric charge is \cle{quantised}: the charge $q$ of any body is an integer multiple of the elementary charge
\[ q = n e, \qquad n \in \mathbb{Z}, \qquad e = 1.60\times 10^{-19}\ \mathrm{C}. \]
The proton carries charge $+e$; the electron carries charge $-e$. The integer $n$ is the number of electrons gained ($n<0$) or lost ($n>0$) by the body.
\end{propriete}

\begin{propriete}[Conservation of charge]
In any \textbf{isolated system} (one that exchanges no charge with its surroundings), the total electric charge remains constant:
\[ \sum q_{\text{before}} = \sum q_{\text{after}}. \]
Charge can be transferred or redistributed, but never created or destroyed. This is one of the great conservation laws of physics, on the same footing as the conservation of energy.
\end{propriete}

\begin{exoresolu}[How many electrons make a charged ruler?]
A plastic ruler rubbed on cloth acquires a charge $q = -8.0\times 10^{-9}\ \mathrm{C}$. (a) How many excess electrons does the ruler carry? (b) What is the charge of the cloth after rubbing, assuming both were initially neutral?
\tcblower
\textbf{(a)} Charge is quantised: $q = ne$ with $n<0$ since electrons were gained. Hence
\[ |n| = \frac{|q|}{e} = \frac{8.0\times 10^{-9}}{1.60\times 10^{-19}} = 5.0\times 10^{10}\ \text{electrons}. \]
That is fifty thousand million excess electrons — an enormous number, yet the charge is tiny on the scale of the coulomb.

\textbf{(b)} The system \{ruler + cloth\} is isolated and initially neutral ($q_{\text{total}} = 0$). By conservation of charge,
\[ q_{\text{cloth}} = -q_{\text{ruler}} = +8.0\times 10^{-9}\ \mathrm{C}. \]
The cloth is positively charged, having lost exactly the electrons gained by the ruler.
\end{exoresolu}

\begin{attention}[Common mistakes with $q = ne$]
\begin{itemize}
\item Do not write $n = q \times e$: the number of electrons is $n = q/e$.
\item $n$ must be an integer; if a calculation gives a non-integer, either the data or the arithmetic is wrong.
\item Keep the \emph{sign} of $q$: a negative $q$ means electrons were \emph{gained}, a positive $q$ means electrons were \emph{lost}.
\item The coulomb is a large unit: everyday static charges are of order $10^{-9}$ to $10^{-6}\ \mathrm{C}$, never several coulombs.
\end{itemize}
\end{attention}

\courssub{Conductors, Insulators and Methods of Charging}

Not all materials respond to rubbing in the same way. The distinction is crucial for understanding how charges move and how we can manipulate them.

\begin{definition}[Conductors and insulators]
A \cle{conductor} is a material that allows electric charges to move freely through it (e.g.\ metals, graphite, electrolytes). In metals, the ``conduction electrons'' are delocalised and form a mobile electron gas.
An \cle{insulator} (or dielectric) is a material in which charges are tightly bound to their atoms and cannot move freely (e.g.\ plastic, glass, dry wood, rubber, air).
\end{definition}

\begin{propriete}[Behaviour of charge on conductors and insulators]
\begin{itemize}
\item On a \textbf{conductor}, any excess charge resides entirely on the \emph{outer surface} and distributes itself so that the electric field inside the material is zero (electrostatic equilibrium).
\item On an \textbf{insulator}, charge stays \emph{where it is put}; it does not spread over the surface.
\end{itemize}
\end{propriete}

There are three principal ways to charge a body:

\begin{description}
\item[\textbf{1. Charging by friction (triboelectric effect).}] Rubbing two different insulators transfers electrons from one to the other (as seen with the plastic ruler). Both bodies become charged with equal and opposite charges.
\item[\textbf{2. Charging by contact (conduction).}] Touching a neutral conductor with a charged conductor allows charge to flow until both reach the same potential. They end up with charges of the \emph{same sign}. If the conductors are identical, they share the total charge equally.
\item[\textbf{3. Charging by induction.}] Bringing a charged object near a neutral conductor (without touching) causes a separation of charges in the conductor: the near side acquires an induced charge opposite to the inducer, the far side acquires the same sign. If the conductor is then \emph{earthed} (connected to ground) while the inducer is present, charge flows to/from the earth. Breaking the earth connection \emph{before} removing the inducer leaves the conductor with a net charge \emph{opposite} to that of the inducer.
\end{description}

\begin{experience}[Charging a metal sphere by induction]
\begin{enumerate}
\item Bring a negatively charged rod near a neutral metal sphere on an insulating stand. Electrons in the sphere are repelled to the far side; the near side becomes positively charged.
\item Touch the far side of the sphere with your finger (earth connection). Electrons flow from the sphere to earth.
\item Remove your finger (break earth connection), \emph{then} remove the rod. The sphere is left with a net \textbf{positive} charge.
\end{enumerate}
This method charges the sphere without ever touching it with the inducing rod, and the induced charge is opposite in sign to the inducer.
\end{experience}

\courssub{Detecting Charge: The Gold-Leaf Electroscope}

How can we tell whether a body is charged, and with which sign? The classic instrument is the \cle{gold-leaf electroscope}. It consists of a metal rod ending in a metal disc (or knob) at the top and carrying two thin gold leaves at the bottom, enclosed in a glass case with earthed metal walls to shield it from draughts and stray influences.

\begin{popfigure}
\begin{center}
\begin{tikzpicture}[scale=0.95]
% ===== uncharged =====
% case
\draw[thick, popmuted] (-1.6,-2.4) rectangle (1.6,1.2);
\draw[thick, popmuted] (-1.6,1.2) -- (-1.6,1.7) (1.6,1.2) -- (1.6,1.7);
\draw[thick, popmuted] (-1.6,1.7) -- (1.6,1.7);
% rod and disc
\draw[thick, popdark] (0,2.6) -- (0,-1.6);
\draw[fill=popgold, draw=popdark, thick] (0,2.9) circle (0.3);
% leaves together
\draw[thick, popgold] (0,-1.6) -- (0.06,-2.25);
\draw[thick, popgold] (0,-1.6) -- (-0.06,-2.25);
\node[popdark] at (0,-3.0) {\small uncharged: leaves hang together};
% ===== charged =====
\begin{scope}[xshift=6cm]
\draw[thick, popmuted] (-1.6,-2.4) rectangle (1.6,1.2);
\draw[thick, popmuted] (-1.6,1.2) -- (-1.6,1.7) (1.6,1.2) -- (1.6,1.7);
\draw[thick, popmuted] (-1.6,1.7) -- (1.6,1.7);
\draw[thick, popdark] (0,2.6) -- (0,-1.6);
\draw[fill=popgold, draw=popdark, thick] (0,2.9) circle (0.3);
\node[popblue] at (0.55,2.9) {$-$};
% leaves apart
\draw[thick, popgold] (0,-1.6) -- (0.55,-2.2);
\draw[thick, popgold] (0,-1.6) -- (-0.55,-2.2);
\node[popblue] at (0.75,-2.3) {\small $-$};
\node[popblue] at (-0.75,-2.3) {\small $-$};
\node[popdark] at (0,-3.0) {\small charged: leaves diverge};
\end{scope}
\end{tikzpicture}
\end{center}
\legende{Gold-leaf electroscope. Left: uncharged, the leaves hang side by side. Right: charged, both leaves carry the same sign and repel each other.}
\end{popfigure}

\textbf{How it works.} When a charged body touches (or is brought near) the disc, charge spreads over the whole conducting system — disc, rod and leaves. The two leaves then carry charges of the \emph{same sign} and repel each other: they \textbf{diverge}. The greater the charge, the wider the divergence.

\begin{methode}[Testing for the presence and sign of a charge with an electroscope]
\begin{enumerate}
\item \textbf{Presence of charge.} Bring the body near (or into contact with) the disc of a neutral electroscope. If the leaves diverge, the body is charged; if nothing happens, it is neutral.
\item \textbf{Sign of an unknown charge.} First charge the electroscope with a \emph{known} sign (e.g.\ touch it with a rubbed plastic rod: the electroscope becomes negative).
\item Bring the unknown charged body \emph{near} the disc (without touching):
\begin{itemize}
\item if the leaves diverge \emph{more}, the body has the \emph{same} sign as the electroscope;
\item if the leaves \emph{collapse} (or diverge less), the body has the \emph{opposite} sign.
\end{itemize}
\item Interpret cautiously: remember that a neutral body can also cause slight movement by induction — large, clear effects are the reliable ones.
\end{enumerate}
\end{methode}

\begin{experience}[Investigating the basic law with simple apparatus]
Suspend two light aluminium-coated pith balls (or two small polystyrene balls wrapped in aluminium foil) side by side from insulating threads.
\begin{enumerate}
\item Touch both balls with the same rubbed plastic rod: they swing apart — both are negative, like charges repel.
\item Touch one ball with rubbed plastic (negative) and the other with glass rubbed with silk (positive): they swing together — unlike charges attract.
\item Bring a charged rod near one untouched ball: it is attracted, although neutral — evidence of induction, not of charge on the ball.
\end{enumerate}
These three simple observations contain the whole of the basic law of electrostatics, together with its most subtle trap.
\end{experience}

\begin{savaistu}[Why lightning strikes]
During a thunderstorm, friction between air currents, water droplets and ice crystals inside a cloud separates enormous quantities of charge: the base of the cloud becomes strongly charged (usually negative), while the ground below acquires an induced opposite charge. When the accumulated charge becomes too great, the air itself becomes conducting (dielectric breakdown, $\approx 3\times 10^6\ \mathrm{V/m}$) and a gigantic spark — lightning — neutralises the charges. The same physics, scaled down by a factor of millions, produces the tiny spark you see when removing a nylon pullover on a dry evening in Bafoussam.
\end{savaistu}

\begin{aretenir}[Key points of this section]
\begin{itemize}
\item Rubbing transfers \textbf{electrons}: a body becomes negative by gaining electrons, positive by losing them.
\item There are exactly \textbf{two types} of charge; the SI unit of charge is the \textbf{coulomb} (C).
\item \textbf{Basic law:} like charges repel, unlike charges attract.
\item Charge is \textbf{quantised}: $q = ne$, with $e = 1.60\times 10^{-19}\ \mathrm{C}$.
\item Charge is \textbf{conserved} in any isolated system.
\item \textbf{Conductors} allow charge movement; \textbf{insulators} do not. Charging methods: friction, contact, induction.
\item The \textbf{gold-leaf electroscope} detects charge and, when pre-charged with a known sign, identifies the sign of an unknown charge.
\item \textbf{Repulsion} proves a body is charged; \textbf{attraction} does not (induction can attract neutral objects).
\end{itemize}
\end{aretenir}

\begin{exercice}[Checking your understanding]
\begin{enumerate}
\item A glass rod rubbed with silk acquires a charge of $+3.2\times 10^{-8}\ \mathrm{C}$. Did the rod gain or lose electrons? How many?
\item Two identical conducting spheres carry charges $+6.0\times 10^{-8}\ \mathrm{C}$ and $-2.0\times 10^{-8}\ \mathrm{C}$. They are brought into contact and then separated. Using conservation of charge, find the charge on each sphere after separation.
\item A student brings a negatively charged rod near a suspended light ball and observes attraction. She concludes the ball is positively charged. Is her conclusion justified? Explain, and describe an experiment that would settle the question.
\item A neutral metal sphere on an insulating stand is charged by induction using a positively charged rod. Describe the steps and state the final charge on the sphere.
\end{enumerate}
\end{exercice}

\begin{corrige}[Exercise 1]
\begin{enumerate}
\item The rod is positive, so it \textbf{lost} electrons. Number lost:
\[ n = \frac{q}{e} = \frac{3.2\times 10^{-8}}{1.60\times 10^{-19}} = 2.0\times 10^{11}\ \text{electrons}. \]
\item The system of the two spheres is isolated, so total charge is conserved:
\[ q_{\text{total}} = (+6.0 - 2.0)\times 10^{-8} = +4.0\times 10^{-8}\ \mathrm{C}. \]
Identical conducting spheres share the charge equally on separation:
\[ q_1 = q_2 = +2.0\times 10^{-8}\ \mathrm{C}. \]
\item No. The ball could be positive \emph{or neutral} (attracted by induction). To decide, bring the rod so as to test for \textbf{repulsion}: charge the ball deliberately (or use a positively charged rod as well). If a rod of known sign repels the ball, the ball carries that same sign; attraction alone can never prove it.
\item Steps: (1) Bring the $+$ rod near the sphere (electrons attracted to near side, far side becomes $+$). (2) Earth the far side (electrons flow from earth to sphere). (3) Break earth connection. (4) Remove rod. The sphere is left with a net \textbf{negative} charge.
\end{enumerate}
\end{corrige}

\coursec{Conductors, Insulators and Methods of Charging}

In the previous section we established that matter carries two kinds of electric charge and that like charges repel while unlike charges attract. This fundamental rule is the \cle{basic law of electrostatics}. But a rubbed balloon sticks to a wall while a rubbed metal spoon held in the hand does not. Why does charge \emph{stay} on the balloon but seem to \emph{disappear} from the spoon? The answer lies in a fundamental distinction between materials: in some of them electric charge can move freely, in others it is locked in place. This distinction governs everything we do in electrostatics, from charging a gold-leaf electroscope in the school laboratory to protecting a fuel tanker on the Douala--Yaound\'e road.

\begin{propriete}[Basic law of electrostatics]
Like charges \textbf{repel} each other; unlike charges \textbf{attract} each other. The force acts along the line joining the charges.
\end{propriete}

\courssub{Conductors and Insulators}

\begin{definition}[Conductor and insulator]
A \cle{conductor} is a material that contains electric charges free to move throughout its volume. In metals, these mobile charges are \cle{free electrons}: each atom of the metal releases one or more of its outer electrons, which then wander through the whole solid like a gas. In electrolytes (e.g. salt solutions), the mobile charges are positive and negative ions. An \cle{insulator} (or dielectric) is a material whose electrons remain tightly \cle{bound} to their atoms; charge placed on an insulator cannot flow through it.
\end{definition}

\begin{center}
\begin{tabularx}{\linewidth}{|l|X|X|}
\hline
\rowcolor{popblueL} & \textbf{Conductors} & \textbf{Insulators} \\
\hline
Electron behaviour & Free electrons (or ions) move through the whole material & Electrons bound to individual atoms \\
\hline
Common examples & Metals (copper, aluminium, iron), graphite, the human body, damp earth, salt solutions & Glass, dry wood, plastics, ebonite, rubber, dry air, paper, ceramics \\
\hline
Effect of charging & Charge spreads instantly over the surface & Charge stays where it was deposited \\
\hline
\end{tabularx}
\end{center}

\begin{attention}[The human body is a conductor]
Because the body conducts, a charged metal object held in the bare hand loses its charge through you to the ground. This is why electrostatics experiments use insulating handles (glass or plastic rods), and why a rubbed balloon (an insulator) keeps its charge while a rubbed spoon does not. Note also that \emph{dry} air is an insulator but \emph{damp} air conducts slightly: electrostatics experiments fail on humid rainy-season days in Cameroon because charge leaks away through the moist air.
\end{attention}

\courssub{Charge Distribution on Conductors and Insulators}

Suppose we deposit an excess of electrons on a copper sphere. Every electron repels every other electron, and because they are free to move, they push each other as far apart as possible. The furthest they can go is the \emph{outer surface} of the sphere. There they stop, spread uniformly, and the interior of the metal remains neutral.

\begin{propriete}[Charge resides on the outer surface]
On a conductor at \cle{electrostatic equilibrium} (no charge in motion), the net charge lies entirely on the \cle{outer surface} of the conductor; the interior carries no net charge. On an insulator, by contrast, charge remains \cle{localised} exactly where it was deposited, because it cannot move.
\end{propriete}

This property is a direct consequence of the repulsion between like charges combined with their mobility; it will be justified more formally when we study the electric field inside a conductor.

\begin{popfigure}
\begin{tikzpicture}[scale=1.0]
% Conductor
\draw[thick, popblue, fill=popblueL] (0,0) circle (1.3);
\foreach \a in {0,30,60,90,120,150,180,210,240,270,300,330}
  \node[popdark] at (\a:1.42) {\scriptsize $-$};
\node[align=center, text width=3.2cm] at (0,-2.1) {\small Conductor: charge spreads over the whole surface};
% Insulator
\begin{scope}[xshift=6.5cm]
\draw[thick, poporange, fill=popgold!30] (0,0) circle (1.3);
\foreach \p in {(0.75,0.55),(0.95,0.2),(0.55,0.9),(1.0,0.75),(0.7,0.1)}
  \node[popdark] at \p {\scriptsize $-$};
\node[align=center, text width=3.2cm] at (0,-2.1) {\small Insulator: charge stays where deposited};
\end{scope}
\end{tikzpicture}
\legende{Excess negative charge on a conducting sphere (left) spreads uniformly over the outer surface; on an insulating sphere (right) it remains localised at the point of deposit.}
\end{popfigure}

\courssub{Charging by Friction}

When two different insulating materials are rubbed together, the contact and the mechanical work tear some electrons from the surface atoms of one material and transfer them to the other. One body ends with an excess of electrons (negative), the other with an equal deficit (positive). This is \cle{charging by friction}.

\begin{exemplebox}[Classic friction pairs]
\begin{itemize}
\item Ebonite rod rubbed with fur: the \textbf{ebonite} becomes \textbf{negative} (it gains electrons), the fur positive.
\item Glass rod rubbed with silk: the \textbf{glass} becomes \textbf{positive} (it loses electrons), the silk negative.
\item A plastic comb pulled through dry hair becomes negative and can then attract small pieces of paper.
\end{itemize}
\end{exemplebox}

\begin{attention}[Only electrons move]
In all charging processes at this level, it is always \textbf{electrons} that are transferred, never protons. A body becomes positive by \emph{losing} electrons, not by gaining protons. Saying ``positive charges moved'' in an exam costs marks.
\end{attention}

\courssub{Charging by Contact (Conduction)}

If a charged body touches a conductor, the mobile charges redistribute over the \emph{combined} surface of the two objects, which momentarily form a single conductor. When they are separated, the charge has been \emph{shared}: this is \cle{charging by contact} (or conduction).

\begin{propriete}[Sharing of charge between identical spheres]
When two identical conducting spheres are brought into contact and then separated, the total charge is shared equally between them:
\[
q'_1 = q'_2 = \frac{q_1 + q_2}{2}.
\]
This follows from symmetry: identical spheres at equilibrium must carry identical charges, and total charge is conserved. For conductors of different sizes (or capacitances), the charge distributes in proportion to their capacitances (for spheres, proportional to their radii).
\end{propriete}

\begin{exoresolu}[Sharing charge between two spheres]
A conducting sphere A carries a charge $q_A = +8.0\ \mathrm{nC}$ and an identical sphere B carries $q_B = -2.0\ \mathrm{nC}$. The spheres are brought into contact and then separated. Find the final charge on each sphere.
\tcblower
\textbf{Solution.} Total charge is conserved:
\[
q_{\text{total}} = q_A + q_B = +8.0 - 2.0 = +6.0\ \mathrm{nC}.
\]
The spheres are identical, so at equilibrium the charge is shared equally:
\[
q'_A = q'_B = \frac{q_A + q_B}{2} = \frac{+6.0}{2} = +3.0\ \mathrm{nC}.
\]
Each sphere finally carries $+3.0\ \mathrm{nC}$. Physically: the $-2.0\ \mathrm{nC}$ of B first neutralises $2.0\ \mathrm{nC}$ of A's positive charge, and the remaining $+6.0\ \mathrm{nC}$ splits into two equal halves.
\end{exoresolu}

Note that contact charging gives the conductor a charge of the \emph{same sign} as the charging body (electrons flow from the negative body onto the conductor, or from the conductor onto a positive body).

\courssub{Charging by Induction}

Can we charge a conductor \emph{without touching it}, and even obtain a charge \emph{opposite} to that of the charging body? Yes --- this is \cle{charging by induction}, the most subtle and most frequently examined method.

The key idea: when a charged rod is brought \emph{near} (but not touching) a neutral conductor, the free electrons inside the conductor migrate. If the rod is negative, electrons are repelled to the far side; the near side becomes positive and the far side negative. The conductor is still neutral overall --- its charges have merely been \emph{separated}. If we now connect the far side to the earth, the repelled electrons escape to the ground. Remove the earth connection, then the rod: the conductor is left with a net \textbf{positive} charge.

\begin{methode}[Charging a conductor by induction in four steps]
To charge a neutral conducting sphere \textbf{negatively} using a \textbf{positively} charged rod:
\begin{enumerate}
\item \textbf{Approach:} bring the positively charged rod near (not touching) the sphere. Electrons in the sphere are attracted to the near side; the far side is left positive.
\item \textbf{Earth:} while holding the rod in place, touch the sphere (or connect it to the ground with a wire). Electrons flow \emph{up} from the earth, attracted by the rod, and neutralise the positive side.
\item \textbf{Remove the earth:} break the earth connection first, keeping the rod in place. The extra electrons are now trapped on the sphere.
\item \textbf{Remove the rod:} take the rod away. The excess electrons spread uniformly over the surface: the sphere is negatively charged.
\end{enumerate}
\end{methode}

\begin{popfigure}
\begin{tikzpicture}[scale=0.86]
% Step 1
\begin{scope}
\draw[thick, popblue, fill=popblueL] (0,0) circle (0.9);
\draw[thick, poporange, fill=popgold!40] (-2.3,-0.15) rectangle (-1.15,0.15);
\node[popdark] at (-1.7,0.45) {\scriptsize $+\,+\,+$};
\node at (-0.55,0.25) {\scriptsize $-\,-$};
\node at (0.55,0.25) {\scriptsize $+\,+$};
\node[align=center, text width=2.6cm] at (0,-1.7) {\small 1. Rod brought near};
\end{scope}
% Step 2
\begin{scope}[xshift=4.2cm]
\draw[thick, popblue, fill=popblueL] (0,0) circle (0.9);
\draw[thick, poporange, fill=popgold!40] (-2.3,-0.15) rectangle (-1.15,0.15);
\node[popdark] at (-1.7,0.45) {\scriptsize $+\,+\,+$};
\node at (-0.55,0.25) {\scriptsize $-\,-$};
\draw[thick, popdark] (0.9,0) -- (1.6,0) -- (1.6,-1.0);
\draw[popdark] (1.3,-1.0) -- (1.9,-1.0);
\draw[popdark] (1.4,-1.12) -- (1.8,-1.12);
\draw[popdark] (1.5,-1.24) -- (1.7,-1.24);
\node[align=center, text width=2.6cm] at (0,-1.9) {\small 2. Earth connected: electrons arrive};
\end{scope}
% Step 3
\begin{scope}[xshift=8.4cm]
\draw[thick, popblue, fill=popblueL] (0,0) circle (0.9);
\draw[thick, poporange, fill=popgold!40] (-2.3,-0.15) rectangle (-1.15,0.15);
\node[popdark] at (-1.7,0.45) {\scriptsize $+\,+\,+$};
\node at (-0.55,0.25) {\scriptsize $-\,-$};
\node at (-0.2,-0.35) {\scriptsize $-$};
\node[align=center, text width=2.6cm] at (0,-1.7) {\small 3. Earth removed first};
\end{scope}
% Step 4
\begin{scope}[xshift=12.6cm]
\draw[thick, popblue, fill=popblueL] (0,0) circle (0.9);
\foreach \a in {20,70,120,170,220,270,320}
  \node[popdark] at (\a:1.02) {\scriptsize $-$};
\node[align=center, text width=2.6cm] at (0,-1.7) {\small 4. Rod removed: sphere negative};
\end{scope}
\end{tikzpicture}
\legende{Charging a neutral sphere negatively by induction with a positive rod. The earth connection must be broken \emph{before} the rod is removed.}
\end{popfigure}

\begin{attention}[The induced charge is always OPPOSITE --- exam trap]
The charge obtained by induction is always \textbf{opposite in sign} to the charge of the inducing rod. A positive rod induces a \emph{negative} charge on the earthed conductor; a negative rod induces a \emph{positive} charge. Students frequently write ``same sign'' by confusion with charging by contact. Remember: contact gives the \emph{same} sign, induction gives the \emph{opposite} sign. A second trap: if you remove the rod \emph{before} breaking the earth connection, the electrons simply flow back to the earth and the sphere stays neutral --- the order of steps 3 and 4 is essential.
\end{attention}

\courssub{Earthing (Grounding)}

\begin{definition}[Earthing]
\cle{Earthing} (or grounding) a conductor means connecting it to the Earth by a conducting path. The Earth is so large that it behaves as a practically \cle{infinite reservoir of charge}: it can accept or supply electrons in enormous numbers without its own electrical state changing noticeably.
\end{definition}

\begin{methode}[Explaining earthing correctly]
To analyse any earthing situation, always reason in terms of \textbf{electron flow}:
\begin{enumerate}
\item Identify the sign of the excess charge on the conductor.
\item If the conductor is \textbf{negative}: the excess electrons are repelled and flow \emph{down to the earth}; the conductor becomes neutral.
\item If the conductor is \textbf{positive}: electrons are attracted \emph{up from the earth} and neutralise the deficit; the conductor becomes neutral.
\item Never say ``positive charges flow to the earth'' --- positive charges (protons in the nuclei) do not move in a solid.
\end{enumerate}
\end{methode}

Earthing also has a vital safety role: any unwanted accumulation of static charge on a metal object is drained away before it can build up enough to produce a spark.

\courssub{Applications and Hazards of Static Electricity}

\begin{itemize}
\item \textbf{Fuel tankers.} A tanker lorry carrying petrol along the Bamenda--Bafoussam road has its fuel sloshing and its tyres rubbing the road; friction charges the metal body. A spark near fuel vapour would be catastrophic. This is why tankers trail an \textbf{earthing chain} (or a conducting strip) that keeps the body connected to the ground, and why they are earthed before fuel is transferred at the depot.
\item \textbf{Photocopiers and laser printers.} These machines use static electricity deliberately: a drum is charged, light from the document discharges selected regions, and charged toner particles are attracted only to the charged image areas, then pressed onto paper.
\item \textbf{Dust on screens.} A television or computer screen acquires static charge in operation and then attracts dust particles from the air --- the same attraction a charged comb exerts on pieces of paper.
\item \textbf{Lightning.} As we shall see in the next section, the lightning conductor protects buildings precisely by combining earthing with the concentration of charge at sharp points.
\end{itemize}

\begin{savaistu}[Why you get a small shock after a car journey]
Passengers sliding across synthetic car seats become charged by friction. When you step out and touch the metal door, the charge flows suddenly through your finger to the earth, producing a small spark and the familiar ``zap''. It is harmless at this scale, but exactly the same spark near petrol vapour at a filling station can be dangerous --- one reason why filling stations forbid re-entering the vehicle during refuelling.
\end{savaistu}

\begin{exercice}[Test your understanding]
\begin{enumerate}
\item A negatively charged ebonite rod is brought near (without touching) a neutral electroscope whose cap is then earthed. State, with reasons, the sign of the charge left on the electroscope after the earth is removed and then the rod withdrawn.
\item Two identical conducting spheres carry $+12\ \mathrm{nC}$ and $-4\ \mathrm{nC}$. They are touched together and separated. Find the final charge on each.
\item Explain why a charged metal sphere loses its charge when touched with the bare hand, but a charged glass rod does not.
\end{enumerate}
\end{exercice}

\begin{corrige}[Exercise 1]
\begin{enumerate}
\item The negative rod repels electrons in the electroscope to the far end; earthing lets them escape to the ground. Breaking the earth traps a deficit of electrons, so the electroscope is left \textbf{positively} charged --- opposite to the inducing rod, as always with induction.
\item Final charge on each sphere: $q' = \dfrac{+12 + (-4)}{2} = +4\ \mathrm{nC}$.
\item The metal sphere, the hand, the body and the earth form one conductor: the charge spreads over this huge system (effectively flowing to earth), so the sphere is discharged. The glass rod is an insulator: its charge is bound where it was deposited and cannot flow through the rod into the hand.
\end{enumerate}
\end{corrige}

\begin{aretenir}[Key points of this section]
\begin{itemize}
\item Conductors contain free electrons (metals, graphite, human body, damp earth, salt solutions); insulators have bound electrons (glass, plastic, ebonite, dry air).
\item On a conductor at equilibrium, net charge resides entirely on the \textbf{outer surface}; on an insulator it stays where deposited.
\item Friction transfers electrons between two insulators; contact shares charge (same sign); induction charges a conductor without contact, giving the \textbf{opposite} sign.
\item Induction procedure: approach rod $\rightarrow$ earth $\rightarrow$ remove earth $\rightarrow$ remove rod. Order matters.
\item Earthing is always explained by \textbf{electron flow} to or from the Earth, an infinite reservoir of charge; it is the basis of electrical safety (tanker chains, lightning conductors).
\end{itemize}
\end{aretenir}

\coursec{Point Action and the Lightning Conductor}

In the previous section we saw that charge can be placed on a conductor by friction, contact or induction, and that on an isolated conductor the charge resides entirely on its outer surface. But \emph{where exactly} on the surface does the charge sit? If the conductor is a perfect sphere, symmetry tells us the charge spreads uniformly. Real conductors, however, are rarely perfect spheres: a key, a nail, a church spire, the tip of a radio antenna. In this section we discover that the shape of a conductor dramatically changes how charge is distributed on it, and that this simple fact — \cle{point action} — is what makes the lightning conductor one of the oldest and most successful applications of electrostatics.

\courssub{Surface Charge Density and Curvature}

\textbf{Observation.} Charge a pear-shaped conductor (or simply a metal sphere with a small pointed rod attached) using an electrostatic machine. Then touch different parts of its surface with a small proof plane (a tiny metal disc on an insulating handle) and transfer the collected charge to an electroscope. The result is striking: the proof plane picks up far more charge from the pointed end than from the rounded body. The charge is not spread uniformly — it \emph{crowds at the sharp point}.

\begin{definition}[Surface charge density]
The \cle{surface charge density} $\sigma$ at a point of a charged conductor is the charge per unit area at that point:
\[
\sigma = \frac{Q}{A}
\]
with $\sigma$ in $\mathrm{C\,m^{-2}}$. On an isolated conductor, $\sigma$ is \textbf{not uniform} unless the conductor is spherical.
\end{definition}

\begin{propriete}[Charge density and curvature — proof not examinable]
On a conductor in electrostatic equilibrium, the surface charge density is \cle{greatest where the radius of curvature is smallest}, i.e. at the sharpest points and edges. Equivalently: the more sharply curved the surface, the more concentrated the charge, and the stronger the electric field just outside the surface.
\end{propriete}

\textbf{Why is this so? (instructive justification).} Consider two conducting spheres of radii $R_1$ and $R_2$ ($R_1 > R_2$), connected by a long thin wire so that they form a single conductor and must be at the same potential $V$. Since $V = \dfrac{Q}{4\pi\varepsilon_0 R}$ for an isolated sphere, equality of potentials gives
\[
\frac{Q_1}{4\pi\varepsilon_0 R_1} = \frac{Q_2}{4\pi\varepsilon_0 R_2}
\quad\Longrightarrow\quad
\frac{Q_1}{Q_2} = \frac{R_1}{R_2}.
\]
The charge densities are $\sigma_1 = \dfrac{Q_1}{4\pi R_1^2}$ and $\sigma_2 = \dfrac{Q_2}{4\pi R_2^2}$, so
\[
\frac{\sigma_2}{\sigma_1} = \frac{Q_2}{Q_1}\cdot\frac{R_1^2}{R_2^2} = \frac{R_2}{R_1}\cdot\frac{R_1^2}{R_2^2} = \frac{R_1}{R_2} > 1.
\]
The \textbf{smaller} sphere (smaller radius of curvature) carries the \textbf{larger} surface charge density. A sharp point is the limiting case of a very small radius of curvature, so $\sigma$ becomes very large there. Since the field just outside a conductor is $E = \sigma/\varepsilon_0$, the field near a sharp point is extremely intense.

\begin{popfigure}
\begin{center}
\begin{tikzpicture}[scale=1.0]
% irregular conductor: rounded left, pointed right
\draw[thick, popblue, fill=popblueL] plot[smooth cycle, tension=0.9] coordinates {(-3.2,0) (-2.6,1.3) (-1.2,1.6) (0.2,1.2) (1.4,0.9) (3.4,0) (1.4,-0.9) (0.2,-1.2) (-1.2,-1.6) (-2.6,-1.3)};
% sparse + on rounded part
\node[popink] at (-2.6,0.9) {\small $+$};
\node[popink] at (-2.2,-1.0) {\small $+$};
\node[popink] at (-1.4,1.25) {\small $+$};
\node[popink] at (-1.3,-1.3) {\small $+$};
\node[popink] at (-0.4,1.15) {\small $+$};
\node[popink] at (-0.3,-1.15) {\small $+$};
\node[popink] at (0.6,0.95) {\small $+$};
\node[popink] at (0.6,-0.95) {\small $+$};
% dense + at the tip
\node[popink] at (1.5,0.72) {\small $+$};
\node[popink] at (1.5,-0.72) {\small $+$};
\node[popink] at (1.9,0.55) {\small $+$};
\node[popink] at (1.9,-0.55) {\small $+$};
\node[popink] at (2.25,0.4) {\small $+$};
\node[popink] at (2.25,-0.4) {\small $+$};
\node[popink] at (2.55,0.25) {\small $+$};
\node[popink] at (2.55,-0.25) {\small $+$};
\node[popink] at (2.8,0.12) {\small $+$};
\node[popink] at (2.8,-0.12) {\small $+$};
\node[popink] at (3.0,0) {\small $+$};
% field arrows at tip
\draw[-{Stealth}, thick, poporange] (3.4,0) -- (4.6,0);
\draw[-{Stealth}, thick, poporange] (3.1,0.45) -- (4.1,0.85);
\draw[-{Stealth}, thick, poporange] (3.1,-0.45) -- (4.1,-0.85);
\node[poporange, align=center] at (4.6,1.25) {\small strong field $\vec{E}$};
\node[popink, align=center] at (-1.6,-2.2) {\small low $\sigma$ (large radius of curvature)};
\node[popink, align=center] at (2.6,-1.5) {\small high $\sigma$ (sharp point)};
\draw[-{Stealth}, popmuted] (2.6,-1.3) -- (2.9,-0.25);
\end{tikzpicture}
\end{center}
\legende{Charge distribution on an irregular conductor: positive charge crowds at the sharp tip where the radius of curvature is smallest, producing an intense local electric field.}
\end{popfigure}

\courssub{Point Action}

\begin{definition}[Point action]
\cle{Point action} is the rapid loss (or gain) of electric charge at the sharp points of a conductor, caused by the very high local surface charge density and the resulting intense electric field, which ionises the surrounding air.
\end{definition}

\textbf{Mechanism.} Air always contains a few ions and free electrons (produced by cosmic rays and natural radioactivity). Near a strongly charged sharp point, the field $E = \sigma/\varepsilon_0$ is so intense that these ions are violently accelerated:

\begin{enumerate}
\item Ions of sign \emph{opposite} to that of the conductor are attracted to the point and neutralise part of its charge — the conductor \textbf{leaks charge} into the air.
\item Ions of the \emph{same} sign are repelled and stream away from the point, carrying charge with them. This flow of charged air molecules is the \cle{electric wind}.
\item The accelerated ions collide with neutral air molecules and knock electrons off them (\emph{ionisation by collision}), multiplying the effect. In the dark, a faint glow (corona discharge) may be seen at the point.
\end{enumerate}

\begin{experience}[Demonstrating point action and the electric wind]
\textbf{(a) Pointed versus rounded conductor.} Charge two identical metal conductors — one ending in a sharp needle, the other perfectly rounded — to the same potential with an electrostatic machine, each connected to an electroscope. The pointed conductor's electroscope collapses much faster: the point leaks charge into the air far more quickly than the rounded surface.

\textbf{(b) The electric wind.} Bring a charged sharp needle close to a candle flame. The flame is blown away from the point, as if a wind were streaming from it. The repelled ions drag neutral air molecules along, producing a genuine draught. A light paper wheel mounted on a pivot with pointed arms spins when charged — the electric wind acts like a tiny reaction engine.
\end{experience}

\begin{attention}[Points work both ways]
Point action does not only \emph{discharge} a conductor. A neutral pointed conductor placed in the field of a nearby charged body readily \emph{collects} charge of opposite sign at its tip by induction and ion capture. Sharp points are therefore excellent both at leaking charge away and at gathering it — both aspects are exploited in the lightning conductor.
\end{attention}

\courssub{Lightning: a Giant Spark}

\begin{savaistu}[Why thunderclouds are charged (enrichment — not examinable)]
Inside a thundercloud, strong updrafts and downdrafts hurl water droplets, ice crystals and hailstones against one another. Collisions and the freezing of supercooled droplets transfer charge: typically the \textbf{top} of the cloud becomes positively charged and the \textbf{base} negatively charged, while the ground beneath is charged positively by induction. The detailed mechanism is still an active research topic — what matters for us is the result: an enormous charge separation, and hence an enormous potential difference, of the order of $10^{8}$ to $10^{9}\ \mathrm{V}$ between cloud and ground.
\end{savaistu}

When the electric field between the base of the cloud and the ground exceeds the insulating strength of air (about $3\times 10^{6}\ \mathrm{V\,m^{-1}}$ in dry air, less in humid air), the air ionises catastrophically and a conducting channel opens: a \textbf{lightning flash} — essentially a giant spark, like the spark of an electrostatic machine but kilometres long. A current of tens of thousands of amperes flows for a fraction of a second, heating the channel to tens of thousands of degrees; the explosive expansion of the heated air is the thunderclap we hear across the hills of the North West or the plains of the Far North during the rainy season.

\courssub{The Lightning Conductor (Franklin Rod)}

The \cle{lightning conductor}, invented by Benjamin Franklin in the 18th century, is a direct application of point action. It consists of three essential parts:

\begin{enumerate}
\item A \textbf{pointed metal rod} (the \emph{air terminal}), fixed at the highest point of the building, ending in one or more sharp tips.
\item A \textbf{thick copper down-conductor} (large cross-section, hence very low resistance), running down the outside of the building with as few bends as possible.
\item A deep \textbf{earth connection}: a copper plate or rod buried in permanently damp ground, ensuring excellent electrical contact with the earth.
\end{enumerate}

\begin{popfigure}
\begin{center}
\begin{tikzpicture}[scale=0.9]
% ground
\fill[popgreenL] (-4.5,0) rectangle (4.5,-0.25);
\draw[thick, popgreen] (-4.5,0) -- (4.5,0);
\node[popgreen!60!black] at (-3.4,-0.55) {\small ground};
% cloud
\draw[thick, popmuted, fill=gray!20] plot[smooth cycle, tension=0.8] coordinates {(-4.2,5.6) (-3.2,6.4) (-1.8,6.1) (-0.6,6.6) (0.8,6.2) (2.2,6.5) (3.6,6.0) (4.2,5.4) (3.0,5.0) (1.2,5.2) (-0.8,4.9) (-2.6,5.1)};
\node[popink] at (-2.6,5.7) {\small $-\ -\ -$};
\node[popink] at (0.2,5.6) {\small $-\ -\ -$};
\node[popink] at (2.6,5.7) {\small $-\ -\ -$};
\node[popink, align=center] at (0,6.9) {\small charged thundercloud};
% building
\draw[thick, popdark, fill=popblueL] (-1.6,0) rectangle (1.6,3.2);
\draw[thick, popdark] (-1.9,3.2) -- (0,4.0) -- (1.9,3.2);
\draw[thick, popdark, fill=white] (-0.5,0) rectangle (0.5,1.1);
\draw[thick, popdark, fill=white] (-1.1,1.7) rectangle (-0.5,2.5);
\draw[thick, popdark, fill=white] (0.5,1.7) rectangle (1.1,2.5);
% lightning rod
\draw[very thick, poporange] (0,4.0) -- (0,5.1);
\draw[very thick, poporange] (0,5.1) -- (0,5.5);
\fill[poporange] (0,5.5) circle (0.03);
\node[poporange, align=center, anchor=west] at (0.25,5.35) {\small pointed rod};
% corona / point action arrows near tip
\draw[-{Stealth}, poppurple] (-0.1,5.45) -- (-0.9,5.9);
\draw[-{Stealth}, poppurple] (0.1,5.45) -- (0.9,5.9);
\node[poppurple, align=center] at (-2.4,6.35) {\small point action: gradual discharge};
% down conductor
\draw[very thick, poporange] (0,4.0) -- (0,3.2);
\draw[very thick, poporange] (1.6,3.2) -- (1.6,0) -- (1.6,-1.3);
\draw[very thick, poporange] (0,3.2) -- (1.6,3.2);
\node[poporange, align=center, anchor=west] at (1.8,1.8) {\small thick copper\\[-2pt]\small down-conductor};
% earth plate
\fill[popgold] (1.1,-1.5) rectangle (2.1,-1.75);
\node[popgold!70!black, align=center, anchor=west] at (2.3,-1.6) {\small earth plate buried\\[-2pt]\small in damp soil};
\draw[very thick, poporange] (1.6,-1.3) -- (1.6,-1.5);
% induced + on ground
\node[popink] at (-0.8,-0.55) {\small $+\ +\ +\ +\ +$};
\end{tikzpicture}
\end{center}
\legende{A building protected by a lightning conductor: pointed air terminal, low-resistance copper down-conductor, and earth plate buried deep in damp ground.}
\end{popfigure}

\begin{methode}[Explaining how a lightning conductor protects (two stages)]
\begin{enumerate}
\item \textbf{Preventive discharge by point action.} As a charged cloud passes overhead, it induces a strong opposite charge on the building and especially at the sharp tip of the rod (charge crowds where curvature is smallest). The intense field at the point ionises the air; a stream of ions flows gently upward, \emph{gradually neutralising} the cloud's charge and reducing the potential difference. Often this silent discharge prevents the strike altogether.
\item \textbf{Safe channelling of the strike.} If the charge accumulation is too rapid and lightning does strike, it strikes the rod — the highest, sharpest, best-earthed point — and the enormous current is led through the thick copper conductor \emph{straight to earth}, instead of passing through the roof, walls, wiring or occupants of the building.
\end{enumerate}
\end{methode}

\begin{attention}[The rod does not dangerously ``attract'' lightning]
A common misconception is that a lightning conductor ``attracts'' lightning and is therefore risky. This is wrong. The rod does not summon lightning from afar; a strike in its vicinity would have occurred anyway. What the rod does is (i) quietly discharge the cloud by point action, often preventing the strike, and (ii) if a strike occurs, offer it the \textbf{safest possible path to earth} — a controlled, low-resistance route that bypasses the structure. A building with a properly earthed conductor is far safer than one without.
\end{attention}

\begin{exoresolu}[Charge density at a tip]
Two conducting spheres, of radii $R_1 = 10.0\ \mathrm{cm}$ and $R_2 = 1.0\ \mathrm{mm}$, are connected by a thin wire and charged. The total charge is such that the large sphere carries $Q_1 = 2.0\times 10^{-8}\ \mathrm{C}$.
\begin{enumerate}
\item Find the charge $Q_2$ on the small sphere.
\item Compare the surface charge densities $\sigma_1$ and $\sigma_2$.
\item Comment on the relevance to a lightning rod.
\end{enumerate}
\tcblower
\textbf{1.} Connected conductors are at the same potential, so $\dfrac{Q_1}{R_1} = \dfrac{Q_2}{R_2}$:
\[
Q_2 = Q_1\,\frac{R_2}{R_1} = 2.0\times 10^{-8}\times\frac{1.0\times 10^{-3}}{1.0\times 10^{-1}} = 2.0\times 10^{-10}\ \mathrm{C}.
\]
\textbf{2.} Surface charge densities:
\[
\sigma_1 = \frac{Q_1}{4\pi R_1^2} = \frac{2.0\times 10^{-8}}{4\pi (0.100)^2} \approx 1.6\times 10^{-7}\ \mathrm{C\,m^{-2}},
\]
\[
\sigma_2 = \frac{Q_2}{4\pi R_2^2} = \frac{2.0\times 10^{-10}}{4\pi (1.0\times 10^{-3})^2} \approx 1.6\times 10^{-5}\ \mathrm{C\,m^{-2}}.
\]
Hence $\dfrac{\sigma_2}{\sigma_1} = \dfrac{R_1}{R_2} = 100$: the density on the small sphere is \textbf{100 times larger}, although its charge is 100 times smaller.

\textbf{3.} A lightning-rod tip behaves like the small sphere: even a modest total charge produces an enormous local charge density and field, enough to ionise air and discharge the cloud gradually — the essence of point action.
\end{exoresolu}

\courssub{Practical Rules and Safety in Cameroon}

\begin{itemize}
\item Tall buildings, churches, mosques, water towers and telecommunication masts — increasingly common in Yaoundé, Douala and Bamenda — should be fitted with lightning conductors whose earth plates reach permanently moist soil.
\item During a storm, \textbf{never shelter under an isolated tree}: the wet tree is a natural pointed conductor, and a strike can pass through anyone beneath it. Open fields and isolated high points are equally dangerous.
\item Indoors, unplug electrical appliances and avoid touching plumbing or landline telephones during a storm, since strikes on power or phone lines can send surges into the house.
\item If caught outdoors with no shelter, crouch low on the balls of your feet (minimising contact area with the ground) rather than lying flat, and stay away from metal fences and poles.
\end{itemize}

\begin{savaistu}[Exam tip]
GCE questions on this topic very often ask candidates to \textbf{justify the pointed shape} of the rod (answer: charge density and field are greatest where the radius of curvature is smallest, so point action discharges the cloud gradually) and \textbf{why the rod must be earthed} (answer: to give the strike current a safe, low-resistance path into the ground, protecting the building). Structure your answer with the two stages of the method box above and you will collect full marks.
\end{savaistu}

\begin{exercice}[Point action and protection]
\begin{enumerate}
\item Explain, in terms of surface charge density, why a charged conductor with a sharp point loses its charge faster than a spherical conductor charged to the same potential.
\item Describe the electric wind and state its cause.
\item A student claims: ``Since the lightning conductor attracts lightning, houses without one are safer.'' Refute this claim in three or four sentences.
\item Why should the down-conductor be thick and run with as few sharp bends as possible?
\end{enumerate}
\end{exercice}

\begin{corrige}[Exercise 1]
\textbf{1.} The surface charge density $\sigma$ is greatest where the radius of curvature is smallest. At a sharp point $\sigma$ — and hence the local field $E=\sigma/\varepsilon_0$ — is very large, so the surrounding air is ionised and charge leaks away rapidly by point action. On the sphere, $\sigma$ is moderate and uniform, so leakage is slow.

\textbf{2.} The electric wind is the stream of air moving away from a charged sharp point. Ions of the same sign as the conductor are repelled from the point and drag neutral air molecules with them, producing a draught capable of deflecting a candle flame.

\textbf{3.} The conductor does not attract lightning from a distance; a strike near it would have occurred regardless. By point action it continuously and quietly discharges the cloud, often preventing the strike. If a strike occurs, it is directed through the rod and thick copper conductor safely to earth, instead of through the building. A house without a conductor has no such protection and is therefore \emph{less} safe.

\textbf{4.} A strike carries an enormous current (tens of thousands of amperes). A thick conductor has very low resistance, so the heating ($P = RI^2$) and the voltage drop along it remain small, avoiding fire and side-flashes. Sharp bends have local high curvature where the current could jump off the conductor (side-flash to the building); smooth routing keeps the current confined to the safe path to earth.
\end{corrige}

\begin{aretenir}[Key points]
\begin{itemize}
\item Surface charge density on a conductor is greatest where the radius of curvature is smallest — charge crowds at sharp points.
\item The intense field at a point ionises the air: the conductor leaks charge (\cle{point action}) and a repelled stream of ions forms the \cle{electric wind}.
\item Lightning is a giant spark between a charged cloud and the ground.
\item The lightning conductor protects in two stages: gradual preventive discharge by point action, then safe channelling of any strike through a thick copper conductor to a deep earth plate.
\item The rod does not dangerously attract lightning — it provides the safest path to earth.
\end{itemize}
\end{aretenir}

\coursec{Coulomb's Law and the Electric Field}

In the previous section we saw how charge concentrates at sharp points, producing the intense local effects behind point action and the lightning conductor. Those observations tell us \emph{where} charges act, but not \emph{how strongly}. To make electrostatics a quantitative science we must answer two questions: what is the exact size of the force between two charges, and how can we describe the influence a charge exerts on the space around it? The first question is answered by \cle{Coulomb's law}; the second leads us to one of the most powerful ideas in all of physics, the \cle{electric field}.

\courssub{Coulomb's Law}

\textbf{From observation to a law.}\par
Rub a balloon and hold it near small pieces of paper: the attraction is obvious. Now move the balloon further away and the pieces fall back. The force clearly \emph{decreases with distance}. Bring two charged balloons close together and they push each other apart more violently the closer they get. In the 1780s, Charles-Augustin de Coulomb measured this force precisely with a torsion balance — a light rod carrying a charged sphere, suspended by a fine wire whose twist measures tiny forces. His results can be summarised in a single mathematical statement.

\begin{propriete}[Coulomb's law]
The electrostatic force between two point charges $q_1$ and $q_2$ separated by a distance $r$ in vacuum has magnitude
\[
F=\frac{1}{4\pi\varepsilon_0}\,\frac{|q_1 q_2|}{r^2}=k\,\frac{|q_1 q_2|}{r^2},
\qquad k=\frac{1}{4\pi\varepsilon_0}\approx \SI{9.0e9}{N.m^2.C^{-2}},
\]
where $\varepsilon_0=\SI{8.85e-12}{C^2.N^{-1}.m^{-2}}$ is the \cle{permittivity of free space}. The force is directed \textbf{along the line joining the two charges}: it is \textbf{repulsive} if the charges have the same sign and \textbf{attractive} if they have opposite signs. This law is examinable and must be quoted with its conditions of validity (point charges, or spherical charge distributions measured from their centres).
\end{propriete}

\begin{attention}[Vector nature of the force]
Coulomb's law as written above gives only the \emph{magnitude}. The direction must always be added by reasoning on the signs: like charges repel, unlike charges attract, and the force on each charge lies on the line joining them. The two forces form an action--reaction pair (Newton's third law): equal in magnitude, opposite in direction, even if the charges themselves are very different in size.
\end{attention}

\textbf{The inverse-square character.}\par
The factor $1/r^2$ is the heart of the law: double the separation and the force falls to a quarter; triple it and the force falls to a ninth. How do we verify this experimentally? If $F\propto 1/r^2$, then a graph of $F$ against $1/r^2$ must be a \emph{straight line through the origin}, whose slope is $k|q_1q_2|$. This is the standard way of testing a power law: plot the quantity against the suspected variable and look for linearity.

\begin{popfigure}
\begin{center}
\begin{tikzpicture}
\begin{axis}[
  width=10cm, height=7cm,
  xlabel={$1/r^2\;(\mathrm{m^{-2}})$},
  ylabel={$F\;(\mathrm{mN})$},
  xmin=0, xmax=110, ymin=0, ymax=10,
  xtick={0,25,50,75,100},
  ytick={0,2,4,6,8,10},
  grid=both, grid style={popline, very thin},
  axis lines=left, thick,
]
\addplot[domain=0:100, samples=2, very thick, popblue] {0.09*x};
\addplot[only marks, mark=*, mark size=1.8pt, poporange] coordinates {(11.1,1.0) (25,2.3) (44.4,4.0) (69.4,6.2) (100,9.0)};
\node[popdark, align=center, font=\small] at (axis cs:72,3.2) {slope $=k|q_1q_2|$};
\end{axis}
\end{tikzpicture}
\end{center}
\legende{Force between two fixed charges measured at several separations, plotted against $1/r^2$. The points fall on a straight line through the origin, confirming $F\propto 1/r^2$.}
\end{popfigure}

\textbf{Dependence on the medium.}\par
Coulomb's experiments were performed in air (for which the result is essentially the vacuum value). When the charges are immersed in an insulating material — oil, glass, water — the force is \emph{reduced}. The material becomes polarised: its molecules develop small induced dipoles that partially screen the charges from each other.

\begin{propriete}[Force in a dielectric medium]
In an insulating medium of \cle{relative permittivity} $\varepsilon_r$ (a pure number, $\varepsilon_r>1$), the permittivity is $\varepsilon=\varepsilon_0\varepsilon_r$ and Coulomb's law becomes
\[
F=\frac{1}{4\pi\varepsilon_0\varepsilon_r}\,\frac{|q_1q_2|}{r^2}=\frac{F_{\text{vacuum}}}{\varepsilon_r}.
\]
The force is divided by $\varepsilon_r$. For example, in water ($\varepsilon_r\approx 80$) the electrostatic force between two ions is about 80 times weaker than in vacuum — one reason why ionic salts dissolve so readily in water.
\end{propriete}

\begin{savaistu}[Electrostatics versus gravitation]
Newton's law of gravitation, $F=Gm_1m_2/r^2$, has exactly the same inverse-square form as Coulomb's law. But the resemblance stops there. Between a proton and an electron ($r=\SI{5.3e-11}{m}$ in a hydrogen atom), the electric attraction is about $2\times10^{39}$ times stronger than their gravitational attraction! Moreover, gravity is always attractive, whereas the electric force can attract \emph{or} repel. It is precisely because matter contains almost exactly equal numbers of positive and negative charges that the vastly stronger electric force usually hides, leaving the feeble but unscreened gravitational force to rule the motion of planets.
\end{savaistu}

\courssub{Superposition of Electrostatic Forces}

What happens when a charge is acted on by \emph{several} other charges? Experiment shows that electrostatic forces combine exactly like all other forces in mechanics.

\begin{propriete}[Principle of superposition]
The resultant force on a charge $q$ due to several other charges is the \textbf{vector sum} of the forces that each charge would exert on $q$ separately, as if the others were absent:
\[
\vec{F}_{\text{res}}=\vec{F}_1+\vec{F}_2+\vec{F}_3+\cdots
\]
\end{propriete}

\begin{methode}[Solving two-charge (or many-charge) problems]
\begin{enumerate}
\item \textbf{Sketch} the situation: mark all charges with their signs and the distances between them.
\item \textbf{Compute magnitudes} of each force with Coulomb's law, $F_i=k|q_iq|/r_i^2$ (take the medium into account if $\varepsilon_r\neq 1$).
\item \textbf{Assign directions} from the signs: attraction towards, repulsion away from each source charge.
\item \textbf{Add vectorially}: if the forces are collinear, add with signs; otherwise resolve into components along perpendicular axes and recombine with Pythagoras' theorem and trigonometry.
\end{enumerate}
\end{methode}

\begin{popfigure}
\begin{center}
\begin{tikzpicture}[scale=1.1, >=stealth]
\fill[popblue] (-3,0) circle (0.14);
\node[white, font=\small\bfseries] at (-3,0) {+};
\node[popdark, below=4pt] at (-3,-0.15) {$q_1=+4\ \mu\mathrm{C}$};
\fill[poporange] (3,0) circle (0.14);
\node[white, font=\small\bfseries] at (3,0) {+};
\node[popdark, below=4pt] at (3,-0.15) {$q_2=+1\ \mu\mathrm{C}$};
\fill[popgreen!70!black] (0.6,0) circle (0.12);
\node[white, font=\small\bfseries] at (0.6,0) {$-$};
\node[popdark, above=4pt] at (0.6,0.15) {$q$ (test)};
\draw[<->, popmuted] (-3,-0.9) -- (0.6,-0.9) node[midway, below, popdark, font=\small]{$r_1$};
\draw[<->, popmuted] (0.6,-0.9) -- (3,-0.9) node[midway, below, popdark, font=\small]{$r_2$};
\draw[->, very thick, popblue] (0.6,0.35) -- (-1.6,0.35) node[midway, above, font=\small]{$\vec{F}_1$ (attraction)};
\draw[->, very thick, poporange] (0.6,-0.35) -- (2.0,-0.35) node[midway, below, font=\small]{$\vec{F}_2$ (attraction)};
\draw[->, ultra thick, popdark] (0.6,0) -- (-0.9,0) node[above left=1pt, font=\small\bfseries]{$\vec{F}_{\text{res}}$};
\end{tikzpicture}
\end{center}
\legende{Superposition on a straight line: a negative test charge $q$ between two positive charges is attracted by both; the resultant is the vector sum $\vec{F}_1+\vec{F}_2$.}
\end{popfigure}

\begin{exoresolu}[Force on a charge between two fixed charges]
A charge $q=-2.0\ \mu\mathrm{C}$ is placed on a straight line between $q_1=+4.0\ \mu\mathrm{C}$ and $q_2=+1.0\ \mu\mathrm{C}$. The distance $q_1q$ is $r_1=\SI{0.30}{m}$ and the distance $qq_2$ is $r_2=\SI{0.20}{m}$ (air, $k=\SI{9.0e9}{N.m^2.C^{-2}}$). Find the resultant force on $q$.
\tcblower
\textbf{Step 1 — Sketch and directions.} $q$ is negative, so it is \emph{attracted} by both positive charges: $\vec{F}_1$ points towards $q_1$ (to the left), $\vec{F}_2$ towards $q_2$ (to the right). The forces are collinear and opposite.

\textbf{Step 2 — Magnitudes.}
\[
F_1=k\frac{|q_1q|}{r_1^2}=9.0\times10^{9}\times\frac{(4.0\times10^{-6})(2.0\times10^{-6})}{(0.30)^2}
=\frac{7.2\times10^{-2}}{0.090}=\SI{0.80}{N}.
\]
\[
F_2=k\frac{|q_2q|}{r_2^2}=9.0\times10^{9}\times\frac{(1.0\times10^{-6})(2.0\times10^{-6})}{(0.20)^2}
=\frac{1.8\times10^{-2}}{0.040}=\SI{0.45}{N}.
\]

\textbf{Step 3 — Vector sum.} Taking the direction towards $q_1$ as positive:
\[
F_{\text{res}}=F_1-F_2=0.80-0.45=\SI{0.35}{N},\quad\text{towards } q_1.
\]
Note how the inverse-square law makes $F_1$ dominant even though $q_1$ is only 4 times bigger than $q_2$: the smaller distance $r_2$ partly compensates, but not enough.
\end{exoresolu}

\begin{attention}[Never add forces as scalars]
If the charges are \emph{not collinear}, the forces point in different directions and you may \textbf{never} simply add their magnitudes. Resolve each force into components, add the components, and rebuild the resultant with $F=\sqrt{F_x^2+F_y^2}$ and $\tan\theta=F_y/F_x$. This is the single most common error in GCE electrostatics questions.
\end{attention}

\courssub{The Concept of Electric Field}

\textbf{Why invent a field?}\par
Coulomb's law describes a force acting at a distance, with nothing tangible between the charges. Physicists find it far more fruitful to think in two steps: a charge $Q$ \emph{modifies the space around it}, creating what we call an \cle{electric field}; any other charge placed in that space then experiences a force because of the field \emph{at its own position}. The field exists whether or not a second charge is present — it is a property of the space surrounding $Q$. To \emph{detect} and \emph{measure} it, we imagine placing a tiny positive \cle{test charge} $q$ at the point of interest and measuring the force on it. The test charge must be small enough not to disturb the source charges.

\begin{definition}[Electric field strength]
The \cle{electric field strength} $\vec{E}$ at a point is the force per unit positive test charge placed at that point:
\[
\vec{E}=\frac{\vec{F}}{q},\qquad\text{unit: }\mathrm{N\,C^{-1}}\ (\text{equivalently } \mathrm{V\,m^{-1}}).
\]
$\vec{E}$ is a vector: its direction is that of the force on a \emph{positive} test charge. A negative charge experiences a force opposite to $\vec{E}$.
\end{definition}

\textbf{Field of a point charge.}\par
Place a test charge $q$ at distance $r$ from a point source charge $Q$. By Coulomb's law, $F=k|Qq|/r^2$, so dividing by $q$:
\[
\boxed{\,E=\frac{F}{q}=\frac{1}{4\pi\varepsilon_0}\,\frac{|Q|}{r^2}\,}
\]
The field points \emph{radially away} from a positive source and \emph{radially towards} a negative source. Notice that $q$ has disappeared: the field depends only on the source $Q$ and the position. Like the force, the field obeys the \textbf{superposition principle}: the resultant field of several charges is the vector sum of their individual fields.

\begin{attention}[Source charge versus test charge]
In $\vec{E}=\vec{F}/q$, the symbol $q$ is the small \textbf{test charge} that feels the field. In $E=kQ/r^2$, the symbol $Q$ is the \textbf{source charge} that creates the field. Confusing the two — for instance dividing by the source charge — destroys the meaning of the formula and loses marks. Always ask yourself: \emph{who creates the field, and who feels it?}
\end{attention}

\begin{exemplebox}[Field at a point]
What is the electric field strength $\SI{0.50}{m}$ from a point charge $Q=+5.0\ \mu\mathrm{C}$ in air?
\[
E=k\frac{|Q|}{r^2}=9.0\times10^{9}\times\frac{5.0\times10^{-6}}{(0.50)^2}
=\SI{1.8e5}{N.C^{-1}},
\]
directed radially \emph{away} from $Q$ (since $Q>0$). A charge of $-3.0\ \mu\mathrm{C}$ placed there would feel $F=|q|E=\SI{0.54}{N}$ directed \emph{towards} $Q$.
\end{exemplebox}

\courssub{Electric Field Lines}

A vector at every point of space is hard to visualise, so we draw \cle{field lines}: curves whose tangent at each point gives the direction of $\vec{E}$, and whose spacing indicates the magnitude.

\begin{aretenir}[Properties of field lines]
\begin{itemize}
\item Field lines \textbf{start on positive} charges and \textbf{end on negative} charges (or go to infinity).
\item The tangent to a field line gives the \textbf{direction of $\vec{E}$} at that point.
\item Field lines \textbf{never cross}: the field has a single, well-defined direction at each point.
\item Where lines are \textbf{close together} the field is \textbf{strong}; where they spread apart it is weak. For a point charge the spreading in three dimensions is exactly why $E\propto 1/r^2$.
\item Lines meet a conductor's surface at right angles.
\end{itemize}
\end{aretenir}

\begin{popfigure}
\begin{center}
\begin{tikzpicture}[scale=0.9, >=stealth]
% Isolated positive charge
\begin{scope}[shift={(-3.6,0)}]
\foreach \a in {0,45,90,135,180,225,270,315}{
  \draw[->, popblue, thick] (\a:0.35) -- (\a:1.9);
}
\fill[popblue] (0,0) circle (0.22);
\node[white, font=\bfseries] at (0,0) {+};
\node[popdark, font=\small] at (0,-2.3) {isolated $+$ charge};
\end{scope}
% Dipole
\begin{scope}[shift={(3.2,0)}]
\foreach \t in {0.12,0.25,0.5,0.75,0.88}{
  \draw[poporange, thick, ->]
    plot[smooth] coordinates {
      ({-1.5+0.35*cos(180-\t*180)},{0.35*sin(\t*180)})
      ({-1.5+1.5*(1-cos(\t*180))*0.5+0.0},{1.1*sin(\t*180)})
      ({1.5-0.35*cos(\t*180)},{0.35*sin(\t*180)})
    };
  \draw[poporange, thick]
    plot[smooth] coordinates {
      ({-1.5+0.35*cos(180-\t*180)},{-0.35*sin(\t*180)})
      ({-1.5+1.5*(1-cos(\t*180))*0.5},{-1.1*sin(\t*180)})
      ({1.5-0.35*cos(\t*180)},{-0.35*sin(\t*180)})
    };
}
\draw[poporange, thick, ->] (-1.15,0) -- (0.4,0);
\draw[poporange, thick] (0.4,0) -- (1.15,0);
\fill[popblue] (-1.5,0) circle (0.22);
\node[white, font=\bfseries] at (-1.5,0) {+};
\fill[poporange] (1.5,0) circle (0.22);
\node[white, font=\bfseries] at (1.5,0) {$-$};
\node[popdark, font=\small] at (0,-2.3) {dipole: $+$ and $-$ pair};
\end{scope}
\end{tikzpicture}
\end{center}
\legende{Field-line maps. Left: an isolated positive point charge — radial lines pointing outwards. Right: an electric dipole — lines leave the positive charge and converge onto the negative one; the field is strongest where the lines are densest.}
\end{popfigure}

For the isolated positive charge the pattern is purely radial: the lines spread out like the spokes of a wheel, and their density falls off as $1/r^2$ — a geometrical picture of the inverse-square law. For the dipole, lines curve from the positive to the negative charge; along the axis between them the fields of the two charges reinforce, while far away the pattern begins to resemble that of a single point charge.

\courssub{The Uniform Field Between Parallel Plates}

One field pattern is so useful that it deserves a first mention here. When two large flat metal plates are placed parallel to each other and charged with equal and opposite charges, the field in the region between them (away from the edges) is \cle{uniform}: the field lines are straight, parallel and equally spaced. This means $\vec{E}$ has the \emph{same magnitude and direction everywhere} between the plates, so a charge placed anywhere in that region experiences the \emph{same force}. Uniform fields are used to steer electron beams — in the cathode-ray tubes of older oscilloscopes and television sets, for instance. We shall study this arrangement quantitatively in the chapters on capacitors and on the motion of charged particles; for now, simply keep the picture: \emph{parallel plates $\Rightarrow$ uniform field $\Rightarrow$ constant force}.

\begin{aretenir}[Key points of this section]
\begin{itemize}
\item Coulomb's law: $F=\dfrac{1}{4\pi\varepsilon_0}\dfrac{|q_1q_2|}{r^2}$, along the line joining the charges; attractive for unlike signs, repulsive for like signs.
\item The inverse-square law is verified by the straight line of $F$ versus $1/r^2$ through the origin.
\item In a dielectric, the force is divided by $\varepsilon_r$.
\item Several charges: add forces \textbf{vectorially} (superposition principle).
\item Electric field: $\vec{E}=\vec{F}/q$ (test charge $q$); field of a point source: $E=k|Q|/r^2$ (source charge $Q$); unit $\mathrm{N\,C^{-1}}$.
\item Field lines leave $+$, end on $-$, never cross, and their density shows the field strength.
\item Between oppositely charged parallel plates the field is uniform — a bridge to the study of capacitors.
\end{itemize}
\end{aretenir}

\begin{exercice}[Forces, fields and field lines]
\begin{enumerate}
\item Two point charges $q_1=+3.0\ \mu\mathrm{C}$ and $q_2=-6.0\ \mu\mathrm{C}$ are $\SI{0.40}{m}$ apart in air. Calculate the force each exerts on the other and state its nature. What does the force become if the gap is filled with oil of $\varepsilon_r=2.5$?
\item A charge $q=+1.0\ \mu\mathrm{C}$ is placed at the midpoint of the segment joining the two charges of question 1. Find the resultant force on it.
\item Calculate the electric field strength at the midpoint due to $q_1$ alone, then due to $q_2$ alone, and deduce the resultant field. Compare with your answer to question 2.
\item Explain, using field lines, why the field midway between two \emph{equal positive} charges is zero, while midway between a positive and a negative charge it is not.
\end{enumerate}
\end{exercice}

\begin{corrige}[Exercise 1]
\textbf{1.} $F=k|q_1q_2|/r^2=9.0\times10^{9}\times(3.0\times10^{-6})(6.0\times10^{-6})/(0.40)^2=\SI{1.0}{N}$, attractive (unlike signs). In oil: $F'=F/\varepsilon_r=1.0/2.5=\SI{0.40}{N}$.

\textbf{2.} At the midpoint $r=0.20\ \mathrm{m}$ from each charge. $q$ is positive: $q_1$ repels it (towards $q_2$) and $q_2$ attracts it (also towards $q_2$): the forces add.
$F_1=9.0\times10^{9}(3.0\times10^{-6})(1.0\times10^{-6})/(0.20)^2=\SI{0.675}{N}$;
$F_2=9.0\times10^{9}(6.0\times10^{-6})(1.0\times10^{-6})/(0.20)^2=\SI{1.35}{N}$.
$F_{\text{res}}=0.675+1.35\approx\SI{2.0}{N}$ towards $q_2$.

\textbf{3.} $E_1=F_1/q=\SI{6.75e5}{N.C^{-1}}$ and $E_2=\SI{1.35e6}{N.C^{-1}}$, both pointing towards $q_2$, so $E_{\text{res}}\approx\SI{2.0e6}{N.C^{-1}}$ towards $q_2$. Check: $F_{\text{res}}=qE_{\text{res}}=(1.0\times10^{-6})(2.0\times10^{6})=\SI{2.0}{N}$ — consistent.

\textbf{4.} With two equal positive charges, field lines point away from each charge; at the midpoint the two fields are equal and opposite, so they cancel ($E=0$, a neutral point). With a $+$ and a $-$ charge, lines go from $+$ to $-$; at the midpoint both fields point the same way (towards $-$), so they reinforce instead of cancelling.
\end{corrige}

\newpage

\coursec{Integration activity}

\coursec{Integration Activity: Lightning Protection of the Science Block}

\courssub{Aims of the Activity}

By the end of this integration activity, you should be able to mobilise, in a single coherent piece of work, all the key competences of this chapter:

\begin{itemize}
\item Use the \cle{two types of charge} and the \cle{basic law of electrostatics} to describe the electrification of a thundercloud and of a building.
\item Explain \cle{charging by induction} in a conductor and apply it to the rooftop of a building under a charged cloud.
\item Justify, using \cle{point action} and the concentration of charge on sharp points, the design of a \cle{Franklin rod} (lightning conductor).
\item Apply \cle{Coulomb's law} quantitatively, including the effect of the \cle{medium} through the relative permittivity $\varepsilon_r$.
\item Sketch and interpret \cle{electric field lines} between two charged objects, and explain how a lightning conductor modifies the field and protects a structure.
\item Communicate scientific results in the form of a short, well-argued \cle{technical note} with diagrams and calculations.
\end{itemize}

\courssub{Situation}

The Parent--Teacher Association (PTA) of a Government Bilingual High School in the North-West Region has just completed a new two-storey science block. The block contains a computer laboratory with equipment bought for $8\,000\,000$ FCFA. The school stands on a small hill, and the region experiences violent thunderstorms every rainy season. Last year, lightning struck a tree only $50\ \mathrm{m}$ from the old classroom block, and the head teacher is worried.

The PTA has decided to fit the science block with a \cle{lightning conductor} (Franklin rod) and has asked the Upper Sixth Physics class to act as \emph{technical consultants}. Your group must produce a \textbf{one-page technical note}, with diagrams and calculations, answering the committee's questions.

\textbf{Data provided by the PTA and the meteorological office}

\begin{center}
\begin{tabularx}{\linewidth}{|l|X|}
\hline
\rowcolor{popblueL} \textbf{Quantity} & \textbf{Value} \\
\hline
Height of the science block & $h = 8.0\ \mathrm{m}$ \\
\hline
Estimated charge on the base of a thundercloud (modelled as a point charge) & $Q_1 = -20\ \mathrm{C}$ \\
\hline
Estimated induced charge on the rooftop (modelled as a point charge) & $Q_2 = +2.0\times 10^{-3}\ \mathrm{C}$ \\
\hline
Mean distance from cloud base to rooftop & $d = 300\ \mathrm{m}$ \\
\hline
Relative permittivity of dry air & $\varepsilon_r = 1.00$ \\
\hline
Relative permittivity of humid rainy-season air & $\varepsilon_r = 1.01$ \\
\hline
Permittivity of free space & $\varepsilon_0 = 8.85\times 10^{-12}\ \mathrm{F\,m^{-1}}$ \\
\hline
Coulomb constant in vacuum & $k = \dfrac{1}{4\pi\varepsilon_0} = 9.0\times 10^{9}\ \mathrm{N\,m^2\,C^{-2}}$ \\
\hline
\end{tabularx}
\end{center}

The PTA committee asks four questions:
\begin{enumerate}
\item How does the negatively charged thundercloud make the rooftop of the block positively charged, even though the cloud never touches the building?
\item Why must the lightning conductor end in a \emph{sharp point}, and why must it be connected by a \emph{thick copper strip} to a metal plate buried deep in the ground?
\item How strong is the electrostatic attraction between the cloud base and the rooftop charge, and does the humid air of the rainy season make it stronger or weaker?
\item What does the electric field look like between the cloud and the building, and how does the conductor change it?
\end{enumerate}

\courssub{Tasks}

\textbf{Task 1: Charging by induction (explain).}
\begin{enumerate}
\item State the two types of electric charge and the basic law of electrostatics.
\item Using clear diagrams (cloud, rooftop, ground), describe step by step how the negatively charged cloud base induces a positive charge on the rooftop. Name this method of charging and state what happens to the electrons in the conducting parts of the building.
\item Explain why this induction is only possible because the building contains conducting materials (metal roof sheets, water pipes, reinforcing iron). What would be different if the block were a perfect insulator?
\end{enumerate}

\textbf{Task 2: Design of the Franklin rod (justify).}
\begin{enumerate}
\item State what is meant by \cle{point action}. Explain why the charge density and the electric field are greatest at sharp points of a charged conductor.
\item Deduce why the lightning conductor is topped with a sharp spike rather than a ball.
\item Explain the role of the thick copper down-conductor and of the earth plate buried in moist ground: what happens to the induced charge, and why must the conductor be thick and continuous?
\end{enumerate}

\textbf{Task 3: Coulomb's law (calculate).}
\begin{enumerate}
\item Write down Coulomb's law for the force between two point charges in a medium of relative permittivity $\varepsilon_r$.
\item Calculate the magnitude of the force between the cloud base and the induced rooftop charge in dry air ($\varepsilon_r = 1.00$). State whether the force is attractive or repulsive.
\item Repeat the calculation for humid air ($\varepsilon_r = 1.01$). Compute the percentage change and comment: does humidity increase or decrease the force? Is the effect large?
\item Calculate the electric field strength $E$ created by the cloud base at the position of the rooftop ($d = 300\ \mathrm{m}$, dry air).
\end{enumerate}

\textbf{Task 4: Field lines and protection (sketch and interpret).}
\begin{enumerate}
\item Sketch the electric field lines between the negatively charged cloud base and the positively charged rooftop \emph{without} a lightning conductor. Indicate the direction of the field lines.
\item Sketch the situation \emph{with} the pointed conductor installed. Show how the field lines concentrate at the tip of the spike.
\item Using your sketches, explain to the committee how the conductor protects the block: (i) by slowly discharging the induced charge through point action, and (ii) by offering lightning a safe, low-resistance path to the ground if a strike does occur.
\end{enumerate}

\textbf{Task 5: Technical note.} Assemble your answers into a one-page technical note addressed to the PTA committee: brief introduction, labelled diagrams, calculations with units, and a clear recommendation. Each group presents its note to the class (the ``PTA committee'') in five minutes.

\courssub{Model Answer}

\textbf{Task 1: Charging by Induction.}

\begin{definition}[Electric Charge and the Basic Law of Electrostatics]
There are two types of electric charge: \cle{positive} ($+$) and \cle{negative} ($-$). The \cle{basic law of electrostatics} states that \emph{like charges repel each other and unlike charges attract each other}. Charge is quantised (multiple of $e = 1.60\times 10^{-19}\ \mathrm{C}$) and conserved.
\end{definition}

\begin{propriete}[Charging by Induction]
When a charged object is brought near a \emph{conductor} without touching it, the free electrons in the conductor move in response to the external electric field. This creates a separation of charge: the near side acquires a charge opposite to the inducer, the far side acquires a charge of the same sign. The conductor as a whole remains electrically neutral unless it is earthed.
\end{propriete}

The building contains conductors (metal roofing sheets, water pipes, reinforcing steel bars) in which \emph{free electrons} can move. When the negatively charged cloud base ($Q_1 < 0$) approaches:

\begin{itemize}
\item The cloud's negative charge \emph{repels} the free electrons in the building's conducting framework.
\item The electrons are pushed as far away as possible, down through the structure into the ground (the Earth acts as an infinite reservoir of charge).
\item The rooftop, having \emph{lost} electrons, is left with a net \emph{positive} charge $Q_2 > 0$.
\end{itemize}

The rooftop has been charged \textbf{by induction}: it acquired a charge of sign \emph{opposite} to that of the influencing charge, \emph{without any contact} with the cloud.

\begin{attention}[Insulator vs Conductor]
If the block were a perfect insulator (e.g. dry wood, plastic), its electrons are tightly bound and cannot move freely. No significant charge separation would occur, no net induced charge would appear on the rooftop, and the electrostatic attraction would be negligible. The presence of conductors is essential for the induction process described here.
\end{attention}

\textbf{Task 2: Design of the Franklin Rod.}

\begin{definition}[Point Action]
On a charged conductor in electrostatic equilibrium, the surface charge density $\sigma$ is greatest where the radius of curvature is smallest, i.e.\ at \emph{sharp points}. Since the electric field just outside a conductor is $E = \sigma/\varepsilon_0$, the field is extremely intense near a point. This phenomenon is called \cle{point action}. The intense field ionises the surrounding air molecules, creating a corona discharge that allows charge to leak into (or out of) the air.
\end{definition}

\begin{propriete}[Design Principles of a Lightning Conductor]
\begin{itemize}
\item \textbf{Sharp spike:} The strong field at the tip ionises the air, allowing the induced rooftop charge to be discharged \emph{gently and continuously} towards the cloud (corona discharge), reducing the potential difference and the chance of a sudden, destructive strike. If a strike does occur, the tip is the most probable point of attachment because the field there exceeds the dielectric strength of air first.
\item \textbf{Thick copper down-conductor:} Copper has a very low resistivity ($\rho \approx 1.7\times 10^{-8}\ \Omega\,\mathrm{m}$). A thick strip offers a very low resistance path, so that an enormous current (up to $200\ \mathrm{kA}$) can pass to the ground in a few microseconds without melting the conductor (Joule heating $I^2Rt$) or causing dangerous side-flashes to the building structure.
\item \textbf{Earth plate in moist ground:} Moist soil has a relatively low resistivity. A large buried plate increases the contact area, lowering the earth resistance. The Earth acts as an effectively infinite reservoir of charge: electrons flow up from the earth to neutralise the induced positive charge during the storm, or a lightning current flows down into the earth harmlessly.
\end{itemize}
\end{propriete}

\begin{savaistu}[Benjamin Franklin and Cameroon]
Benjamin Franklin (1705--1790) invented the lightning rod in 1752 after his famous kite experiment. In Cameroon, the rainy season (March--October in the North-West) brings frequent thunderstorms. The MINESEC guidelines require all public buildings, especially those housing expensive equipment like computer labs, to be fitted with certified lightning protection systems conforming to the NFC 17-102 / IEC 62305 standards.
\end{savaistu}

\textbf{Task 3: Coulomb's Law Calculations.}

\begin{propriete}[Coulomb's Law in a Dielectric Medium]
The magnitude of the electrostatic force between two point charges $Q_1$ and $Q_2$ separated by a distance $r$ in a homogeneous, isotropic medium of relative permittivity $\varepsilon_r$ is:
\[
F = \frac{1}{4\pi\varepsilon_0\varepsilon_r}\,\frac{|Q_1 Q_2|}{r^2} = \frac{k}{\varepsilon_r}\,\frac{|Q_1 Q_2|}{r^2}
\]
where $k = 9.0\times 10^{9}\ \mathrm{N\,m^2\,C^{-2}}$. The force is attractive if the charges have opposite signs, repulsive if they have the same sign. The direction is along the line joining the charges.
\end{propriete}

\textbf{(i) Dry air ($\varepsilon_r = 1.00$):}
\[
F = \frac{9.0\times 10^{9}\times |(-20)\times (2.0\times 10^{-3})|}{(300)^2}
= \frac{9.0\times 10^{9}\times 4.0\times 10^{-2}}{9.0\times 10^{4}}
= \frac{3.6\times 10^{8}}{9.0\times 10^{4}}
= 4.0\times 10^{3}\ \mathrm{N}.
\]
The charges have opposite signs ($Q_1<0$, $Q_2>0$), so the force is \textbf{attractive}. The force on the rooftop charge is directed \emph{vertically upwards} towards the cloud base.

\textbf{(ii) Humid air ($\varepsilon_r = 1.01$):}
\[
F' = \frac{F}{\varepsilon_r} = \frac{4.0\times 10^{3}}{1.01} = 3.96\times 10^{3}\ \mathrm{N}\quad(\text{3 s.f.}).
\]
Percentage change:
\[
\frac{F' - F}{F}\times 100\% = \left(\frac{1}{1.01} - 1\right)\times 100\% \approx -0.99\% \approx -1.0\%.
\]
The humid air \emph{slightly decreases} the force (by about $1\%$). The force is inversely proportional to $\varepsilon_r$. The effect is small because the permittivity of air, even when saturated with water vapour, remains very close to that of vacuum ($\varepsilon_r \approx 1$).

\textbf{(iii) Electric field strength at the rooftop due to the cloud base (dry air):}
\[
E = \frac{k|Q_1|}{\varepsilon_r d^2} = \frac{9.0\times 10^{9}\times 20}{1.00 \times (300)^2}
= \frac{1.8\times 10^{11}}{9.0\times 10^{4}}
= 2.0\times 10^{6}\ \mathrm{N\,C^{-1}}\ (\text{or}\ \mathrm{V\,m^{-1}}).
\]
The direction is \emph{upwards} (towards the negative cloud), because electric field lines point \emph{towards} negative charges.

\begin{aretenir}[Summary of Numerical Results]
\begin{itemize}
\item Force (dry air): $F = 4.0\times 10^{3}\ \mathrm{N}$ (attractive).
\item Force (humid air): $F' = 3.96\times 10^{3}\ \mathrm{N}$ (attractive), $\approx 1\%$ weaker.
\item Electric field at rooftop: $E = 2.0\times 10^{6}\ \mathrm{N\,C^{-1}}$ directed upwards.
\end{itemize}
\end{aretenir}

\textbf{Task 4: Field Lines and Protection.}

\begin{experience}[Sketching Electric Field Lines]
\textbf{Without the conductor:} Field lines originate on the induced positive charges (rooftop, ground) and terminate on the negative cloud base. They are roughly parallel and uniform over the flat roof surface, indicating a fairly uniform field $E \approx 2.0\times 10^{6}\ \mathrm{N\,C^{-1}}$. The whole roof is exposed to this field.

\textbf{With the pointed conductor:} The conductor, being earthed, is at zero potential. The sharp tip has a very small radius of curvature, so the field lines crowd together intensely at the tip. The field at the tip far exceeds the dielectric strength of air ($\approx 3\times 10^{6}\ \mathrm{V\,m^{-1}}$), causing corona discharge.
\end{experience}

\begin{popfigure}
\begin{center}
\begin{tikzpicture}[scale=0.85]
% Cloud
\draw[fill=popblueL, draw=popblue, thick] (-3,3.5) ellipse (2.5 and 0.8);
\node[popdark] at (-3,3.5) {$Q_1 = -20\ \mathrm{C}$};
% Building
\draw[fill=popmuted!30, draw=popdark, thick] (-4.5,0) rectangle (-1.5,1.8);
% Roof detail
\draw[popdark, thick] (-4.7,1.8) -- (-3,2.6) -- (-1.3,1.8);
% Lightning conductor (rod)
\draw[poporange, very thick] (-3,2.6) -- (-3,3.3);
\fill[poporange] (-3,3.35) circle (0.06);
\node[poporange, right] at (-3,3.35) {\small Sharp tip};
% Down conductor
\draw[poporange, thick] (-3,0) -- (-3,-0.8);
\draw[poporange, thick] (-3,-0.8) -- (-5.5,-0.8);
\node[popdark, below] at (-4.5,-0.8) {Earth plate};
% Ground line
\draw[popdark, very thick] (-6,0) -- (1,0);
% Field lines WITHOUT rod (faint, dashed) - for comparison
\foreach \x in {-4.8, -4.0, -2.0, -1.2} {
  \draw[->, popgreen!50!popmuted, dashed, thick] (\x,0.15) -- (\x,2.8);
}
% Field lines WITH rod (solid, concentrated at tip)
\foreach \x in {-4.8, -4.0, -2.0, -1.2} {
  \draw[->, popgreen, thick] (\x,0.15) -- (\x,2.5);
}
% Concentrated lines at tip
\draw[->, popgreen, very thick] (-3,2.65) -- (-3,3.1);
\draw[->, popgreen, thick] (-3.2,2.65) -- (-3.1,3.1);
\draw[->, popgreen, thick] (-2.8,2.65) -- (-2.9,3.1);
% Labels
\node[popgreen, right, text width=4cm, align=left] at (0.5,1.5) {Field lines converge\\ on sharp tip};
\node[popgreen!50!popmuted, right, text width=4cm, align=left] at (0.5,0.5) {Uniform field without rod (dashed)};
\node[popdark, anchor=north] at (-3, -1.2) {Science Block ($h=8\ \mathrm{m}$)};
\node[popdark, anchor=south] at (-3, 3.8) {Thundercloud ($d=300\ \mathrm{m}$)};
\end{tikzpicture}
\end{center}
\legende{Electric field lines between the charged cloud and the protected rooftop. Solid lines: with Franklin rod (concentration at tip). Dashed lines: without rod (uniform field).}
\end{popfigure}

\textbf{Protection Mechanism (for the Committee):}

\begin{enumerate}
\item \textbf{Prevention (Corona Discharge):} The intense field at the sharp tip ($E_{\text{tip}} \gg E_{\text{roof}}$) ionises the surrounding air. Positive ions are repelled upwards towards the negative cloud, while electrons are attracted to the tip from the ground. This creates a small \emph{corona current} that gradually neutralises the induced charge on the rooftop, lowering the potential difference between cloud and building. This reduces the probability of a lightning strike initiating.
\item \textbf{Safe Channelling (Lightning Strike):} If the electric field builds up too rapidly for corona discharge to keep pace (e.g. during a very violent storm), a lightning strike occurs. The strike attaches preferentially to the tip (highest field). The massive current ($I \sim 10^4$--$10^5\ \mathrm{A}$) flows down the thick copper strip (low $R$, low $L$) to the earth plate, bypassing the building structure, the wiring, and the computer laboratory. The building and its $8\,000\,000$ FCFA equipment remain safe.
\end{enumerate}

\textbf{Task 5: Technical Note to the PTA Committee.}

\begin{center}
\textbf{\large TECHNICAL NOTE}\\[0.5em]
\textbf{To:} PTA Committee, GBHS [School Name]\\[0.2em]
\textbf{From:} Upper Sixth Physics Consultancy Group\\[0.2em]
\textbf{Date:} \today\\[0.2em]
\textbf{Subject:} Lightning Protection for New Science Block
\end{center}

\vspace{0.5em}
\noindent\textbf{1. Risk Assessment.}
The science block ($h=8\ \mathrm{m}$) on a hilltop is a prime target for lightning. Meteorological data estimates a cloud base charge $Q_1=-20\ \mathrm{C}$ at $d=300\ \mathrm{m}$. This induces $Q_2=+2\ \mathrm{mC}$ on the rooftop, creating an attractive force of $4.0\ \mathrm{kN}$ and an electric field of $2.0\ \mathrm{MV/m}$ at the roof. This field approaches the breakdown strength of air ($3\ \mathrm{MV/m}$), posing a high risk of a direct strike.

\noindent\textbf{2. Protection Principle.}
We recommend a \textbf{Franklin Rod System} based on three physics principles:
\begin{itemize}
\item \emph{Induction:} The cloud charge induces opposite charge on the earthed conductor.
\item \emph{Point Action:} A sharp tip ($r_{\text{tip}} \approx 1\ \mathrm{mm}$) creates a local field $>3\ \mathrm{MV/m}$, initiating corona discharge to leak charge safely.
\item \emph{Safe Conduction:} A thick copper tape ($50\ \mathrm{mm}^2$ cross-section) provides a low-impedance path ($R < 0.1\ \Omega$) for any strike current to a buried copper earth plate ($1\ \mathrm{m}^2$, $1\ \mathrm{m}$ deep in moist soil).
\end{itemize}

\noindent\textbf{3. Calculations.}
\begin{align*}
F_{\text{dry}} &= \frac{k|Q_1Q_2|}{d^2} = 4.0\times 10^{3}\ \mathrm{N}\ (\text{attractive}) \\
F_{\text{humid}} &= F_{\text{dry}} / 1.01 = 3.96\times 10^{3}\ \mathrm{N}\ (-1.0\%\ \text{change}) \\
E_{\text{roof}} &= k|Q_1|/d^2 = 2.0\times 10^{6}\ \mathrm{N\,C^{-1}}
\end{align*}
Humidity slightly reduces the force but does not eliminate the risk.

\noindent\textbf{4. Recommendation.}
Install a pointed air terminal (Franklin rod) at the highest point of the roof, bonded to a continuous $25\times 3\ \mathrm{mm}$ copper down-conductor fixed to the walls, terminating in a $1\ \mathrm{m}^2$ copper earth plate buried at $1\ \mathrm{m}$ depth in a moist location. Interconnect all metal services (pipes, window frames) to the down-conductor to prevent side-flashes. Estimated cost: $\ll 1\%$ of equipment value. This system, compliant with IEC 62305, will safeguard the laboratory throughout the rainy season.

\vspace{1em}
\noindent\textit{End of Note.}

\newpage

\coursec{Methods and worked examples}

\coursec{Electrostatics and Electric Fields}

\courssub{Electric Charge: Origin, Types and the Basic Law of Electrostatics}

\begin{definition}[Electric Charge]
Electric charge is a fundamental property of matter that determines its electromagnetic interaction. It exists in two types, conventionally called \textbf{positive} and \textbf{negative}. The SI unit of charge is the \textbf{coulomb (C)}.
\end{definition}

\begin{propriete}[Elementary Charge and Quantisation]
Electric charge is \cle{quantised}: any free charge $q$ is an integer multiple of the elementary charge $e$:
\[ q = n\,e \qquad (n \in \mathbb{Z}), \quad e = 1.60 \times 10^{-19}\ \mathrm{C}. \]
The electron carries charge $-e$, the proton $+e$. In solid conductors, only electrons move; positive charges (atomic nuclei) remain fixed in the lattice.
\end{propriete}

\begin{propriete}[Conservation of Charge]
In any isolated system, the total electric charge remains constant. Charge can be transferred from one body to another, but it can never be created or destroyed.
\end{propriete}

\begin{propriete}[Basic Law of Electrostatics]
\begin{itemize}
\item \textbf{Like charges repel} (positive-positive or negative-negative).
\item \textbf{Unlike charges attract} (positive-negative).
\end{itemize}
\textbf{Important:} attraction can also occur between a charged body and a neutral one (by induction). Only \textbf{repulsion} is a conclusive test for the sign of a charge.
\end{propriete}

\begin{popfigure}
\begin{tikzpicture}[scale=1.2]
% Atom diagram
\draw[popblue, thick] (0,0) circle (1.2);
\draw[poporange, thick] (0,0) circle (0.3);
\node at (0,0) {\small Nucleus ($+Ze$)};
\foreach \angle in {0,60,120,180,240,300} {
  \draw[popgreen, ->] (0,0) -- (\angle:1.0);
  \node at (\angle:1.15) {\small $e^-$};
}
\node at (0,-1.6) {Neutral atom: $Z$ protons, $Z$ electrons};
% Charged bodies
\begin{scope}[xshift=5cm]
\draw[popblue, thick] (0,0) circle (1.2);
\draw[poporange, thick] (0,0) circle (0.3);
\node at (0,0) {\small Nucleus ($+Ze$)};
\foreach \angle in {0,90,180,270} {
  \draw[popgreen, ->] (0,0) -- (\angle:1.0);
  \node at (\angle:1.15) {\small $e^-$};
}
\node at (0,-1.6) {Positive ion: lost electrons};
\end{scope}
\begin{scope}[xshift=10cm]
\draw[popblue, thick] (0,0) circle (1.2);
\draw[poporange, thick] (0,0) circle (0.3);
\node at (0,0) {\small Nucleus ($+Ze$)};
\foreach \angle in {0,45,90,135,180,225,270,315} {
  \draw[popgreen, ->] (0,0) -- (\angle:1.0);
  \node at (\angle:1.15) {\small $e^-$};
}
\node at (0,-1.6) {Negative ion: gained electrons};
\end{scope}
\end{tikzpicture}
\legende{Atomic origin of charge: neutral atom, positive ion (electron deficit), negative ion (electron excess).}
\end{popfigure}

\courssub{Conductors and Insulators}

\begin{definition}[Conductors and Insulators]
\begin{itemize}
\item A \cle{conductor} (metal, graphite, electrolyte) contains \textbf{free charge carriers} (electrons in metals, ions in electrolytes) that can move through the material.
\item An \cle{insulator} (dielectric) has \textbf{no free charge carriers}; charges deposited on it remain localised.
\end{itemize}
\end{definition}

\begin{propriete}[Behaviour of Conductors in Electrostatics]
\begin{enumerate}
\item Excess charge resides entirely on the \textbf{outer surface}.
\item The electric field \textbf{inside} the conductor is zero.
\item The conductor is an \textbf{equipotential} body.
\item Surface charge density is highest where the curvature is greatest (sharp points) --- this is \cle{point action}.
\end{enumerate}
\end{propriete}

\begin{popfigure}
\begin{tikzpicture}[scale=1.2]
% Conductor with charge
\draw[popblue, thick] (0,0) ellipse (1.5 and 1);
\foreach \x in {-1.2,-0.6,0,0.6,1.2} {
  \node[popred, circle, fill=popred, inner sep=1.5pt] at (\x,0.8) {};
  \node[popred, circle, fill=popred, inner sep=1.5pt] at (\x,-0.8) {};
}
\draw[->, popgreen] (0,0) -- (0,1.5) node[above] {$\vec{E}=0$};
\node at (0,-1.5) {Conductor: charge on surface, $\vec{E}_{\text{int}}=0$};
% Insulator with charge
\begin{scope}[xshift=5cm]
\draw[poporange, thick] (0,0) ellipse (1.5 and 1);
\foreach \x in {-1.2,-0.6,0,0.6,1.2} {
  \foreach \y in {-0.6,0,0.6} {
    \node[popred, circle, fill=popred, inner sep=1.5pt] at (\x,\y) {};
  }
}
\node at (0,-1.5) {Insulator: charge fixed throughout volume};
\end{scope}
\end{tikzpicture}
\legende{Distribution of excess charge on a conductor vs an insulator.}
\end{popfigure}

\courssub{Methods of Charging}

\begin{methode}[Analysing a Charging Experiment]
When a problem describes an electrostatics experiment (rubbing, contact, induction), proceed as follows.
\begin{enumerate}
\item \textbf{Identify the materials} involved and classify them: conductor or insulator. Remember that charge can only move through a \cle{conductor}; on an \cle{insulator} it stays where it is deposited.
\item \textbf{Identify the mechanism}:
\begin{itemize}
\item \emph{Friction (rubbing)}: electrons are torn from one body and transferred to the other; the two bodies acquire \textbf{equal and opposite} charges. The triboelectric series predicts which material becomes positive/negative.
\item \emph{Contact}: a charged body touches a conductor; charge is \textbf{shared} between them (same sign as the original charge). For identical conductors, $q_{\text{final}} = \frac{q_1+q_2}{2}$.
\item \emph{Induction (influence)}: a charged body is brought \textbf{near} a conductor \textbf{without touching}; charges inside the conductor separate, and with an earth connection the conductor can be charged with the \textbf{opposite} sign.
\end{itemize}
\item \textbf{Apply conservation of charge}: the total charge of an isolated system is constant. Only \cle{electrons} move; positive charges (protons in nuclei) never do.
\item \textbf{Conclude} on the sign of the charge on each body, and justify any observation (attraction, repulsion, movement of gold leaves).
\end{enumerate}
\textbf{Reflex:} attraction alone does \emph{not} prove that two bodies carry opposite charges --- a charged body always attracts a light neutral body by induction. Only \textbf{repulsion} is a sure test of charges of the same sign.
\end{methode}

\begin{experience}[Charging an Electroscope by Induction]
\textbf{Apparatus:} gold-leaf electroscope, ebonite rod, wool cloth, metal plate (earth).
\textbf{Procedure:}
\begin{enumerate}
\item Charge the ebonite rod negatively by rubbing with wool.
\item Bring the rod near (without touching) the cap of a neutral electroscope. Observe leaf divergence.
\item While the rod is held in place, touch the cap with a finger (earth connection). Leaves collapse.
\item Remove the finger \textbf{first}, then take away the rod. Leaves diverge again.
\end{enumerate}
\textbf{Observation:} the electroscope is left positively charged (opposite to the inducing rod).
\end{experience}

\begin{exoresolu}[Charging an Electroscope by Induction]
A negatively charged ebonite rod is brought near (without touching) the cap of a neutral gold-leaf electroscope. The leaves diverge. While the rod is held in place, the cap is touched with a finger; the leaves collapse. The finger is removed first, then the rod is taken away: the leaves diverge again.
\begin{enumerate}
\item Explain why the leaves diverge when the rod is first brought near.
\item State, with reasons, the sign of the charge left on the electroscope at the end.
\end{enumerate}
\tcblower
\textbf{1. Divergence of the leaves.} The ebonite rod carries an excess of electrons (negative charge). Brought near the metal cap, it exerts a repulsive force on the free electrons of the metal (cap, stem and leaves form one conductor). The electrons are pushed as far from the rod as possible, i.e. down into the leaves. The cap is left with a deficit of electrons (positive), and the two leaves both become \textbf{negative}. Carrying charges of the same sign, the leaves repel each other and diverge. This is charging \emph{by induction}: no contact, hence no transfer of charge between rod and electroscope; the total charge of the electroscope is still zero at this stage.

\textbf{2. Final charge.} With the rod still in place, the finger connects the electroscope to the earth, which acts as an enormous reservoir of charge. The repelled electrons now escape to earth rather than accumulate in the leaves, so the leaves collapse. The electroscope is then positive at the cap (held there by the attraction of the rod) and neutral elsewhere. Removing the finger \emph{first} breaks the earth connection while the rod still holds the positive charge on the cap. When the rod is finally removed, this positive charge redistributes over the whole conductor (cap, stem, leaves). Both leaves become positive, repel each other, and diverge again.

\textbf{Conclusion:} the electroscope is left \textbf{positively charged}, i.e. with a charge of sign \emph{opposite} to that of the inducing rod --- the characteristic feature of charging by induction. Note the essential order of operations: the finger must be removed \emph{before} the rod, otherwise the charge would simply flow back to earth and the electroscope would remain neutral.
\end{exoresolu}

\begin{popfigure}
\begin{tikzpicture}[scale=1.0]
% Step 1: rod near neutral electroscope
\begin{scope}[yshift=0]
\node[align=center] at (0,3.5) {\textbf{Step 1: Rod brought near}};
% Electroscope
\draw[popblue, thick] (0,0) -- (0,2.5); % stem
\draw[popblue, thick] (-0.5,2.5) -- (0.5,2.5); % cap
\draw[popgreen, thick] (0,1.5) -- (-0.8,0.5); % left leaf
\draw[popgreen, thick] (0,1.5) -- (0.8,0.5); % right leaf
% Charges on electroscope
\foreach \x in {-0.4,0,0.4} \node[popred, circle, fill=popred, inner sep=1.5pt] at (\x,2.6) {}; % positive on cap
\foreach \x in {-0.6,-0.2,0.2,0.6} \node[popblue, circle, fill=popblue, inner sep=1.5pt] at (\x,0.7) {}; % negative on leaves
% Rod
\draw[poporange, thick] (2.5,2.5) -- (2.5,0.5);
\foreach \y in {0.7,1.1,1.5,1.9,2.3} \node[popblue, circle, fill=popblue, inner sep=1.5pt] at (2.5,\y) {};
\node at (2.5,3.0) {$-$ rod};
\node at (0,-0.5) {Leaves diverge (both $-$)}; 
\end{scope}
% Step 2: finger touches cap
\begin{scope}[xshift=6cm]
\node[align=center] at (0,3.5) {\textbf{Step 2: Finger touches cap (earth)}};
\draw[popblue, thick] (0,0) -- (0,2.5);
\draw[popblue, thick] (-0.5,2.5) -- (0.5,2.5);
\draw[popgreen, thick] (0,1.5) -- (-0.3,1.2); % leaves collapsed
\draw[popgreen, thick] (0,1.5) -- (0.3,1.2);
\foreach \x in {-0.4,0,0.4} \node[popred, circle, fill=popred, inner sep=1.5pt] at (\x,2.6) {}; % positive on cap
\draw[popmuted, thick] (0,2.5) -- (-1.5,2.5) node[left] {Finger (earth)};
\draw[poporange, thick] (2.5,2.5) -- (2.5,0.5);
\foreach \y in {0.7,1.1,1.5,1.9,2.3} \node[popblue, circle, fill=popblue, inner sep=1.5pt] at (2.5,\y) {};
\node at (2.5,3.0) {$-$ rod};
\node at (0,-0.5) {Leaves collapse};
\end{scope}
% Step 3: finger removed, then rod
\begin{scope}[xshift=12cm]
\node[align=center] at (0,3.5) {\textbf{Step 3: Finger removed, then rod}};
\draw[popblue, thick] (0,0) -- (0,2.5);
\draw[popblue, thick] (-0.5,2.5) -- (0.5,2.5);
\draw[popgreen, thick] (0,1.5) -- (-0.8,0.5);
\draw[popgreen, thick] (0,1.5) -- (0.8,0.5);
\foreach \x in {-0.6,-0.2,0.2,0.6} \node[popred, circle, fill=popred, inner sep=1.5pt] at (\x,0.7) {}; % positive on leaves
\foreach \x in {-0.4,0,0.4} \node[popred, circle, fill=popred, inner sep=1.5pt] at (\x,2.6) {}; % positive on cap
\node at (0,-0.5) {Leaves diverge again (both $+$)};
\end{scope}
\end{tikzpicture}
\legende{Sequence of charging a gold-leaf electroscope by induction using a negatively charged rod.}
\end{popfigure}

\courssub{Point Action and the Lightning Conductor}

\begin{methode}[Using Point Action in an Explanation]
\begin{enumerate}
\item \textbf{Recall the fact}: on a charged conductor, charge accumulates where the curvature is greatest; the surface charge density, and hence the field just outside, is \textbf{largest at sharp points}.
\item \textbf{Link to ionisation}: near a sharp point the field can exceed the dielectric strength of air ($\approx 3 \times 10^{6}\ \mathrm{V\,m^{-1}}$), ionising the air: a \emph{corona discharge} leaks charge away from the point.
\item \textbf{Apply to the situation}: lightning conductor (pointed rod earthed by a thick conductor), pointed electrodes, or why a charged conductor loses charge faster if it has sharp edges.
\item \textbf{Conclude} on the protective or undesirable role of the point.
\end{enumerate}
\end{methode}

\begin{exoresolu}[How the Lightning Conductor Protects a Building]
During the rainy season in Cameroon, tall buildings are fitted with a lightning conductor: a sharp metal spike at the top of the building, connected by a thick copper cable to a metal plate buried in the ground. Explain, using the ideas of induction and point action, how this device protects the building during a thunderstorm.
\tcblower
\textbf{Situation.} The base of a thundercloud carries a large charge (say negative). By \textbf{induction}, it repels the free electrons of the ground and of the building: the top of the building, and in particular the spike, becomes positively charged while negative charge is pushed into the earth.

\textbf{Point action.} Because the spike is sharply pointed, the induced positive charge concentrates there with a very high surface density. The electric field just outside the point becomes extremely intense --- far stronger than above a flat roof --- and can exceed the dielectric strength of the surrounding air.

\textbf{Corona discharge.} The air around the point is ionised. Positive ions are repelled from the spike and drift towards the cloud, while electrons are attracted to the spike and conducted harmlessly to earth. This gentle, continuous leakage of charge (\emph{point discharge}) progressively \textbf{reduces the charge difference} between cloud and building, lowering the probability that a violent lightning strike will occur at all.

\textbf{If lightning does strike.} Should the discharge nevertheless happen, it preferentially hits the spike (the highest, sharpest, best-connected point). The enormous current is then channelled through the low-resistance copper cable directly into the earth, instead of passing through the roof, walls or electrical wiring of the building, where it would cause fire and destruction.

\textbf{Conclusion.} The lightning conductor protects in two complementary ways: by \emph{point action} it quietly discharges the induced charge and may prevent the strike; and if the strike occurs, it offers it a \emph{safe conducting path to earth}. This is why the cable must be thick, continuous and well earthed: any break or high resistance would force the current through the building itself.
\end{exoresolu}

\begin{popfigure}
\begin{tikzpicture}[scale=0.8]
% Cloud
\draw[popblue, thick, decorate, decoration={snake, amplitude=2pt, segment length=10pt}] (-3,4) -- (3,4);
\foreach \x in {-2.5,-1.5,-0.5,0.5,1.5,2.5} \node[popblue, circle, fill=popblue, inner sep=2pt] at (\x,4.2) {};
\node at (0,4.8) {Thundercloud ($-$ charge at base)};
% Building
\draw[popdark, thick] (-1.5,0) -- (-1.5,3) -- (1.5,3) -- (1.5,0);
\draw[popdark, thick] (-2,0) -- (2,0);
% Lightning conductor
\draw[popgold, very thick] (0,3) -- (0,4.5);
\draw[popgold, thick] (0,0) -- (0,-1.5);
\draw[popgold, thick] (-0.5,-1.5) -- (0.5,-1.5);
\node[popgold] at (0,4.7) {$\triangle$}; % spike
% Induced charges
\foreach \y in {2.5,2.8,3.1} \node[popred, circle, fill=popred, inner sep=1.5pt] at (0,\y) {};
\node[popred] at (1.2,3) {$+$};
\node[popblue] at (1.2,0) {$-$ (to earth)};
% Corona discharge
\draw[poporange, ->, thick] (0,4.5) .. controls (0.5,4.8) and (1,4.5) .. (1.5,4.2);
\draw[poporange, ->, thick] (0,4.5) .. controls (-0.5,4.8) and (-1,4.5) .. (-1.5,4.2);
\node[poporange] at (0,5.0) {Corona discharge (point action)};
% Lightning strike (dashed)
\draw[poppink, dashed, thick] (-2,4) .. controls (-1,3) .. (-1.5,3);
\draw[poppink, dashed, thick] (2,4) .. controls (1,3) .. (1.5,3);
\node[poppink] at (-2.5,3.5) {Lightning prefers spike};
\end{tikzpicture}
\legende{Lightning conductor: induction charges the spike, point action creates corona discharge, and the thick cable provides a safe path to earth.}
\end{popfigure}

\courssub{Coulomb's Law and the Inverse-Square Nature}

\begin{propriete}[Coulomb's Law]
The electrostatic force $\vec{F}_{12}$ exerted by a point charge $q_1$ on a point charge $q_2$ separated by distance $r$ in vacuum (or air) is:
\[ \vec{F}_{12} = k\,\frac{q_1 q_2}{r^2}\,\hat{r}_{12}, \qquad k = \frac{1}{4\pi\varepsilon_0} \approx 9.0 \times 10^{9}\ \mathrm{N\,m^2\,C^{-2}}, \]
where $\hat{r}_{12}$ is the unit vector from $q_1$ to $q_2$. The force is \textbf{attractive} if $q_1 q_2 < 0$ and \textbf{repulsive} if $q_1 q_2 > 0$. Newton's third law holds: $\vec{F}_{21} = -\vec{F}_{12}$.
\end{propriete}

\begin{propriete}[Effect of a Dielectric Medium]
In a homogeneous, isotropic dielectric medium of relative permittivity $\varepsilon_r$ (dielectric constant), the force is reduced by a factor $\varepsilon_r$:
\[ F_{\text{medium}} = \frac{F_{\text{vacuum}}}{\varepsilon_r} = \frac{1}{4\pi\varepsilon_0\varepsilon_r}\,\frac{|q_1 q_2|}{r^2}. \]
Typical values: air $\varepsilon_r \approx 1.0006 \approx 1$, kerosene $\approx 2$, water $\approx 80$.
\end{propriete}

\begin{methode}[Computing the Force Between Two Point Charges]
\begin{enumerate}
\item \textbf{Check the conditions}: the charges must be point charges (or spherically symmetric distributions, which behave as if all the charge were at the centre).
\item \textbf{Convert all data to SI units}: charges in coulombs ($1\ \mu\mathrm{C} = 10^{-6}\ \mathrm{C}$; $1\ \mathrm{nC} = 10^{-9}\ \mathrm{C}$), distances in metres.
\item \textbf{Write Coulomb's law}:
\[ F = k\,\frac{|q_1 q_2|}{r^2}, \qquad k = \frac{1}{4\pi\varepsilon_0} \approx 9.0\times 10^{9}\ \mathrm{N\,m^2\,C^{-2}} \quad (\text{in vacuum/air}). \]
\item \textbf{Compute the magnitude} with absolute values; determine the \textbf{direction} separately: same signs $\Rightarrow$ repulsion; opposite signs $\Rightarrow$ attraction. The force acts along the line joining the charges.
\item \textbf{In a medium} other than vacuum, replace $\varepsilon_0$ by $\varepsilon = \varepsilon_0 \varepsilon_r$: the force is divided by the relative permittivity $\varepsilon_r$.
\end{enumerate}
\textbf{Reflexes:} the force varies as $1/r^2$: doubling $r$ divides $F$ by 4; doubling \emph{both} charges multiplies $F$ by 4. Newton's third law applies: the two forces are equal in magnitude even if the charges are very different.
\end{methode}

\begin{exoresolu}[Force Between Two Charged Spheres]
Two small identical conducting spheres carry charges $q_1 = +4.0\ \mu\mathrm{C}$ and $q_2 = -2.0\ \mu\mathrm{C}$. Their centres are $r = 30\ \mathrm{cm}$ apart in air. Take $k = 9.0\times 10^{9}\ \mathrm{N\,m^2\,C^{-2}}$.
\begin{enumerate}
\item Calculate the magnitude of the electrostatic force between them and state its nature.
\item The spheres are briefly connected by a thin wire, then separated again to $30\ \mathrm{cm}$. Calculate the new force and state its nature.
\end{enumerate}
\tcblower
\textbf{1. Force before contact.} Convert the data:
\[ q_1 = 4.0\times 10^{-6}\ \mathrm{C}, \quad q_2 = -2.0\times 10^{-6}\ \mathrm{C}, \quad r = 0.30\ \mathrm{m}. \]
Apply Coulomb's law with the absolute values:
\[ F = k\,\frac{|q_1 q_2|}{r^2} = 9.0\times 10^{9}\times \frac{(4.0\times 10^{-6})(2.0\times 10^{-6})}{(0.30)^2}. \]
Numerator: $9.0\times 10^{9}\times 8.0\times 10^{-12} = 7.2\times 10^{-2}$. Denominator: $(0.30)^2 = 0.090$. Hence
\[ F = \frac{7.2\times 10^{-2}}{9.0\times 10^{-2}} = 0.80\ \mathrm{N}. \]
The charges have opposite signs, so the force is \textbf{attractive}: each sphere is pulled towards the other along the line joining their centres.

\textbf{2. After contact.} The two spheres are identical conductors joined by a wire: the total charge is shared equally between them. By conservation of charge:
\[ q_{\text{total}} = q_1 + q_2 = +4.0 - 2.0 = +2.0\ \mu\mathrm{C}, \qquad q'_1 = q'_2 = \frac{+2.0}{2} = +1.0\ \mu\mathrm{C}. \]
The new force is
\[ F' = k\,\frac{q'_1 q'_2}{r^2} = 9.0\times 10^{9}\times \frac{(1.0\times 10^{-6})^2}{(0.30)^2} = \frac{9.0\times 10^{-3}}{0.090} = 0.10\ \mathrm{N}. \]
Both spheres now carry charges of the \textbf{same sign}, so the force is \textbf{repulsive}, of magnitude $\SI{0.10}{N}$.

\textbf{Remark:} the attractive force before contact ($0.80\ \mathrm{N}$) is much larger than the repulsive force after contact ($0.10\ \mathrm{N}$), because contact redistributed the charge and reduced the product $|q_1 q_2|$ from $8.0\times 10^{-12}$ to $1.0\times 10^{-12}\ \mathrm{C^2}$.
\end{exoresolu}

\courssub{Using the Inverse-Square Dependence and the Effect of the Medium}

\begin{methode}[Proportionality Problems on Coulomb's Law]
\begin{enumerate}
\item \textbf{Write the ratio} rather than recomputing everything: for the same charges,
\[ \frac{F_2}{F_1} = \frac{r_1^2}{r_2^2} \qquad (\text{same medium}). \]
\item \textbf{Changing the medium}: in a medium of relative permittivity $\varepsilon_r$,
\[ F_{\text{medium}} = \frac{F_{\text{vacuum}}}{\varepsilon_r}. \]
\item \textbf{Combined changes}: if both distance and medium change,
\[ \frac{F_2}{F_1} = \frac{\varepsilon_{r,1}}{\varepsilon_{r,2}}\times\frac{r_1^2}{r_2^2}. \]
\item \textbf{Finding an unknown distance}: rearrange to $r_2 = r_1\sqrt{F_1/F_2}$ (same charges, same medium).
\end{enumerate}
\textbf{Reflex:} the force \emph{decreases} when the charges are immersed in an insulating liquid (e.g. kerosene, $\varepsilon_r \approx 2$) because the medium is polarised and partially screens the interaction. The force never depends on the medium for \emph{conductors} --- charges simply redistribute.
\end{methode}

\begin{exoresolu}[Force in Air and in Kerosene]
Two point charges $q_1 = +5.0\ \mathrm{nC}$ and $q_2 = +8.0\ \mathrm{nC}$ are held $10\ \mathrm{cm}$ apart in air ($\varepsilon_r \approx 1$).
\begin{enumerate}
\item Calculate the force between them in air.
\item The charges are now immersed in kerosene of relative permittivity $\varepsilon_r = 2.0$, keeping the same separation. Calculate the new force.
\item At what separation in air would the charges experience the same force as in kerosene at $10\ \mathrm{cm}$?
\end{enumerate}
\tcblower
\textbf{1. Force in air.}
\[ F_{\text{air}} = k\,\frac{q_1 q_2}{r^2} = 9.0\times 10^{9}\times \frac{(5.0\times 10^{-9})(8.0\times 10^{-9})}{(0.10)^2}. \]
The product of the charges is $4.0\times 10^{-17}\ \mathrm{C^2}$, so
\[ F_{\text{air}} = \frac{9.0\times 10^{9}\times 4.0\times 10^{-17}}{1.0\times 10^{-2}} = \frac{3.6\times 10^{-7}}{1.0\times 10^{-2}} = 3.6\times 10^{-5}\ \mathrm{N}. \]
The force is repulsive (charges of the same sign).

\textbf{2. Force in kerosene.} The medium divides the force by $\varepsilon_r$:
\[ F_{\text{ker}} = \frac{F_{\text{air}}}{\varepsilon_r} = \frac{3.6\times 10^{-5}}{2.0} = 1.8\times 10^{-5}\ \mathrm{N}. \]

\textbf{3. Equivalent separation in air.} We require $F_{\text{air}}(r') = F_{\text{ker}}$, with the same charges:
\[ k\,\frac{q_1 q_2}{r'^2} = 1.8\times 10^{-5} \quad\Longrightarrow\quad r'^2 = \frac{9.0\times 10^{9}\times 4.0\times 10^{-17}}{1.8\times 10^{-5}} = 2.0\times 10^{-2}\ \mathrm{m^2}. \]
\[ r' = \sqrt{2.0\times 10^{-2}} = 0.141\ \mathrm{m} \approx 14\ \mathrm{cm}. \]
\textbf{Check by ratio:} $r' = r\sqrt{F_{\text{air}}/F_{\text{ker}}} = 10\times\sqrt{2} = 14.1\ \mathrm{cm}$. Consistent. Physically: to weaken the force by a factor of 2 without the dielectric, the distance must be increased by a factor $\sqrt{2}$.
\end{exoresolu}

\courssub{Superposition of Forces on a Charge}

\begin{methode}[Resultant Force on a Charge Due to Several Others]
\begin{enumerate}
\item \textbf{Draw a clear diagram}: mark all charges with their signs and positions; choose axes.
\item \textbf{Compute each force separately} with Coulomb's law (magnitudes only, using $|q_i q_j|$).
\item \textbf{Draw each force vector} on the charge of interest: attraction towards the source charge, repulsion away from it.
\item \textbf{Add the vectors}: if they are collinear, use signed addition along the line; if not, resolve into components or use the parallelogram rule / Pythagoras.
\item \textbf{State the resultant}: magnitude and direction.
\end{enumerate}
\textbf{Reflex:} the \cle{principle of superposition} says each pair interacts as if the other charges were absent. Never add magnitudes blindly when the forces are not in the same direction.
\end{methode}

\begin{exoresolu}[Three Charges on a Line]
Three point charges are fixed on the $x$-axis: $q_A = +2.0\ \mu\mathrm{C}$ at $x = 0$, $q_B = -4.0\ \mu\mathrm{C}$ at $x = 0.20\ \mathrm{m}$ and $q_C = +3.0\ \mu\mathrm{C}$ at $x = 0.50\ \mathrm{m}$. Calculate the resultant force on $q_B$. Take $k = 9.0\times 10^{9}\ \mathrm{N\,m^2\,C^{-2}}$.
\tcblower
\textbf{Step 1: distances.} $r_{AB} = 0.20\ \mathrm{m}$ and $r_{CB} = 0.50 - 0.20 = 0.30\ \mathrm{m}$.

\textbf{Step 2: force exerted by $q_A$ on $q_B$.} The charges have opposite signs: attraction, so $\vec F_{A/B}$ points from $B$ towards $A$, i.e. in the \textbf{negative} $x$-direction.
\[ F_{A/B} = k\,\frac{|q_A q_B|}{r_{AB}^2} = 9.0\times 10^{9}\times\frac{(2.0\times 10^{-6})(4.0\times 10^{-6})}{(0.20)^2} = \frac{7.2\times 10^{-2}}{0.040} = 1.8\ \mathrm{N}. \]

\textbf{Step 3: force exerted by $q_C$ on $q_B$.} Opposite signs again: attraction, so $\vec F_{C/B}$ points from $B$ towards $C$, i.e. in the \textbf{positive} $x$-direction.
\[ F_{C/B} = k\,\frac{|q_C q_B|}{r_{CB}^2} = 9.0\times 10^{9}\times\frac{(3.0\times 10^{-6})(4.0\times 10^{-6})}{(0.30)^2} = \frac{1.08\times 10^{-1}}{0.090} = 1.2\ \mathrm{N}. \]

\textbf{Step 4: resultant.} The two forces are collinear and opposite in direction. Taking the positive $x$-sense:
\[ F_B = F_{C/B} - F_{A/B} = 1.2 - 1.8 = -0.60\ \mathrm{N}. \]
\textbf{Conclusion:} the resultant force on $q_B$ has magnitude $\SI{0.60}{N}$ and points in the \textbf{negative $x$-direction} (towards $q_A$), because the attraction by $q_A$ (nearer) outweighs the attraction by $q_C$.

\textbf{Check:} although $|q_C| > |q_A|$, $q_C$ is farther away; the $1/r^2$ factor dominates here: $1.8 > 1.2$. This illustrates the inverse-square nature of Coulomb's law.
\end{exoresolu}

\begin{popfigure}
\begin{tikzpicture}[scale=1.2]
% Charges on x-axis
\draw[->, thick] (-0.5,0) -- (3.5,0) node[right] {$x$};
\foreach \x/\lbl in {0/A, 1/B, 2.5/C} {
  \draw (\x,-0.1) -- (\x,0.1);
  \node[below] at (\x,-0.2) {$\lbl$};
}
% Charge A
\node[popred, circle, fill=popred, inner sep=3pt] at (0,0) {};
\node[above] at (0,0.3) {$+2.0\ \mu\mathrm{C}$};
% Charge B
\node[popblue, circle, fill=popblue, inner sep=3pt] at (1,0) {};
\node[above] at (1,0.3) {$-4.0\ \mu\mathrm{C}$};
% Charge C
\node[popred, circle, fill=popred, inner sep=3pt] at (2.5,0) {};
\node[above] at (2.5,0.3) {$+3.0\ \mu\mathrm{C}$};
% Forces on B
\draw[popgreen, thick, ->] (1,0.5) -- (0.2,0.5) node[midway, above] {$F_{A/B}=1.8\ \mathrm{N}$};
\draw[popgreen, thick, ->] (1,0.5) -- (1.8,0.5) node[midway, above] {$F_{C/B}=1.2\ \mathrm{N}$};
\node at (1,-0.8) {Resultant on $q_B$: $0.60\ \mathrm{N}$ towards $A$};
\end{tikzpicture}
\legende{Forces on $q_B$ due to $q_A$ and $q_C$. Both are attractive; the nearer charge $q_A$ wins.}
\end{popfigure}

\courssub{Electric Field of a Point Charge and Field Superposition}

\begin{definition}[Electric Field Strength]
The \cle{electric field strength} $\vec{E}$ at a point is the force per unit positive test charge placed at that point:
\[ \vec{E} = \frac{\vec{F}}{q_0} \quad (q_0 > 0,\ \text{infinitesimal}). \]
Units: $\mathrm{N\,C^{-1}}$ or $\mathrm{V\,m^{-1}}$.
\end{definition}

\begin{propriete}[Field of a Point Charge]
The field due to a source charge $Q$ at distance $r$ is
\[ E = k\,\frac{|Q|}{r^2} \qquad (\mathrm{N\,C^{-1}}). \]
Direction: $\vec{E}$ points \emph{away} from $Q$ if $Q > 0$, \emph{towards} $Q$ if $Q < 0$. A positive test charge would move along $\vec{E}$; the force on a charge $q$ is $\vec{F} = q\vec{E}$.
\end{propriete}

\begin{propriete}[Superposition of Electric Fields]
The total field at a point is the \emph{vector} sum of the fields due to each source charge:
\[ \vec{E} = \vec{E}_1 + \vec{E}_2 + \cdots \]
Field lines never cross; they start on positive charges and end on negative charges; their density shows the field strength.
\end{propriete}

\begin{methode}[Computing and Combining Electric Fields]
\begin{enumerate}
\item \textbf{Definition}: the field due to a source charge $Q$ at distance $r$ is
\[ E = k\,\frac{|Q|}{r^2} \qquad (\mathrm{N\,C^{-1}} \text{ or } \mathrm{V\,m^{-1}}). \]
\item \textbf{Direction}: $\vec{E}$ points \emph{away} from $Q$ if $Q > 0$, \emph{towards} $Q$ if $Q < 0$. A positive test charge would move along $\vec{E}$; the force on a charge $q$ is $\vec{F} = q\vec{E}$.
\item \textbf{Superposition}: the total field at a point is the \emph{vector} sum $\vec{E} = \vec{E}_1 + \vec{E}_2 + \cdots$
\item \textbf{Zero-field points}: for two charges, look for a point where the two field vectors are \emph{collinear and opposite} with equal magnitudes. Between two charges of the \emph{same} sign; outside the segment (on the side of the weaker charge) for charges of \emph{opposite} sign.
\end{enumerate}
\textbf{Reflex:} field lines never cross; they start on positive charges and end on negative charges; their density shows the field strength.
\end{methode}

\begin{exoresolu}[Field at the Midpoint and Zero-Field Point]
Two point charges $q_1 = +6.0\ \mathrm{nC}$ and $q_2 = -6.0\ \mathrm{nC}$ are placed at $A$ and $B$, with $AB = 40\ \mathrm{cm}$ in air.
\begin{enumerate}
\item Calculate the electric field (magnitude and direction) at the midpoint $M$ of $AB$.
\item A charge $q_3 = +2.0\ \mathrm{nC}$ is placed at $M$. Calculate the force on it.
\end{enumerate}
\tcblower
\textbf{1. Field at $M$.} $M$ is at $r = 0.20\ \mathrm{m}$ from each charge. Magnitude of each contribution:
\[ E_1 = E_2 = k\,\frac{|q|}{r^2} = 9.0\times 10^{9}\times\frac{6.0\times 10^{-9}}{(0.20)^2} = \frac{54}{0.040} = 1.35\times 10^{3}\ \mathrm{N\,C^{-1}}. \]
Directions: $q_1$ is positive, so $\vec E_1$ at $M$ points \emph{away} from $A$, i.e. from $A$ towards $B$. $q_2$ is negative, so $\vec E_2$ at $M$ points \emph{towards} $B$, i.e. also from $A$ towards $B$. The two fields \textbf{add}:
\[ E_M = E_1 + E_2 = 2\times 1.35\times 10^{3} = 2.7\times 10^{3}\ \mathrm{N\,C^{-1}}, \]
directed from $A$ to $B$ (from the positive charge towards the negative charge), as expected for the field between the charges of a dipole.

\textbf{2. Force on $q_3$ at $M$.}
\[ F = q_3 E_M = (2.0\times 10^{-9})(2.7\times 10^{3}) = 5.4\times 10^{-6}\ \mathrm{N}. \]
Since $q_3 > 0$, $\vec F$ is parallel to $\vec E_M$: the force points from $A$ towards $B$.

\textbf{Remark:} one could equally compute the two Coulomb forces on $q_3$ directly: $F_1 = k|q_1 q_3|/r^2 = 2.7\ \mu\mathrm{N}$ (repulsion, towards $B$) and $F_2 = 2.7\ \mu\mathrm{N}$ (attraction, towards $B$), giving the same total $5.4\ \mu\mathrm{N}$. The field method is faster once $\vec E$ is known.
\end{exoresolu}

\begin{popfigure}
\begin{tikzpicture}[scale=1.0]
% Dipole field lines
\foreach \angle in {0,30,60,90,120,150,180,210,240,270,300,330} {
  \draw[popblue, ->] (0,0) -- (\angle:2);
}
\foreach \angle in {0,30,60,90,120,150,180,210,240,270,300,330} {
  \draw[popred, ->] (3,0) -- ++(\angle:2);
}
\node[popblue, circle, fill=popblue, inner sep=3pt] at (0,0) {};
\node[popred, circle, fill=popred, inner sep=3pt] at (3,0) {};
\node at (0,-0.5) {$+Q$};
\node at (3,-0.5) {$-Q$};
\node at (1.5,2.2) {Field lines: $+Q \to -Q$};
% Field at midpoint
\draw[popgreen, thick, ->] (1.5,0) -- (1.5,1) node[above] {$\vec{E}_M$};
\node at (1.5,-0.5) {$M$};
\end{tikzpicture}
\legende{Electric field lines for an electric dipole ($+Q$ and $-Q$). At the midpoint $M$, the field points from $+Q$ to $-Q$.}
\end{popfigure}

\courssub{Locating a Point of Zero Field}

\begin{methode}[Finding Where the Resultant Field Vanishes]
\begin{enumerate}
\item \textbf{Decide the region} where cancellation is possible (fields must be in opposite directions there):
\begin{itemize}
\item same-sign charges: \emph{between} the charges;
\item opposite-sign charges: on the line, \emph{outside} the segment, on the side of the \emph{smaller} magnitude.
\end{itemize}
\item \textbf{Set the magnitudes equal}: $\dfrac{|q_1|}{x^2} = \dfrac{|q_2|}{(d \pm x)^2}$.
\item \textbf{Take the square root} to avoid solving a full quadratic: $\dfrac{d \pm x}{x} = \sqrt{|q_2|/|q_1|}$.
\item \textbf{Solve for $x$}, check that the solution lies in the expected region, and reject spurious roots.
\end{enumerate}
\end{methode}

\begin{exoresolu}[Zero Field Near Two Opposite Charges]
A charge $q_1 = +9.0\ \mathrm{nC}$ is placed at the origin and $q_2 = -4.0\ \mathrm{nC}$ at $x = 0.30\ \mathrm{m}$ on the $x$-axis. Find the point(s) on the axis where the resultant electric field is zero.
\tcblower
\textbf{Step 1: choose the region.} The charges have \emph{opposite} signs, so the fields point in the \emph{same} direction between them (both point from $q_1$ towards $q_2$ there): no cancellation in $0 < x < 0.30\ \mathrm{m}$. Outside the segment, cancellation can only occur on the side of the \emph{weaker} charge, i.e. beyond $q_2$, at $x > 0.30\ \mathrm{m}$. (On the left of $q_1$, the stronger charge is also the nearer one, so its field always wins.)

\textbf{Step 2: equality of magnitudes.} Let $P$ be at position $x > 0.30\ \mathrm{m}$. Its distances are $x$ from $q_1$ and $(x - 0.30)$ from $q_2$:
\[ \frac{|q_1|}{x^2} = \frac{|q_2|}{(x-0.30)^2} \quad\Longrightarrow\quad \frac{9.0}{x^2} = \frac{4.0}{(x-0.30)^2}. \]

\textbf{Step 3: take square roots.}
\[ \frac{x}{x - 0.30} = \sqrt{\frac{9.0}{4.0}} = \frac{3}{2} = 1.5. \]

\textbf{Step 4: solve.}
\[ x = 1.5\,(x - 0.30) \;\Longrightarrow\; x = 1.5x - 0.45 \;\Longrightarrow\; 0.5x = 0.45 \;\Longrightarrow\; x = 0.90\ \mathrm{m}. \]
Check: $x = 0.90 > 0.30\ \mathrm{m}$: the point lies in the expected region. Verification of magnitudes:
\[ E_1 = \frac{k\times 9.0\times 10^{-9}}{0.90^2} = 11.1\ k\times 10^{-9}, \qquad E_2 = \frac{k\times 4.0\times 10^{-9}}{0.60^2} = 11.1\ k\times 10^{-9}. \]
Equal magnitudes, opposite directions ($\vec E_1$ points away from $q_1$, i.e. $+x$; $\vec E_2$ points towards $q_2$, i.e. $-x$). \textbf{The field vanishes at $x = 0.90\ \mathrm{m}$}, i.e. $0.60\ \mathrm{m}$ beyond $q_2$.
\end{exoresolu}

\begin{popfigure}
\begin{tikzpicture}[scale=1.2]
\draw[->, thick] (-1,0) -- (4,0) node[right] {$x$};
\foreach \x/\lbl in {0/q_1, 1/q_2, 3/P} {
  \draw (\x,-0.1) -- (\x,0.1);
  \node[below] at (\x,-0.3) {$\lbl$};
}
\node[popred, circle, fill=popred, inner sep=3pt] at (0,0) {};
\node[popblue, circle, fill=popblue, inner sep=3pt] at (1,0) {};
\node[popgreen, circle, fill=popgreen, inner sep=3pt] at (3,0) {};
\node at (0,0.5) {$+9\ \mathrm{nC}$};
\node at (1,0.5) {$-4\ \mathrm{nC}$};
\node at (3,0.5) {$P$ (zero field)};
% Field vectors at P
\draw[popred, thick, ->] (3,0.3) -- (2.2,0.3) node[midway, above] {$\vec{E}_1$};
\draw[popblue, thick, ->] (3,0.3) -- (3.8,0.3) node[midway, above] {$\vec{E}_2$};
\node at (1.5,-1) {$E_1 = E_2$ at $x=0.90\ \mathrm{m}$};
\end{tikzpicture}
\legende{Zero-field point for opposite charges: outside the segment, on the side of the smaller charge.}
\end{popfigure}

\courssub{Equilibrium of a Charged Body in an Electric Field}

\begin{methode}[Charged Body in Equilibrium in an Electric Field]
\begin{enumerate}
\item \textbf{List all forces} on the body: weight $\vec P = m\vec g$ (vertical, downwards), electric force $\vec F = q\vec E$ (parallel to $\vec E$ if $q>0$, antiparallel if $q<0$), tension $\vec T$ along the string, normal reaction, friction, etc.
\item \textbf{Draw the free-body diagram} and choose axes (often horizontal/vertical).
\item \textbf{Write the equilibrium conditions}: $\sum F_x = 0$ and $\sum F_y = 0$.
\item \textbf{Solve} for the unknown (charge, field, angle). For a pendulum deflected by a horizontal field: $\tan\theta = \dfrac{qE}{mg}$.
\item \textbf{Check} the sign of $q$ against the direction of deflection.
\end{enumerate}
\textbf{Reflex:} the electric force on a charge in a field is \emph{independent} of the other forces; treat it exactly like weight in a force balance.
\end{methode}

\begin{exoresolu}[Charged Pendulum in a Uniform Field]
A small sphere of mass $m = 2.0\ \mathrm{g}$ carrying a charge $q$ hangs from an insulating thread. A uniform horizontal electric field $E = 3.0\times 10^{4}\ \mathrm{N\,C^{-1}}$ is applied, and the thread settles at an angle $\theta = 25^{\circ}$ to the vertical, the sphere being deflected in the direction of the field. Take $g = 9.81\ \mathrm{m\,s^{-2}}$, $\tan 25^{\circ} = 0.466$.
\begin{enumerate}
\item State the sign of $q$, justifying your answer.
\item Calculate the value of $q$ and the tension in the thread.
\end{enumerate}
\tcblower
\textbf{1. Sign of $q$.} The sphere is deflected \emph{in the direction of $\vec E$}, so the electric force $\vec F = q\vec E$ is parallel to $\vec E$; this requires $q > 0$. The charge is \textbf{positive}.

\textbf{2. Force balance.} Three forces act on the sphere: the weight $P = mg$ (vertical), the electric force $F = qE$ (horizontal), and the tension $T$ along the thread. At equilibrium, resolving horizontally and vertically:
\[ T\sin\theta = qE \qquad\text{and}\qquad T\cos\theta = mg. \]
Dividing the two equations:
\[ \tan\theta = \frac{qE}{mg} \quad\Longrightarrow\quad q = \frac{mg\tan\theta}{E}. \]
Numerically, with $m = 2.0\times 10^{-3}\ \mathrm{kg}$:
\[ q = \frac{(2.0\times 10^{-3})(9.81)(0.466)}{3.0\times 10^{4}} = \frac{9.14\times 10^{-3}}{3.0\times 10^{4}} = 3.0\times 10^{-7}\ \mathrm{C} = 0.30\ \mu\mathrm{C}. \]
Tension, from $T = mg/\cos\theta$ with $\cos 25^{\circ} = 0.906$:
\[ T = \frac{(2.0\times 10^{-3})(9.81)}{0.906} = \frac{1.96\times 10^{-2}}{0.906} = 2.2\times 10^{-2}\ \mathrm{N}. \]
\textbf{Check:} the horizontal component $T\sin\theta = 2.2\times 10^{-2}\times 0.423 = 9.2\times 10^{-3}\ \mathrm{N}$ equals $qE = (3.0\times 10^{-7})(3.0\times 10^{4}) = 9.0\times 10^{-3}\ \mathrm{N}$. Consistent (small rounding).
\end{exoresolu}

\begin{popfigure}
\begin{tikzpicture}[scale=1.2]
% Pendulum in field
\draw[popdark, thick] (0,3) -- (0,0); % ceiling
\draw[popblue, thick] (0,3) -- (1.27,1.87); % thread at 25 deg
\node[popred, circle, fill=popred, inner sep=3pt] at (1.27,1.87) {}; % sphere
% Forces
\draw[popgreen, thick, ->] (1.27,1.87) -- (1.27,0.87) node[midway, right] {$mg$};
\draw[poporange, thick, ->] (1.27,1.87) -- (2.27,1.87) node[midway, above] {$qE$};
\draw[poppurple, thick, ->] (1.27,1.87) -- (0.5,2.5) node[midway, left] {$T$};
% Angle
\draw[popmuted, ->] (0,2.5) arc (90:65:0.5) node[midway, left] {$\theta$};
% Field direction
\draw[popblue, dashed, ->] (2.5,1.87) -- (3.5,1.87) node[right] {$\vec{E}$};
\node at (1.27,1.2) {$m=2.0\ \mathrm{g}$};
\node at (1.27,2.3) {$q=+0.30\ \mu\mathrm{C}$};
\end{tikzpicture}
\legende{Free-body diagram of a charged pendulum in a uniform horizontal electric field.}
\end{popfigure}

\begin{attention}[Frequent Errors in This Chapter]
\begin{itemize}
\item Forgetting to convert $\mu\mathrm{C}$, $\mathrm{nC}$ or $\mathrm{cm}$ to SI units before applying Coulomb's law.
\item Squaring only the numerical value of $r$ but not its power of ten: $(0.30)^2 = 0.090$, not $0.9$.
\item Adding force or field \emph{magnitudes} when the vectors are not parallel: always add vectors.
\item Believing that a heavier or larger charge feels a larger force in a pair: Newton's third law gives equal magnitudes.
\item Claiming two bodies attract because they carry opposite charges: attraction can also occur by induction on a neutral body; only repulsion is conclusive.
\item Forgetting that in a dielectric medium the force is \emph{divided} by $\varepsilon_r$, not multiplied.
\item Saying that "positive charges move" when a body is charged: in solids, only electrons are transferred.
\item Confusing the direction of the electric field: it points \emph{away} from positive, \emph{towards} negative.
\item Forgetting to take the square root when solving for zero-field points, leading to unnecessary quadratic equations.
\end{itemize}
\end{attention}

\begin{savaistu}[Did You Know?]
\begin{itemize}
\item The \textbf{triboelectric series} ranks materials by their tendency to become positive or negative when rubbed. For example, rubbing glass with silk makes glass positive; rubbing ebonite with wool makes ebonite negative.
\item \textbf{Lightning} is a giant spark: a typical lightning bolt carries $\sim 15\ \mathrm{C}$ of charge, reaches $30\,000\ \mathrm{K}$ (hotter than the Sun's surface), and releases $\sim 10^9\ \mathrm{J}$ of energy.
\item The \textbf{Faraday cage} effect: a closed conducting shell shields its interior from external electric fields. This is why you are safe inside a car during a thunderstorm.
\item \textbf{Photocopiers and laser printers} use electrostatics: a charged drum attracts toner particles only where light has discharged it.
\item \textbf{Electrostatic precipitators} in cement factories (like CIMENCAM in Figuil) remove dust from exhaust gases by charging particles and collecting them on oppositely charged plates.
\end{itemize}
\end{savaistu}

\begin{exercice}[Practice Problems]
\begin{enumerate}
\item A glass rod rubbed with silk acquires a charge of $+8.0\ \mathrm{nC}$. How many electrons were transferred? Is the silk positive or negative?
\item Two identical conducting spheres carry charges $+6.0\ \mu\mathrm{C}$ and $-2.0\ \mu\mathrm{C}$. They are touched together and separated to $20\ \mathrm{cm}$. Find the force between them.
\item Three charges $+q$, $+q$, and $-q$ are placed at the vertices of an equilateral triangle of side $a$. Find the magnitude and direction of the resultant force on the negative charge.
\item A uniform electric field $E = 2.0\times 10^4\ \mathrm{N\,C^{-1}}$ points horizontally. A $1.0\ \mathrm{g}$ ball with charge $-5.0\ \mu\mathrm{C}$ is suspended by a thread. Find the angle the thread makes with the vertical.
\item Two point charges $+4.0\ \mu\mathrm{C}$ and $+9.0\ \mu\mathrm{C}$ are $30\ \mathrm{cm}$ apart. Where on the line joining them is the electric field zero?
\end{enumerate}
\end{exercice}

\begin{corrige}[Exercise 1]
$n = \frac{8.0\times 10^{-9}}{1.60\times 10^{-19}} = 5.0\times 10^{10}$ electrons transferred from glass to silk. Silk becomes negative ($-8.0\ \mathrm{nC}$).
\end{corrige}

\begin{corrige}[Exercise 2]
Total charge $= +4.0\ \mu\mathrm{C}$. Each sphere gets $+2.0\ \mu\mathrm{C}$. $F = 9.0\times 10^9 \times \frac{(2.0\times 10^{-6})^2}{(0.20)^2} = 0.90\ \mathrm{N}$ (repulsive).
\end{corrige}

\begin{corrige}[Exercise 3]
Forces from each $+q$ on $-q$ are attractive, magnitude $F = k q^2/a^2$, at $60^\circ$ to each other. Resultant $F_R = \sqrt{F^2+F^2+2F^2\cos 60^\circ} = \sqrt{3}F = \sqrt{3}k q^2/a^2$, directed along the line bisecting the angle (towards the midpoint of the two positive charges).
\end{corrige}

\begin{corrige}[Exercise 4]
$q$ negative $\Rightarrow$ force opposite to $\vec{E}$. $\tan\theta = \frac{|q|E}{mg} = \frac{(5.0\times 10^{-6})(2.0\times 10^4)}{(1.0\times 10^{-3})(9.81)} = 10.2$. $\theta \approx 84.4^\circ$ from vertical (almost horizontal).
\end{corrige}

\begin{corrige}[Exercise 5]
Same sign $\Rightarrow$ zero field between them. Let $x$ be distance from $+4\ \mu\mathrm{C}$. $\frac{4}{x^2} = \frac{9}{(0.30-x)^2} \Rightarrow \frac{0.30-x}{x} = \frac{3}{2} \Rightarrow 0.30-x = 1.5x \Rightarrow 2.5x = 0.30 \Rightarrow x = 0.12\ \mathrm{m} = 12\ \mathrm{cm}$ from the $+4\ \mu\mathrm{C}$ charge.
\end{corrige}

\newpage

\coursec{Exercises}

\courssub{I check my knowledge}

\begin{exercice}[Multiple choice questions]
For each question, choose the single correct answer and justify your choice in one sentence.
\begin{enumerate}
\item Two identical conducting spheres, one carrying $+8\ \mu\mathrm{C}$ and the other $-2\ \mu\mathrm{C}$, are brought into contact and then separated. The charge on each sphere after separation is:
\begin{enumerate}
\item $+5\ \mu\mathrm{C}$
\item $+3\ \mu\mathrm{C}$
\item $-3\ \mu\mathrm{C}$
\item $+10\ \mu\mathrm{C}$
\end{enumerate}
\item The force between two point charges in air is $F$. If both charges are doubled and the separation is also doubled, the new force is:
\begin{enumerate}
\item $F/4$
\item $F/2$
\item $F$
\item $4F$
\end{enumerate}
\item The SI unit of electric field strength is:
\begin{enumerate}
\item $\mathrm{N}$
\item $\mathrm{N\,C^{-1}}$
\item $\mathrm{C\,N^{-1}}$
\item $\mathrm{J\,C^{-1}}$
\end{enumerate}
\item A lightning conductor protects a building mainly because of:
\begin{enumerate}
\item the large mass of the copper rod
\item point action at its sharp tip
\item the insulation of the rod
\item the attraction of clouds by copper
\end{enumerate}
\end{enumerate}
\end{exercice}

\begin{exercice}[True or false]
State whether each statement is true or false, and justify your answer briefly.
\begin{enumerate}
\item A negatively charged body has gained electrons.
\item When a glass rod is rubbed with silk, the silk acquires a positive charge.
\item The force between two point charges increases when they are placed in a medium of relative permittivity greater than 1.
\item The electric field due to a point charge varies as $1/r^{2}$, just like the force between two point charges.
\item Charge resides on the outer surface of a hollow conductor in electrostatic equilibrium.
\end{enumerate}
\end{exercice}

\begin{exercice}[Definitions]
Define each of the following terms precisely, as you would in the GCE examination:
\begin{enumerate}
\item electric field strength at a point;
\item a point charge;
\item relative permittivity of a medium;
\item point action (action of points).
\end{enumerate}
\end{exercice}

\begin{exercice}[Direct calculations]
\begin{enumerate}
\item Calculate the charge carried by $2.5\times 10^{13}$ electrons. ($e = 1.60\times 10^{-19}\ \mathrm{C}$)
\item Two point charges of $+3.0\ \mu\mathrm{C}$ each are placed $10\ \mathrm{cm}$ apart in vacuum. Calculate the magnitude of the force on each charge. ($k = \dfrac{1}{4\pi\varepsilon_{0}} = 9.0\times 10^{9}\ \mathrm{N\,m^{2}\,C^{-2}}$)
\item Calculate the electric field strength at a point $25\ \mathrm{cm}$ from a charge of $+2.0\ \mathrm{nC}$ in air.
\end{enumerate}
\end{exercice}

\courssub{I practise}

\begin{exercice}[Force between two charges]
Two point charges $q_{1} = +4.0\ \mu\mathrm{C}$ and $q_{2} = -6.0\ \mu\mathrm{C}$ are placed $30\ \mathrm{cm}$ apart in air.
\begin{enumerate}
\item Calculate the magnitude of the force between them and state its nature (attractive or repulsive).
\item Find the new force if the separation is doubled.
\item The system is now immersed in a medium of relative permittivity $\varepsilon_{r} = 4$, with the original separation of $30\ \mathrm{cm}$. Calculate the new force and comment on the effect of the medium.
\end{enumerate}
\end{exercice}

\begin{exercice}[Resultant force on a charge]
Three point charges $q_{A} = +2.0\ \mu\mathrm{C}$, $q_{B} = +2.0\ \mu\mathrm{C}$ and $q_{C} = -4.0\ \mu\mathrm{C}$ are placed at the corners $A$, $B$ and $C$ of an equilateral triangle of side $20\ \mathrm{cm}$ in air.
\begin{enumerate}
\item Draw a clear diagram showing the forces exerted on the charge at $C$.
\item Calculate the magnitude of each force acting on the charge at $C$.
\item Using vector components, determine the magnitude and direction of the resultant force on the charge at $C$.
\end{enumerate}
\end{exercice}

\begin{exercice}[Charged sphere hanging from a thread]
A small conducting sphere of mass $0.5\ \mathrm{g}$, carrying an unknown charge $q$, hangs from an insulating silk thread near a fixed charge of $+2.0\ \mu\mathrm{C}$. At equilibrium the thread makes an angle of $30^{\circ}$ with the vertical, and the centre of the sphere is $10\ \mathrm{cm}$ from the fixed charge, at the same horizontal level. Take $g = 9.81\ \mathrm{m\,s^{-2}}$.
\begin{enumerate}
\item Draw the free-body diagram of the sphere.
\item By resolving forces horizontally and vertically, determine the electrostatic force on the sphere.
\item Deduce the magnitude and sign of $q$.
\end{enumerate}
\end{exercice}

\begin{exercice}[Electric field of two charges]
A charge $q_{1} = +5.0\ \mathrm{nC}$ and a charge $q_{2} = -5.0\ \mathrm{nC}$ are fixed $1.0\ \mathrm{m}$ apart in air.
\begin{enumerate}
\item Calculate the electric field strength at a point $50\ \mathrm{cm}$ from $q_{1}$ on the side opposite to $q_{2}$.
\item Investigate whether there exists a point on the line joining the two charges where the resultant electric field is zero. Set up the equation, discuss its solutions, and conclude.
\item Repeat the discussion for the case where both charges are $+5.0\ \mathrm{nC}$, and find the position of the zero-field point if it exists.
\end{enumerate}
\end{exercice}

\courssub{I prepare for the GCE}

\begin{exercice}[Charging by induction — structured question (12 marks)]
A gold-leaf electroscope is initially uncharged.
\begin{enumerate}
\item Describe, with the aid of clearly labelled diagrams, how the electroscope can be given a \emph{positive} charge by induction using a negatively charged ebonite rod. \hfill (6 marks)
\item Explain each step of the process in terms of the movement of electrons within the electroscope. \hfill (4 marks)
\item State what happens to the leaf if a negatively charged body is afterwards brought near the cap of the positively charged electroscope, and explain why. \hfill (2 marks)
\end{enumerate}
\end{exercice}

\begin{exercice}[Thunderstorm safety — essay question (10 marks)]
During the rainy season in Cameroon, thunderstorms are frequent, and people are often advised not to shelter under trees.
\begin{enumerate}
\item Explain, using the ideas of conductors, insulators and charge distribution, why a person is safer inside a closed car than under a tree during a thunderstorm. \hfill (5 marks)
\item Describe the construction of a lightning conductor and explain, with reference to point action, how it protects a building. \hfill (5 marks)
\end{enumerate}
\end{exercice}

\begin{exercice}[Data analysis — verification of Coulomb's law (18 marks)]
In an experiment, two small charged spheres are held at different separations $r$, and the force $F$ between them is measured. The results are given below.

\begin{center}
\begin{tabularx}{\linewidth}{|l|X|X|X|X|X|X|}
\hline
\rowcolor{popblueL} $r$ / cm & 5.0 & 8.0 & 10.0 & 15.0 & 20.0 & 25.0 \\
\hline
$F$ / $10^{-3}\ \mathrm{N}$ & 14.4 & 5.6 & 3.6 & 1.6 & 0.90 & 0.58 \\
\hline
$1/r^{2}$ / $\mathrm{m^{-2}}$ & & & & & & \\
\hline
\end{tabularx}
\end{center}

\begin{enumerate}
\item State Coulomb's law and explain why a graph of $F$ against $1/r^{2}$ should be a straight line through the origin. \hfill (4 marks)
\item Copy and complete the table by calculating the values of $1/r^{2}$ (in $\mathrm{m^{-2}}$). \hfill (3 marks)
\item Plot the graph of $F$ (vertical axis) against $1/r^{2}$ (horizontal axis), choosing suitable scales, and draw the best straight line. \hfill (5 marks)
\item Determine the gradient of the line and deduce the product $q_{1}q_{2}$ of the two charges. Take $k = 9.0\times 10^{9}\ \mathrm{N\,m^{2}\,C^{-2}}$. \hfill (4 marks)
\item State two precautions that should be taken in this experiment to obtain reliable results. \hfill (2 marks)
\end{enumerate}
\end{exercice}

\newpage

\coursec{Solutions to the exercises}

\begin{corrige}[Exercise 1 — Multiple choice questions]
\begin{enumerate}
\item \textbf{Answer: (b) $+3\,\mu\mathrm{C}$.} When identical conducting spheres touch, they form a single conductor and the total charge is shared equally between them. Total charge $Q_{\text{tot}} = +8\,\mu\mathrm{C} + (-2\,\mu\mathrm{C}) = +6\,\mu\mathrm{C}$. Charge on each sphere after separation: $q = Q_{\text{tot}}/2 = +3\,\mu\mathrm{C}$.
\item \textbf{Answer: (c) $F$.} By Coulomb's law, $F = k\dfrac{q_{1}q_{2}}{r^{2}}$. If both charges are doubled and the separation is doubled, the new force is
\[
F' = k\dfrac{(2q_{1})(2q_{2})}{(2r)^{2}} = k\dfrac{4q_{1}q_{2}}{4r^{2}} = k\dfrac{q_{1}q_{2}}{r^{2}} = F.
\]
\item \textbf{Answer: (b) $\mathrm{N\,C^{-1}}$.} Electric field strength is defined as $\vec{E} = \vec{F}/q$. The SI unit of force is the newton ($\mathrm{N}$) and that of charge is the coulomb ($\mathrm{C}$), hence the unit of $E$ is $\mathrm{N\,C^{-1}}$ (equivalently $\mathrm{V\,m^{-1}}$).
\item \textbf{Answer: (b) point action at its sharp tip.} On a conductor, surface charge density is greatest at sharply curved points. This produces a very intense local electric field which ionises the surrounding air, allowing charge to leak away gradually (or providing a preferred path for a lightning strike to earth).
\end{enumerate}
\end{corrige}

\begin{attention}[Examiner's expectations for Exercise 1]
In (i), remember that charge conservation and equal sharing apply only to \emph{identical} conductors. In (ii), check the combined effect of changes in charge and distance carefully — do not treat them separately. In (iv), the key phrase is ``point action''; ``sharp point'' alone is insufficient.
\end{attention}

\begin{corrige}[Exercise 2 — True or false]
\begin{enumerate}
\item \textbf{True.} A body becomes negatively charged by gaining electrons (which carry negative charge); a positively charged body has lost electrons.
\item \textbf{False.} Glass rubbed with silk becomes \emph{positive} because it \emph{loses} electrons to the silk; the silk gains those electrons and therefore acquires a \emph{negative} charge.
\item \textbf{False.} In a material medium, Coulomb's law reads $F = \dfrac{1}{4\pi\varepsilon_{0}\varepsilon_{r}}\dfrac{q_{1}q_{2}}{r^{2}}$. Since the relative permittivity $\varepsilon_{r} > 1$ for any material medium, the force is \emph{reduced} by the factor $\varepsilon_{r}$ compared with its value in vacuum.
\item \textbf{True.} For a point charge $Q$, $E = \dfrac{1}{4\pi\varepsilon_{0}}\dfrac{Q}{r^{2}}$. This inverse-square dependence is the same as that of the Coulomb force because $\vec{F} = q\vec{E}$.
\item \textbf{True.} In electrostatic equilibrium the electric field inside the material of a conductor is zero. By Gauss's law, this implies no net charge resides in the interior; all excess charge resides on the outer surface.
\end{enumerate}
\end{corrige}

\begin{attention}[Examiner's expectations for Exercise 2]
For (iii), distinguish between \emph{permittivity} $\varepsilon$ and \emph{relative permittivity} $\varepsilon_{r}$. The force is reduced by $\varepsilon_{r}$, not increased. For (v), the statement applies only in \emph{electrostatic equilibrium}; if a current flows, the field inside is not zero.
\end{attention}

\begin{corrige}[Exercise 3 — Definitions]
\begin{enumerate}
\item \textbf{Electric field strength at a point} is the electrostatic force per unit positive charge that would act on a small test charge placed at that point:
\[
\vec{E} = \frac{\vec{F}}{q},
\]
a vector quantity whose SI unit is the newton per coulomb ($\mathrm{N\,C^{-1}}$), equivalently the volt per metre ($\mathrm{V\,m^{-1}}$).
\item \textbf{A point charge} is a charged body whose physical dimensions are negligible compared with the distances at which its electrical effects are considered, so that all its charge may be treated as concentrated at a single point.
\item \textbf{The relative permittivity $\varepsilon_{r}$ of a medium} is the ratio of the permittivity of the medium to the permittivity of free space, $\varepsilon_{r} = \varepsilon/\varepsilon_{0}$. It equals the factor by which the electrostatic force between two given charges is reduced when vacuum is replaced by the medium. It is a dimensionless number (no unit).
\item \textbf{Point action} is the concentration of electric charge (and hence of electric field) at sharply curved points of a conductor; the intense field near a sharp point ionises the surrounding air, allowing charge to be discharged rapidly through the point.
\end{enumerate}
\end{corrige}

\begin{attention}[Examiner's expectations for Exercise 3]
Definitions must be precise. For (i), include ``per unit \emph{positive} charge'' and ``small test charge''. For (ii), the comparison of dimensions with distance is essential. For (iii), state clearly that $\varepsilon_{r}$ is a ratio and has no unit. For (iv), mention both charge concentration and air ionisation.
\end{attention}

\begin{corrige}[Exercise 4 — Direct calculations]
\begin{enumerate}
\item Charge is quantised: $Q = ne$, where $e = 1.60\times 10^{-19}\,\mathrm{C}$.
\[
Q = (2.5\times 10^{13})\times(1.60\times 10^{-19}\,\mathrm{C}) = 4.0\times 10^{-6}\,\mathrm{C} = -4.0\,\mu\mathrm{C}.
\]
The charge is negative because electrons carry negative charge.
\item Coulomb's law for two equal charges $q = 3.0\times 10^{-6}\,\mathrm{C}$ at distance $r = 0.10\,\mathrm{m}$:
\[
F = k\frac{q^{2}}{r^{2}} = (9.0\times 10^{9}\,\mathrm{N\,m^{2}\,C^{-2}})\frac{(3.0\times 10^{-6}\,\mathrm{C})^{2}}{(0.10\,\mathrm{m})^{2}}
= 9.0\times 10^{9}\times\frac{9.0\times 10^{-12}}{0.010} = 8.1\,\mathrm{N}.
\]
The force on each charge has magnitude $8.1\,\mathrm{N}$; it is repulsive because the charges have the same sign.
\item Electric field of a point charge $Q = 2.0\times 10^{-9}\,\mathrm{C}$ at distance $r = 0.25\,\mathrm{m}$:
\[
E = \frac{kQ}{r^{2}} = \frac{(9.0\times 10^{9})\times(2.0\times 10^{-9})}{(0.25)^{2}} = \frac{18}{0.0625} = 2.88\times 10^{2}\,\mathrm{N\,C^{-1}} \approx 2.9\times 10^{2}\,\mathrm{N\,C^{-1}}.
\]
The direction is radially away from the positive charge.
\end{enumerate}
\end{corrige}

\begin{attention}[Examiner's expectations for Exercise 4]
Always include units in every step. In (i), the sign of the charge must be stated explicitly. In (ii) and (iii), give the final answer to two or three significant figures consistent with the data. Direction of $\vec{E}$ is required for vector quantities.
\end{attention}

\begin{corrige}[Exercise 5 — Force between two charges]
\begin{enumerate}
\item Magnitude of the force between $q_{1} = +4.0\,\mu\mathrm{C}$ and $q_{2} = -6.0\,\mu\mathrm{C}$ at $r = 0.30\,\mathrm{m}$:
\[
F = k\frac{|q_{1}q_{2}|}{r^{2}} = (9.0\times 10^{9})\frac{(4.0\times 10^{-6})(6.0\times 10^{-6})}{(0.30)^{2}}
= 9.0\times 10^{9}\times\frac{2.4\times 10^{-11}}{0.090} = 2.4\,\mathrm{N}.
\]
The charges have opposite signs, so the force is \textbf{attractive}.
\item Doubling the separation to $0.60\,\mathrm{m}$ divides the force by $2^{2} = 4$:
\[
F' = \frac{F}{4} = \frac{2.4}{4} = 0.60\,\mathrm{N}.
\]
\item In a medium of relative permittivity $\varepsilon_{r} = 4$, the force is reduced by the factor $\varepsilon_{r}$:
\[
F'' = \frac{F}{\varepsilon_{r}} = \frac{2.4}{4} = 0.60\,\mathrm{N}.
\]
\textbf{Comment:} The molecules of the medium polarise in the field of the charges, producing an opposing field that partially shields the interaction.
\end{enumerate}
\end{corrige}

\begin{attention}[Examiner's expectations for Exercise 5]
In (i), use the absolute values of the charges for the magnitude and state the nature (attractive/repulsive) separately. In (iii), write the formula $F_{\text{medium}} = F_{\text{vacuum}}/\varepsilon_{r}$ explicitly.
\end{attention}

\begin{corrige}[Exercise 6 — Resultant force on a charge]
\begin{enumerate}
\item The three charges form an equilateral triangle of side $0.20\,\mathrm{m}$. $q_{A} = q_{B} = +2.0\,\mu\mathrm{C}$, $q_{C} = -4.0\,\mu\mathrm{C}$. Since $q_{C}$ is negative, both $q_{A}$ and $q_{B}$ \emph{attract} it. The force $\vec{F}_{A}$ on $C$ points from $C$ towards $A$, and $\vec{F}_{B}$ points from $C$ towards $B$.
\begin{popfigure}
\begin{center}
\begin{tikzpicture}[scale=2.2]
\coordinate (A) at (0,0);
\coordinate (B) at (2,0);
\coordinate (C) at (1,1.732);
\draw[popmuted] (A)--(B)--(C)--cycle;
\fill[popblue] (A) circle (0.05) node[below left] {$A\ (+\SI{2}{\micro\coulomb})$};
\fill[popblue] (B) circle (0.05) node[below right] {$B\ (+\SI{2}{\micro\coulomb})$};
\fill[poporange] (C) circle (0.05) node[above] {$C\ (-\SI{4}{\micro\coulomb})$};
\draw[->, very thick, popdark] (C) -- ($(C)!0.45!(A)$) node[midway, left] {$\vec{F}_{A}$};
\draw[->, very thick, popdark] (C) -- ($(C)!0.45!(B)$) node[midway, right] {$\vec{F}_{B}$};
\draw[->, very thick, popgreen] (C) -- (1,0.9) node[midway, right] {$\vec{F}$};
\end{tikzpicture}
\end{center}
\legende{Forces on the charge at $C$: attraction towards both $A$ and $B$; the resultant points along the symmetry axis.}
\end{popfigure}
\item Magnitude of each force ($r = 0.20\,\mathrm{m}$):
\[
F_{A} = F_{B} = k\frac{|q_{A}q_{C}|}{r^{2}} = (9.0\times 10^{9})\frac{(2.0\times 10^{-6})(4.0\times 10^{-6})}{(0.20)^{2}} = 1.8\,\mathrm{N}.
\]
\item By symmetry, the resultant lies along the perpendicular bisector of $AB$ (the $y$-axis in the diagram). Each force makes an angle of $30^{\circ}$ with this axis. Horizontal components cancel; vertical components add:
\[
F = 2F_{A}\cos 30^{\circ} = 2\times 1.8\times \frac{\sqrt{3}}{2} = 1.8\sqrt{3} \approx 3.1\,\mathrm{N}.
\]
\textbf{Direction:} along the perpendicular bisector of $AB$, from $C$ towards the midpoint of $AB$ (vertically downwards in the figure).
\end{enumerate}
\end{corrige}

\begin{attention}[Examiner's expectations for Exercise 6]
Draw a clear free-body diagram showing the forces on the charge of interest. Resolve vectors along symmetry axes. Give the final magnitude to two significant figures and describe the direction unambiguously (e.g., ``towards the midpoint of $AB$'').
\end{attention}

\begin{corrige}[Exercise 7 — Charged sphere hanging from a thread]
\begin{enumerate}
\item Free-body diagram: the sphere of mass $m = 0.50\,\mathrm{g} = 5.0\times 10^{-4}\,\mathrm{kg}$ is in equilibrium under three forces — its weight $m\vec{g}$ (vertically down), the tension $\vec{T}$ (along the thread at $30^{\circ}$ to the vertical), and the horizontal electrostatic repulsion $\vec{F}_{e}$ from the fixed positive charge.
\begin{popfigure}
\begin{center}
\begin{tikzpicture}[scale=1.6]
\coordinate (O) at (0,2);
\coordinate (S) at (1,0);
\draw[popmuted, dashed] (O) -- (0,0);
\draw[thick] (O) -- (S);
\fill[poporange] (S) circle (0.09);
\draw[->, very thick, popdark] (S) -- (1,-1.0) node[below] {$m\vec{g}$};
\draw[->, very thick, popblue] (S) -- (0.45,1.05) node[above left] {$\vec{T}$};
\draw[->, very thick, popgreen] (S) -- (2.0,0) node[right] {$\vec{F}_{e}$};
\draw (0,1.4) arc[start angle=-90, end angle=-63, radius=0.6] node[midway, below] {$30^{\circ}$};
\end{tikzpicture}
\end{center}
\legende{Equilibrium of the charged sphere under weight, tension and electrostatic force.}
\end{popfigure}
\item Resolving the tension: vertical component $T\cos 30^{\circ}$, horizontal component $T\sin 30^{\circ}$. Equilibrium conditions:
\[
T\cos 30^{\circ} = mg, \qquad T\sin 30^{\circ} = F_{e}.
\]
Dividing the second by the first eliminates $T$:
\[
F_{e} = mg\tan 30^{\circ} = (5.0\times 10^{-4}\,\mathrm{kg})(9.81\,\mathrm{m\,s^{-2}})(0.5774) = 2.83\times 10^{-3}\,\mathrm{N} \approx 2.8\times 10^{-3}\,\mathrm{N}.
\]
\item Coulomb's law gives $F_{e} = k\dfrac{|q|\,Q}{r^{2}}$ with $Q = +2.0\,\mu\mathrm{C} = 2.0\times 10^{-6}\,\mathrm{C}$ and $r = 0.10\,\mathrm{m}$:
\[
|q| = \frac{F_{e}\,r^{2}}{kQ} = \frac{(2.83\times 10^{-3})(0.10)^{2}}{(9.0\times 10^{9})(2.0\times 10^{-6})} = 1.57\times 10^{-9}\,\mathrm{C} \approx 1.6\,\mathrm{nC}.
\]
Since the sphere is \emph{repelled} by the fixed positive charge (it hangs displaced away from it), $q$ must be \textbf{positive}: $q = +1.6\,\mathrm{nC}$.
\end{enumerate}
\end{corrige}

\begin{attention}[Examiner's expectations for Exercise 7]
The free-body diagram must show all three forces with correct directions and labels. Resolve tension correctly: $T\cos\theta$ vertical, $T\sin\theta$ horizontal. The sign of the charge is deduced from the direction of the electrostatic force (repulsion $\Rightarrow$ same sign).
\end{attention}

\begin{corrige}[Exercise 8 — Electric field of two charges]
\begin{enumerate}
\item Charges $q_{1} = +5.0\,\mathrm{nC}$, $q_{2} = -5.0\,\mathrm{nC}$ separated by $1.0\,\mathrm{m}$. Point $P$ lies on the line, on the side of $q_{1}$ opposite to $q_{2}$. Distances: $r_{1} = 0.50\,\mathrm{m}$ from $q_{1}$, $r_{2} = 1.5\,\mathrm{m}$ from $q_{2}$.
Field due to $q_{1}$ (positive) at $P$ points \emph{away} from $q_{1}$ (i.e., away from $q_{2}$). Field due to $q_{2}$ (negative) points \emph{towards} $q_{2}$. At $P$ these are in opposite directions.
\[
E_{1} = \frac{k|q_{1}|}{r_{1}^{2}} = \frac{(9.0\times 10^{9})(5.0\times 10^{-9})}{(0.50)^{2}} = 180\,\mathrm{N\,C^{-1}},
\]
\[
E_{2} = \frac{k|q_{2}|}{r_{2}^{2}} = \frac{(9.0\times 10^{9})(5.0\times 10^{-9})}{(1.5)^{2}} = 20\,\mathrm{N\,C^{-1}}.
\]
Resultant: $E = E_{1} - E_{2} = 160\,\mathrm{N\,C^{-1}}$, directed away from $q_{1}$ (away from the dipole).
\item Search for a point on the line where $E=0$. Let $x$ be the distance from $q_{1}$.
\begin{itemize}
\item \textbf{Between the charges ($0 < x < 1.0$):} $\vec{E}_{1}$ points towards $q_{2}$ (right), $\vec{E}_{2}$ (negative charge) also points towards $q_{2}$ (right). They add, cannot cancel.
\item \textbf{Beyond $q_{1}$ ($x < 0$):} Distance to $q_{1}$ is $|x| = -x$, to $q_{2}$ is $1.0 - x$. $\vec{E}_{1}$ points left, $\vec{E}_{2}$ points right. Cancellation requires
\[
\frac{k|q_{1}|}{x^{2}} = \frac{k|q_{2}|}{(1.0 - x)^{2}} \;\Rightarrow\; \frac{1}{x^{2}} = \frac{1}{(1.0 - x)^{2}} \;\Rightarrow\; x^{2} = (1.0 - x)^{2}.
\]
Solutions: $x = 1.0 - x \Rightarrow x = 0.50$ (between, already excluded) or $x = -(1.0 - x) \Rightarrow 0 = -1.0$ (impossible). No solution in this region.
\item \textbf{Beyond $q_{2}$ ($x > 1.0$):} Distance to $q_{1}$ is $x$, to $q_{2}$ is $x - 1.0$. $\vec{E}_{1}$ points right, $\vec{E}_{2}$ points left. Same equation $1/x^{2} = 1/(x-1.0)^{2}$ gives $x = x - 1.0$, impossible.
\end{itemize}
\textbf{Conclusion:} For two charges of equal magnitude and opposite sign, there is \emph{no} point on the line joining them where the electric field is zero.
\item If both charges are $+5.0\,\mathrm{nC}$, the fields oppose only \emph{between} the charges. With $x$ measured from $q_{1}$ ($0 < x < 1.0$):
\[
\frac{kq}{x^{2}} = \frac{kq}{(1.0 - x)^{2}} \;\Rightarrow\; x = 1.0 - x \;\Rightarrow\; x = 0.50\,\mathrm{m}.
\]
The zero-field point is at the \textbf{midpoint} of the segment joining the two equal like charges.
\end{enumerate}
\end{corrige}

\begin{attention}[Examiner's expectations for Exercise 8]
In (b), examine all three regions separately. State clearly the direction of each field contribution in each region. The algebraic equation must use distances (positive quantities). For equal and opposite charges, the field is never zero on the line; for equal like charges, it is zero at the midpoint.
\end{attention}

\begin{corrige}[Exercise 9 — Charging by induction (12 marks)]
\textbf{(a) Procedure (6 marks).}
\begin{enumerate}
\item Bring the negatively charged ebonite rod \emph{near} (without touching) the cap of the uncharged gold-leaf electroscope: the leaf diverges. \hfill (1)
\item With the rod held in place, touch the cap with a finger (earthing the electroscope): the leaf falls back to the vertical. \hfill (1)
\item Remove the finger \emph{first}, while the rod is still held near the cap. \hfill (1)
\item Remove the rod: the leaf diverges again — the electroscope is now positively charged. \hfill (1)
\item Diagrams: four clearly labelled sketches showing the rod, the cap, the leaf and the charge distribution ($+$ and $-$ signs) at each step. \hfill (2)
\end{enumerate}
\textbf{(b) Explanation in terms of electrons (4 marks).}
\begin{itemize}
\item The negative rod repels free electrons in the metal cap down into the stem and leaf; the cap becomes positively charged (deficit of electrons) and the leaf (now negatively charged) diverges. \hfill (1)
\item On earthing, the repelled electrons flow from the electroscope through the finger to earth; the leaf collapses because it loses its excess negative charge. \hfill (1)
\item Removing the finger first traps a deficit of electrons on the electroscope (the rod still holds the remaining electrons away from the cap). \hfill (1)
\item When the rod is removed, the net positive charge redistributes uniformly over the cap, stem and leaf, so the leaf diverges: the electroscope carries a permanent \emph{positive} charge, opposite to that of the inducing rod. \hfill (1)
\end{itemize}
\textbf{(c) (2 marks).} Bringing a negatively charged body near the cap of the \emph{positively charged} electroscope repels free electrons from the cap down towards the leaf. These electrons partially neutralise the positive charge on the leaf, reducing its net positive charge, so the leaf \textbf{falls (divergence decreases)}. \hfill (1) This happens because like charges repel: the negative body pushes mobile electrons away from the cap towards the leaf. \hfill (1)
\end{corrige}

\begin{attention}[Examiner's expectations for Exercise 9]
The order of operations (earth, remove finger, \emph{then} remove rod) must be stated explicitly — reversing the last two steps discharges the electroscope and is the classic error. Diagrams must show charge signs on rod, cap and leaf at each stage. In (c), the electroscope is already positively charged; the approaching negative rod drives electrons \emph{towards} the leaf, reducing its divergence.
\end{attention}

\begin{corrige}[Exercise 10 — Thunderstorm safety (10 marks)]
\textbf{(a) Safety in a car (5 marks).}
\begin{itemize}
\item A tree is tall and, when wet, a partial conductor; lightning tends to strike the tallest conducting path to earth, so sheltering under a tree places a person close to the strike path and to dangerous ground currents (step potential). \hfill (2)
\item The metal body of a closed car is a hollow conductor: in electrostatic equilibrium, charge resides on its \emph{outer surface} and the electric field inside is zero (electrostatic shielding / Faraday cage effect). \hfill (2)
\item Hence if lightning strikes the car, the current flows over the outer metal shell to the ground, and the occupants inside remain safe. \hfill (1)
\end{itemize}
\textbf{(b) The lightning conductor (5 marks).}
\begin{itemize}
\item Construction: a pointed copper rod fixed at the highest point of the building, connected by a thick low-resistance copper strip to a metal plate buried deep in moist earth. \hfill (2)
\item By \textbf{point action}, charge density and field are extremely high at the sharp tip; the surrounding air is ionised. \hfill (1)
\item The ionised air allows the induced charge on the building to leak away gently, reducing the chance of a strike. \hfill (1)
\item If a strike does occur, the conductor offers a low-resistance path that channels the large current safely to earth, protecting the building. \hfill (1)
\end{itemize}
\end{corrige}

\begin{attention}[Examiner's expectations for Exercise 10]
Key words expected: \emph{hollow conductor, charge on outer surface, zero internal field, electrostatic shielding, Faraday cage, point action, ionisation of air, low-resistance path to earth}. Vague statements such as ``copper attracts lightning'' earn no credit. Distinguish between the \emph{preventive} role (leakage) and the \emph{protective} role (safe path) of the lightning conductor.
\end{attention}

\begin{corrige}[Exercise 11 — Data analysis: verification of Coulomb's law (18 marks)]
\textbf{(a) (4 marks).} Coulomb's law: the force between two point charges is directly proportional to the product of the charges and inversely proportional to the square of their separation:
\[
F = k\,\frac{q_{1}q_{2}}{r^{2}} \quad\text{with}\quad k = \frac{1}{4\pi\varepsilon_{0}} = 9.0\times 10^{9}\,\mathrm{N\,m^{2}\,C^{-2}}. \hfill (2)
\]
Since $q_{1}q_{2}$ and $k$ are constant, $F = (kq_{1}q_{2})\times\dfrac{1}{r^{2}}$ is of the form $y = mx$: plotting $F$ against $1/r^{2}$ gives a straight line through the origin with gradient $kq_{1}q_{2}$. \hfill (2)

\textbf{(b) (3 marks).} With $r$ in metres, $1/r^{2}$ in $\mathrm{m^{-2}}$:
\begin{center}
\begin{tabularx}{\linewidth}{|l|X|X|X|X|X|X|}
\hline
\rowcolor{popblueL} $r$ / cm & 5.0 & 8.0 & 10.0 & 15.0 & 20.0 & 25.0 \\
\hline
$F$ / $10^{-3}\,\mathrm{N}$ & 14.4 & 5.6 & 3.6 & 1.6 & 0.90 & 0.58 \\
\hline
$1/r^{2}$ / $\mathrm{m^{-2}}$ & 400 & 156 & 100 & 44.4 & 25.0 & 16.0 \\
\hline
\end{tabularx}
\end{center}
(1 mark per two correct values, rounded to 3 significant figures.)

\textbf{(c) (5 marks).}
\begin{popfigure}
\begin{center}
\begin{tikzpicture}
\begin{axis}[
width=0.85\linewidth, height=7cm,
xlabel={$1/r^{2}$ / $\mathrm{m^{-2}}$}, ylabel={$F$ / $10^{-3}\,\mathrm{N}$},
xmin=0, xmax=420, ymin=0, ymax=16,
grid=major, grid style={popline},
tick label style={font=\small}, label style={font=\small},
]
\addplot[only marks, mark=*, mark size=1.8pt, poporange] coordinates {
(16.0,0.58) (25.0,0.90) (44.4,1.6) (100,3.6) (156,5.6) (400,14.4)};
\addplot[domain=0:420, thick, popblue] {0.036*x};
\end{axis}
\end{tikzpicture}
\end{center}
\legende{Graph of $F$ against $1/r^{2}$: the points lie on a straight line through the origin, confirming the inverse-square law.}
\end{popfigure}
Marks: axes correctly labelled with units (1); suitable uniform scales (1); all six points correctly plotted (2); best-fit straight line through the origin (1).

\textbf{(d) (4 marks).} Choose two well-separated points on the line, e.g.\ $(0,\,0)$ and $(400,\,14.4\times 10^{-3})$:
\[
\text{gradient} = \frac{14.4\times 10^{-3}\,\mathrm{N}}{400\,\mathrm{m^{-2}}} = 3.6\times 10^{-5}\,\mathrm{N\,m^{2}}. \hfill (2)
\]
Since gradient $= kq_{1}q_{2}$:
\[
q_{1}q_{2} = \frac{3.6\times 10^{-5}\,\mathrm{N\,m^{2}}}{9.0\times 10^{9}\,\mathrm{N\,m^{2}\,C^{-2}}} = 4.0\times 10^{-15}\,\mathrm{C^{2}}. \hfill (2)
\]
(For example, if $q_{1} = q_{2}$, each charge would be $\sqrt{4.0\times 10^{-15}} \approx 6.3\times 10^{-8}\,\mathrm{C}$.)

\textbf{(e) (2 marks).} Any two of:
\begin{itemize}
\item ensure the spheres are small compared with the separation so they behave as point charges;
\item keep the charges constant (avoid leakage by using dry air / insulating supports);
\item measure $r$ between the \emph{centres} of the spheres;
\item keep other charged objects and earthed conductors far away. \hfill (2)
\end{itemize}
\end{corrige}

\begin{attention}[Examiner's expectations for Exercise 11]
The gradient must carry units ($\mathrm{N\,m^{2}}$) and the product $q_{1}q_{2}$ units of $\mathrm{C^{2}}$ — missing units cost marks. The line of best fit must balance points on both sides and should pass through the origin, since the law predicts direct proportionality. In (e), precautions must be specific to the Coulomb's law experiment (e.g., measuring centre-to-centre distance, avoiding charge leakage).
\end{attention}

\newpage

\coursec{Summary sheet}

\begin{popfigure}
\begin{tikzpicture}[
  central/.style={circle, draw=popdark, fill=popblueL, text=popdark, minimum size=1.5cm, align=center},
  branch/.style={circle, draw=popdark, fill=poporange, text=popdark, minimum size=1.2cm, align=center},
  line/.style={draw=popline, thick}
]
\node[central] (center) {Electrostatics and Electric Fields};
\foreach \angle/\text/\col in {30/Electric Charge/poporange,90/Conductors \& Insulators/popgreen,150/Charging Methods/poppurple,210/Point Action \& Lightning Conductor/poppink,270/Coulomb's Law/popgold,330/Electric Field/popblueL} {
  \node[branch, fill=\col] at (\angle:4cm) (branch-\angle) {\text};
  \draw[line] (center) -- (branch-\angle);
}
\end{tikzpicture}
\legende{Mind map of the chapter Electrostatics and Electric Fields}
\end{popfigure}

\begin{bilan}
\textbf{Key Definitions}
\begin{itemize}
\item \cle{Electric charge}: fundamental property of matter; exists as positive (proton) and negative (electron); quantised in units of $e=1.60\times10^{-19}\,\mathrm{C}$.
\item \cle{Conductor}: material with free charges that move easily (e.g., metals).
\item \cle{Insulator}: material with no free charges; charges remain localised.
\item \cle{Electric field} $\vec{E}$: force per unit positive test charge, $\vec{E}=\vec{F}/q_0$; units $\mathrm{N\,C^{-1}}$ or $\mathrm{V\,m^{-1}}$.
\item \cle{Electric field lines}: directed from positive to negative charges; density proportional to field magnitude.
\end{itemize}

\textbf{Laws, Formulas and Theorems}
\begin{itemize}
\item \cle{Coulomb's Law}: $F = \dfrac{1}{4\pi\varepsilon_0\varepsilon_r}\dfrac{|q_1q_2|}{r^2}$; $\varepsilon_0=8.85\times10^{-12}\,\mathrm{F\,m^{-1}}$.
\item \cle{Principle of superposition}: net force/field is vector sum of individual contributions.
\item \cle{Electric field of a point charge}: $E = \dfrac{1}{4\pi\varepsilon_0\varepsilon_r}\dfrac{|q|}{r^2}$.
\item \cle{Field of a uniformly charged sphere}: outside, as point charge at centre; inside, $E=0$ (conductor) or $E\propto r$ (uniform volume charge).
\item \cle{Point action}: surface charge density $\sigma$ is highest at sharp points; $E \propto \sigma$; used in lightning conductors.
\item \cle{Lightning conductor}: provides a low-resistance path to ground, protecting structures by corona discharge at sharp tip.
\end{itemize}

\textbf{Methods}
\begin{itemize}
\item \textbf{Charging by friction}: rubbing transfers electrons; one body gains electrons (negative), the other loses (positive).
\item \textbf{Charging by conduction}: contact with a charged conductor shares charge.
\item \textbf{Charging by induction}: bring charged object near conductor, ground it, then remove ground and charged object; conductor retains opposite charge.
\item \textbf{Calculating net force/field}: draw free-body diagram, resolve vectors into components, sum components.
\item \textbf{Using symmetry}: for continuous charge distributions, exploit symmetry to simplify integration (e.g., ring, disk, sphere).
\end{itemize}

\textbf{Common Mistakes}
\begin{itemize}
\item Forgetting the vector nature of force and field; adding magnitudes instead of components.
\item Using $\varepsilon_0$ instead of $\varepsilon_0\varepsilon_r$ in a dielectric medium.
\item Confusing electric field direction: field points away from positive, toward negative.
\item Assuming field inside a conductor is non-zero in electrostatic equilibrium.
\item Misapplying Coulomb's law to non-point charges without integration.
\item Neglecting units: charge in coulombs, distance in metres, force in newtons.
\end{itemize}

\textbf{GCE Tips}
\begin{itemize}
\item Always state the law in words before writing the formula.
\item Show vector addition clearly; use component method or scale drawing.
\item For point action, explain why sharp points have high field: small radius of curvature $\Rightarrow$ high $\sigma$.
\item In lightning conductor questions, mention corona discharge and safe path to earth.
\item Check significant figures: use given constants ($G$, $e$, $h$, $c$, $m_e$, $g$) as provided.
\item Practice past paper questions on charging methods and field calculations; they often combine multiple concepts.
\end{itemize}
\end{bilan}

\findecours

\end{document}
